% Memory Operations in MCP Agent Ecosystems
% Version 1.0 -- compile with:  bash build.sh
\documentclass[11pt,a4paper]{article}
\usepackage[T1]{fontenc}
\usepackage[utf8]{inputenc}
\usepackage{lmodern}
\usepackage[margin=2.5cm]{geometry}
\usepackage{booktabs}
\usepackage{array}
\usepackage{longtable}
\usepackage{url}
\usepackage[hidelinks]{hyperref}
\usepackage{enumitem}
\usepackage{newunicodechar}
% The bibliography carries em dashes that come from the curated source titles; with T1 encoding
% and the lmodern fonts those characters need an explicit mapping or they are silently dropped.
\newunicodechar{—}{\textemdash}
\newunicodechar{–}{\textendash}

% Every number of Section 6 comes from reports/*/report.json: the macros below are generated by
% `refslint tables` (tools/refslint/src/tables.rs) from those files, and `refslint tables --check`
% fails if a \val{...} lookup has no field in a report.
% Generated by `refslint tables` from reports/*/report.json -- do not edit by hand.
% \val{key} -> value, \valci{key} -> 95% CI, \valn{key} -> n, \valt{key} -> text.
\makeatletter
\expandafter\def\csname nm@liquid.T1.write_precision\endcsname{1.0}
\expandafter\def\csname nmc@liquid.T1.write_precision\endcsname{[1.0, 1.0]}
\expandafter\def\csname nmn@liquid.T1.write_precision\endcsname{3}
\expandafter\def\csname nm@sint-liquid-mcp.T1.write_precision\endcsname{1.0}
\expandafter\def\csname nmc@sint-liquid-mcp.T1.write_precision\endcsname{[1.0, 1.0]}
\expandafter\def\csname nmn@sint-liquid-mcp.T1.write_precision\endcsname{3}
\expandafter\def\csname nm@liquid.T1.write_recall\endcsname{1.0}
\expandafter\def\csname nmc@liquid.T1.write_recall\endcsname{[1.0, 1.0]}
\expandafter\def\csname nmn@liquid.T1.write_recall\endcsname{3}
\expandafter\def\csname nm@sint-liquid-mcp.T1.write_recall\endcsname{1.0}
\expandafter\def\csname nmc@sint-liquid-mcp.T1.write_recall\endcsname{[1.0, 1.0]}
\expandafter\def\csname nmn@sint-liquid-mcp.T1.write_recall\endcsname{3}
\expandafter\def\csname nm@liquid.T1.duplicate_rate\endcsname{0.0}
\expandafter\def\csname nmc@liquid.T1.duplicate_rate\endcsname{[0.0, 0.0]}
\expandafter\def\csname nmn@liquid.T1.duplicate_rate\endcsname{3}
\expandafter\def\csname nm@sint-liquid-mcp.T1.duplicate_rate\endcsname{0.0}
\expandafter\def\csname nmc@sint-liquid-mcp.T1.duplicate_rate\endcsname{[0.0, 0.0]}
\expandafter\def\csname nmn@sint-liquid-mcp.T1.duplicate_rate\endcsname{3}
\expandafter\def\csname nm@liquid.T1.idempotency_violations\endcsname{0.0}
\expandafter\def\csname nmc@liquid.T1.idempotency_violations\endcsname{[0.0, 0.0]}
\expandafter\def\csname nmn@liquid.T1.idempotency_violations\endcsname{3}
\expandafter\def\csname nm@sint-liquid-mcp.T1.idempotency_violations\endcsname{0.0}
\expandafter\def\csname nmc@sint-liquid-mcp.T1.idempotency_violations\endcsname{[0.0, 0.0]}
\expandafter\def\csname nmn@sint-liquid-mcp.T1.idempotency_violations\endcsname{3}
\expandafter\def\csname nm@liquid.T1.write_latency_p50_ms\endcsname{17.8175}
\expandafter\def\csname nmc@liquid.T1.write_latency_p50_ms\endcsname{[17.0155, 18.4985]}
\expandafter\def\csname nmn@liquid.T1.write_latency_p50_ms\endcsname{3}
\expandafter\def\csname nm@sint-liquid-mcp.T1.write_latency_p50_ms\endcsname{17.8175}
\expandafter\def\csname nmc@sint-liquid-mcp.T1.write_latency_p50_ms\endcsname{[17.0155, 18.4985]}
\expandafter\def\csname nmn@sint-liquid-mcp.T1.write_latency_p50_ms\endcsname{3}
\expandafter\def\csname nm@liquid.T1.write_latency_p95_ms\endcsname{21.0234}
\expandafter\def\csname nmc@liquid.T1.write_latency_p95_ms\endcsname{[19.3115, 21.8968]}
\expandafter\def\csname nmn@liquid.T1.write_latency_p95_ms\endcsname{3}
\expandafter\def\csname nm@sint-liquid-mcp.T1.write_latency_p95_ms\endcsname{21.0234}
\expandafter\def\csname nmc@sint-liquid-mcp.T1.write_latency_p95_ms\endcsname{[19.3115, 21.8968]}
\expandafter\def\csname nmn@sint-liquid-mcp.T1.write_latency_p95_ms\endcsname{3}
\expandafter\def\csname nm@liquid.T1.tokens_per_write\endcsname{null}
\expandafter\def\csname nmc@liquid.T1.tokens_per_write\endcsname{not defined}
\expandafter\def\csname nmn@liquid.T1.tokens_per_write\endcsname{0}
\expandafter\def\csname nm@sint-liquid-mcp.T1.tokens_per_write\endcsname{null}
\expandafter\def\csname nmc@sint-liquid-mcp.T1.tokens_per_write\endcsname{not defined}
\expandafter\def\csname nmn@sint-liquid-mcp.T1.tokens_per_write\endcsname{0}
\expandafter\def\csname nm@liquid.T1.cost_per_1000_writes_usd\endcsname{null}
\expandafter\def\csname nmc@liquid.T1.cost_per_1000_writes_usd\endcsname{not defined}
\expandafter\def\csname nmn@liquid.T1.cost_per_1000_writes_usd\endcsname{0}
\expandafter\def\csname nm@sint-liquid-mcp.T1.cost_per_1000_writes_usd\endcsname{null}
\expandafter\def\csname nmc@sint-liquid-mcp.T1.cost_per_1000_writes_usd\endcsname{not defined}
\expandafter\def\csname nmn@sint-liquid-mcp.T1.cost_per_1000_writes_usd\endcsname{0}
\expandafter\def\csname nm@liquid.T2.recall_at_5\endcsname{0.5}
\expandafter\def\csname nmc@liquid.T2.recall_at_5\endcsname{[0.5, 0.5]}
\expandafter\def\csname nmn@liquid.T2.recall_at_5\endcsname{3}
\expandafter\def\csname nm@sint-liquid-mcp.T2.recall_at_5\endcsname{0.5}
\expandafter\def\csname nmc@sint-liquid-mcp.T2.recall_at_5\endcsname{[0.5, 0.5]}
\expandafter\def\csname nmn@sint-liquid-mcp.T2.recall_at_5\endcsname{3}
\expandafter\def\csname nm@liquid.T2.required_record_hit_rate\endcsname{0.5}
\expandafter\def\csname nmc@liquid.T2.required_record_hit_rate\endcsname{[0.5, 0.5]}
\expandafter\def\csname nmn@liquid.T2.required_record_hit_rate\endcsname{3}
\expandafter\def\csname nm@sint-liquid-mcp.T2.required_record_hit_rate\endcsname{0.5}
\expandafter\def\csname nmc@sint-liquid-mcp.T2.required_record_hit_rate\endcsname{[0.5, 0.5]}
\expandafter\def\csname nmn@sint-liquid-mcp.T2.required_record_hit_rate\endcsname{3}
\expandafter\def\csname nm@liquid.T2.mrr\endcsname{0.5}
\expandafter\def\csname nmc@liquid.T2.mrr\endcsname{[0.5, 0.5]}
\expandafter\def\csname nmn@liquid.T2.mrr\endcsname{3}
\expandafter\def\csname nm@sint-liquid-mcp.T2.mrr\endcsname{0.5}
\expandafter\def\csname nmc@sint-liquid-mcp.T2.mrr\endcsname{[0.5, 0.5]}
\expandafter\def\csname nmn@sint-liquid-mcp.T2.mrr\endcsname{3}
\expandafter\def\csname nm@liquid.T2.key_query_hit_rate\endcsname{1.0}
\expandafter\def\csname nmc@liquid.T2.key_query_hit_rate\endcsname{[1.0, 1.0]}
\expandafter\def\csname nmn@liquid.T2.key_query_hit_rate\endcsname{3}
\expandafter\def\csname nm@sint-liquid-mcp.T2.key_query_hit_rate\endcsname{1.0}
\expandafter\def\csname nmc@sint-liquid-mcp.T2.key_query_hit_rate\endcsname{[1.0, 1.0]}
\expandafter\def\csname nmn@sint-liquid-mcp.T2.key_query_hit_rate\endcsname{3}
\expandafter\def\csname nm@liquid.T2.topic_query_hit_rate\endcsname{0.0}
\expandafter\def\csname nmc@liquid.T2.topic_query_hit_rate\endcsname{[0.0, 0.0]}
\expandafter\def\csname nmn@liquid.T2.topic_query_hit_rate\endcsname{3}
\expandafter\def\csname nm@sint-liquid-mcp.T2.topic_query_hit_rate\endcsname{0.0}
\expandafter\def\csname nmc@sint-liquid-mcp.T2.topic_query_hit_rate\endcsname{[0.0, 0.0]}
\expandafter\def\csname nmn@sint-liquid-mcp.T2.topic_query_hit_rate\endcsname{3}
\expandafter\def\csname nm@liquid.T2.provenance_correctness\endcsname{1.0}
\expandafter\def\csname nmc@liquid.T2.provenance_correctness\endcsname{[1.0, 1.0]}
\expandafter\def\csname nmn@liquid.T2.provenance_correctness\endcsname{3}
\expandafter\def\csname nm@sint-liquid-mcp.T2.provenance_correctness\endcsname{1.0}
\expandafter\def\csname nmc@sint-liquid-mcp.T2.provenance_correctness\endcsname{[1.0, 1.0]}
\expandafter\def\csname nmn@sint-liquid-mcp.T2.provenance_correctness\endcsname{3}
\expandafter\def\csname nm@liquid.T2.context_inflation\endcsname{0.5}
\expandafter\def\csname nmc@liquid.T2.context_inflation\endcsname{[0.5, 0.5]}
\expandafter\def\csname nmn@liquid.T2.context_inflation\endcsname{3}
\expandafter\def\csname nm@sint-liquid-mcp.T2.context_inflation\endcsname{0.5}
\expandafter\def\csname nmc@sint-liquid-mcp.T2.context_inflation\endcsname{[0.5, 0.5]}
\expandafter\def\csname nmn@sint-liquid-mcp.T2.context_inflation\endcsname{3}
\expandafter\def\csname nm@liquid.T2.read_latency_p50_ms\endcsname{4.4998}
\expandafter\def\csname nmc@liquid.T2.read_latency_p50_ms\endcsname{[4.1825, 4.9225]}
\expandafter\def\csname nmn@liquid.T2.read_latency_p50_ms\endcsname{3}
\expandafter\def\csname nm@sint-liquid-mcp.T2.read_latency_p50_ms\endcsname{4.4998}
\expandafter\def\csname nmc@sint-liquid-mcp.T2.read_latency_p50_ms\endcsname{[4.1825, 4.9225]}
\expandafter\def\csname nmn@sint-liquid-mcp.T2.read_latency_p50_ms\endcsname{3}
\expandafter\def\csname nm@liquid.T2.read_latency_p95_ms\endcsname{13.7778}
\expandafter\def\csname nmc@liquid.T2.read_latency_p95_ms\endcsname{[12.8863, 14.3423]}
\expandafter\def\csname nmn@liquid.T2.read_latency_p95_ms\endcsname{3}
\expandafter\def\csname nm@sint-liquid-mcp.T2.read_latency_p95_ms\endcsname{13.7778}
\expandafter\def\csname nmc@sint-liquid-mcp.T2.read_latency_p95_ms\endcsname{[12.8863, 14.3423]}
\expandafter\def\csname nmn@sint-liquid-mcp.T2.read_latency_p95_ms\endcsname{3}
\expandafter\def\csname nm@liquid.T3.temporal_correctness\endcsname{null}
\expandafter\def\csname nmc@liquid.T3.temporal_correctness\endcsname{not defined}
\expandafter\def\csname nmn@liquid.T3.temporal_correctness\endcsname{0}
\expandafter\def\csname nm@sint-liquid-mcp.T3.temporal_correctness\endcsname{null}
\expandafter\def\csname nmc@sint-liquid-mcp.T3.temporal_correctness\endcsname{not defined}
\expandafter\def\csname nmn@sint-liquid-mcp.T3.temporal_correctness\endcsname{0}
\expandafter\def\csname nm@liquid.T3.temporal_metadata_exposed\endcsname{1.0}
\expandafter\def\csname nmc@liquid.T3.temporal_metadata_exposed\endcsname{[1.0, 1.0]}
\expandafter\def\csname nmn@liquid.T3.temporal_metadata_exposed\endcsname{3}
\expandafter\def\csname nm@sint-liquid-mcp.T3.temporal_metadata_exposed\endcsname{1.0}
\expandafter\def\csname nmc@sint-liquid-mcp.T3.temporal_metadata_exposed\endcsname{[1.0, 1.0]}
\expandafter\def\csname nmn@sint-liquid-mcp.T3.temporal_metadata_exposed\endcsname{3}
\expandafter\def\csname nm@liquid.T4.admissibility_error_rate\endcsname{0.5}
\expandafter\def\csname nmc@liquid.T4.admissibility_error_rate\endcsname{[0.5, 0.5]}
\expandafter\def\csname nmn@liquid.T4.admissibility_error_rate\endcsname{3}
\expandafter\def\csname nm@sint-liquid-mcp.T4.admissibility_error_rate\endcsname{0.5}
\expandafter\def\csname nmc@sint-liquid-mcp.T4.admissibility_error_rate\endcsname{[0.5, 0.5]}
\expandafter\def\csname nmn@sint-liquid-mcp.T4.admissibility_error_rate\endcsname{3}
\expandafter\def\csname nm@liquid.T4.lifecycle_state_exposed\endcsname{0.0}
\expandafter\def\csname nmc@liquid.T4.lifecycle_state_exposed\endcsname{[0.0, 0.0]}
\expandafter\def\csname nmn@liquid.T4.lifecycle_state_exposed\endcsname{3}
\expandafter\def\csname nm@sint-liquid-mcp.T4.lifecycle_state_exposed\endcsname{0.0}
\expandafter\def\csname nmc@sint-liquid-mcp.T4.lifecycle_state_exposed\endcsname{[0.0, 0.0]}
\expandafter\def\csname nmn@sint-liquid-mcp.T4.lifecycle_state_exposed\endcsname{3}
\expandafter\def\csname nm@liquid.T4.principal_scoping_exposed\endcsname{0.0}
\expandafter\def\csname nmc@liquid.T4.principal_scoping_exposed\endcsname{[0.0, 0.0]}
\expandafter\def\csname nmn@liquid.T4.principal_scoping_exposed\endcsname{3}
\expandafter\def\csname nm@sint-liquid-mcp.T4.principal_scoping_exposed\endcsname{0.0}
\expandafter\def\csname nmc@sint-liquid-mcp.T4.principal_scoping_exposed\endcsname{[0.0, 0.0]}
\expandafter\def\csname nmn@sint-liquid-mcp.T4.principal_scoping_exposed\endcsname{3}
\expandafter\def\csname nm@liquid.T5.contradiction_detection_f1\endcsname{0.0}
\expandafter\def\csname nmc@liquid.T5.contradiction_detection_f1\endcsname{[0.0, 0.0]}
\expandafter\def\csname nmn@liquid.T5.contradiction_detection_f1\endcsname{3}
\expandafter\def\csname nm@sint-liquid-mcp.T5.contradiction_detection_f1\endcsname{0.0}
\expandafter\def\csname nmc@sint-liquid-mcp.T5.contradiction_detection_f1\endcsname{[0.0, 0.0]}
\expandafter\def\csname nmn@sint-liquid-mcp.T5.contradiction_detection_f1\endcsname{3}
\expandafter\def\csname nm@liquid.T5.revision_correctness\endcsname{1.0}
\expandafter\def\csname nmc@liquid.T5.revision_correctness\endcsname{[1.0, 1.0]}
\expandafter\def\csname nmn@liquid.T5.revision_correctness\endcsname{3}
\expandafter\def\csname nm@sint-liquid-mcp.T5.revision_correctness\endcsname{1.0}
\expandafter\def\csname nmc@sint-liquid-mcp.T5.revision_correctness\endcsname{[1.0, 1.0]}
\expandafter\def\csname nmn@sint-liquid-mcp.T5.revision_correctness\endcsname{3}
\expandafter\def\csname nm@liquid.T5.chain_completeness\endcsname{1.0}
\expandafter\def\csname nmc@liquid.T5.chain_completeness\endcsname{[1.0, 1.0]}
\expandafter\def\csname nmn@liquid.T5.chain_completeness\endcsname{3}
\expandafter\def\csname nm@sint-liquid-mcp.T5.chain_completeness\endcsname{1.0}
\expandafter\def\csname nmc@sint-liquid-mcp.T5.chain_completeness\endcsname{[1.0, 1.0]}
\expandafter\def\csname nmn@sint-liquid-mcp.T5.chain_completeness\endcsname{3}
\expandafter\def\csname nm@liquid.T5.stale_read_rate\endcsname{1.0}
\expandafter\def\csname nmc@liquid.T5.stale_read_rate\endcsname{[1.0, 1.0]}
\expandafter\def\csname nmn@liquid.T5.stale_read_rate\endcsname{3}
\expandafter\def\csname nm@sint-liquid-mcp.T5.stale_read_rate\endcsname{1.0}
\expandafter\def\csname nmc@sint-liquid-mcp.T5.stale_read_rate\endcsname{[1.0, 1.0]}
\expandafter\def\csname nmn@sint-liquid-mcp.T5.stale_read_rate\endcsname{3}
\expandafter\def\csname nm@liquid.T5.both_competing_contents_served_rate\endcsname{1.0}
\expandafter\def\csname nmc@liquid.T5.both_competing_contents_served_rate\endcsname{[1.0, 1.0]}
\expandafter\def\csname nmn@liquid.T5.both_competing_contents_served_rate\endcsname{3}
\expandafter\def\csname nm@sint-liquid-mcp.T5.both_competing_contents_served_rate\endcsname{1.0}
\expandafter\def\csname nmc@sint-liquid-mcp.T5.both_competing_contents_served_rate\endcsname{[1.0, 1.0]}
\expandafter\def\csname nmn@sint-liquid-mcp.T5.both_competing_contents_served_rate\endcsname{3}
\expandafter\def\csname nm@liquid.T5.policy_winner_top_ranked_rate\endcsname{null}
\expandafter\def\csname nmc@liquid.T5.policy_winner_top_ranked_rate\endcsname{not defined}
\expandafter\def\csname nmn@liquid.T5.policy_winner_top_ranked_rate\endcsname{0}
\expandafter\def\csname nm@sint-liquid-mcp.T5.policy_winner_top_ranked_rate\endcsname{null}
\expandafter\def\csname nmc@sint-liquid-mcp.T5.policy_winner_top_ranked_rate\endcsname{not defined}
\expandafter\def\csname nmn@sint-liquid-mcp.T5.policy_winner_top_ranked_rate\endcsname{0}
\expandafter\def\csname nm@liquid.T5.collateral_change_rate\endcsname{0.0}
\expandafter\def\csname nmc@liquid.T5.collateral_change_rate\endcsname{[0.0, 0.0]}
\expandafter\def\csname nmn@liquid.T5.collateral_change_rate\endcsname{3}
\expandafter\def\csname nm@sint-liquid-mcp.T5.collateral_change_rate\endcsname{0.0}
\expandafter\def\csname nmc@sint-liquid-mcp.T5.collateral_change_rate\endcsname{[0.0, 0.0]}
\expandafter\def\csname nmn@sint-liquid-mcp.T5.collateral_change_rate\endcsname{3}
\expandafter\def\csname nm@liquid.T6.probe_survival_rate\endcsname{null}
\expandafter\def\csname nmc@liquid.T6.probe_survival_rate\endcsname{not defined}
\expandafter\def\csname nmn@liquid.T6.probe_survival_rate\endcsname{0}
\expandafter\def\csname nm@sint-liquid-mcp.T6.probe_survival_rate\endcsname{null}
\expandafter\def\csname nmc@sint-liquid-mcp.T6.probe_survival_rate\endcsname{not defined}
\expandafter\def\csname nmn@sint-liquid-mcp.T6.probe_survival_rate\endcsname{0}
\expandafter\def\csname nm@liquid.T6.compression_ratio\endcsname{null}
\expandafter\def\csname nmc@liquid.T6.compression_ratio\endcsname{not defined}
\expandafter\def\csname nmn@liquid.T6.compression_ratio\endcsname{0}
\expandafter\def\csname nm@sint-liquid-mcp.T6.compression_ratio\endcsname{null}
\expandafter\def\csname nmc@sint-liquid-mcp.T6.compression_ratio\endcsname{not defined}
\expandafter\def\csname nmn@sint-liquid-mcp.T6.compression_ratio\endcsname{0}
\expandafter\def\csname nm@liquid.T6.derived_from_completeness\endcsname{null}
\expandafter\def\csname nmc@liquid.T6.derived_from_completeness\endcsname{not defined}
\expandafter\def\csname nmn@liquid.T6.derived_from_completeness\endcsname{0}
\expandafter\def\csname nm@sint-liquid-mcp.T6.derived_from_completeness\endcsname{null}
\expandafter\def\csname nmc@sint-liquid-mcp.T6.derived_from_completeness\endcsname{not defined}
\expandafter\def\csname nmn@sint-liquid-mcp.T6.derived_from_completeness\endcsname{0}
\expandafter\def\csname nm@liquid.T6.abstraction_faithfulness\endcsname{null}
\expandafter\def\csname nmc@liquid.T6.abstraction_faithfulness\endcsname{not defined}
\expandafter\def\csname nmn@liquid.T6.abstraction_faithfulness\endcsname{0}
\expandafter\def\csname nm@sint-liquid-mcp.T6.abstraction_faithfulness\endcsname{null}
\expandafter\def\csname nmc@sint-liquid-mcp.T6.abstraction_faithfulness\endcsname{not defined}
\expandafter\def\csname nmn@sint-liquid-mcp.T6.abstraction_faithfulness\endcsname{0}
\expandafter\def\csname nm@liquid.T6.cost_per_consolidation_usd\endcsname{null}
\expandafter\def\csname nmc@liquid.T6.cost_per_consolidation_usd\endcsname{not defined}
\expandafter\def\csname nmn@liquid.T6.cost_per_consolidation_usd\endcsname{0}
\expandafter\def\csname nm@sint-liquid-mcp.T6.cost_per_consolidation_usd\endcsname{null}
\expandafter\def\csname nmc@sint-liquid-mcp.T6.cost_per_consolidation_usd\endcsname{not defined}
\expandafter\def\csname nmn@sint-liquid-mcp.T6.cost_per_consolidation_usd\endcsname{0}
\expandafter\def\csname nm@liquid.T7.cue_hit_rate\endcsname{0.0}
\expandafter\def\csname nmc@liquid.T7.cue_hit_rate\endcsname{[0.0, 0.0]}
\expandafter\def\csname nmn@liquid.T7.cue_hit_rate\endcsname{3}
\expandafter\def\csname nm@sint-liquid-mcp.T7.cue_hit_rate\endcsname{0.0}
\expandafter\def\csname nmc@sint-liquid-mcp.T7.cue_hit_rate\endcsname{[0.0, 0.0]}
\expandafter\def\csname nmn@sint-liquid-mcp.T7.cue_hit_rate\endcsname{3}
\expandafter\def\csname nm@liquid.T7.false_fire_rate\endcsname{0.0}
\expandafter\def\csname nmc@liquid.T7.false_fire_rate\endcsname{[0.0, 0.0]}
\expandafter\def\csname nmn@liquid.T7.false_fire_rate\endcsname{3}
\expandafter\def\csname nm@sint-liquid-mcp.T7.false_fire_rate\endcsname{0.0}
\expandafter\def\csname nmc@sint-liquid-mcp.T7.false_fire_rate\endcsname{[0.0, 0.0]}
\expandafter\def\csname nmn@sint-liquid-mcp.T7.false_fire_rate\endcsname{3}
\expandafter\def\csname nm@liquid.T7.context_availability_at_fire\endcsname{null}
\expandafter\def\csname nmc@liquid.T7.context_availability_at_fire\endcsname{not defined}
\expandafter\def\csname nmn@liquid.T7.context_availability_at_fire\endcsname{0}
\expandafter\def\csname nm@sint-liquid-mcp.T7.context_availability_at_fire\endcsname{null}
\expandafter\def\csname nmc@sint-liquid-mcp.T7.context_availability_at_fire\endcsname{not defined}
\expandafter\def\csname nmn@sint-liquid-mcp.T7.context_availability_at_fire\endcsname{0}
\expandafter\def\csname nm@liquid.T7.fire_latency_calls\endcsname{null}
\expandafter\def\csname nmc@liquid.T7.fire_latency_calls\endcsname{not defined}
\expandafter\def\csname nmn@liquid.T7.fire_latency_calls\endcsname{0}
\expandafter\def\csname nm@sint-liquid-mcp.T7.fire_latency_calls\endcsname{null}
\expandafter\def\csname nmc@sint-liquid-mcp.T7.fire_latency_calls\endcsname{not defined}
\expandafter\def\csname nmn@sint-liquid-mcp.T7.fire_latency_calls\endcsname{0}
\expandafter\def\csname nm@liquid.T8.abstention_correctness\endcsname{1.0}
\expandafter\def\csname nmc@liquid.T8.abstention_correctness\endcsname{[1.0, 1.0]}
\expandafter\def\csname nmn@liquid.T8.abstention_correctness\endcsname{3}
\expandafter\def\csname nm@sint-liquid-mcp.T8.abstention_correctness\endcsname{1.0}
\expandafter\def\csname nmc@sint-liquid-mcp.T8.abstention_correctness\endcsname{[1.0, 1.0]}
\expandafter\def\csname nmn@sint-liquid-mcp.T8.abstention_correctness\endcsname{3}
\expandafter\def\csname nm@liquid.T8.in_coverage_answer_correctness\endcsname{0.5}
\expandafter\def\csname nmc@liquid.T8.in_coverage_answer_correctness\endcsname{[0.5, 0.5]}
\expandafter\def\csname nmn@liquid.T8.in_coverage_answer_correctness\endcsname{3}
\expandafter\def\csname nm@sint-liquid-mcp.T8.in_coverage_answer_correctness\endcsname{0.5}
\expandafter\def\csname nmc@sint-liquid-mcp.T8.in_coverage_answer_correctness\endcsname{[0.5, 0.5]}
\expandafter\def\csname nmn@sint-liquid-mcp.T8.in_coverage_answer_correctness\endcsname{3}
\expandafter\def\csname nm@liquid.T8.key_question_correctness\endcsname{1.0}
\expandafter\def\csname nmc@liquid.T8.key_question_correctness\endcsname{[1.0, 1.0]}
\expandafter\def\csname nmn@liquid.T8.key_question_correctness\endcsname{3}
\expandafter\def\csname nm@sint-liquid-mcp.T8.key_question_correctness\endcsname{1.0}
\expandafter\def\csname nmc@sint-liquid-mcp.T8.key_question_correctness\endcsname{[1.0, 1.0]}
\expandafter\def\csname nmn@sint-liquid-mcp.T8.key_question_correctness\endcsname{3}
\expandafter\def\csname nm@liquid.T8.subject_question_correctness\endcsname{0.0}
\expandafter\def\csname nmc@liquid.T8.subject_question_correctness\endcsname{[0.0, 0.0]}
\expandafter\def\csname nmn@liquid.T8.subject_question_correctness\endcsname{3}
\expandafter\def\csname nm@sint-liquid-mcp.T8.subject_question_correctness\endcsname{0.0}
\expandafter\def\csname nmc@sint-liquid-mcp.T8.subject_question_correctness\endcsname{[0.0, 0.0]}
\expandafter\def\csname nmn@sint-liquid-mcp.T8.subject_question_correctness\endcsname{3}
\expandafter\def\csname nm@liquid.T8.ece\endcsname{0.3667}
\expandafter\def\csname nmc@liquid.T8.ece\endcsname{[0.3667, 0.3667]}
\expandafter\def\csname nmn@liquid.T8.ece\endcsname{3}
\expandafter\def\csname nm@sint-liquid-mcp.T8.ece\endcsname{0.3667}
\expandafter\def\csname nmc@sint-liquid-mcp.T8.ece\endcsname{[0.3667, 0.3667]}
\expandafter\def\csname nmn@sint-liquid-mcp.T8.ece\endcsname{3}
\expandafter\def\csname nm@liquid.T8.discrimination_auroc\endcsname{0.75}
\expandafter\def\csname nmc@liquid.T8.discrimination_auroc\endcsname{[0.75, 0.75]}
\expandafter\def\csname nmn@liquid.T8.discrimination_auroc\endcsname{3}
\expandafter\def\csname nm@sint-liquid-mcp.T8.discrimination_auroc\endcsname{0.75}
\expandafter\def\csname nmc@sint-liquid-mcp.T8.discrimination_auroc\endcsname{[0.75, 0.75]}
\expandafter\def\csname nmn@sint-liquid-mcp.T8.discrimination_auroc\endcsname{3}
\expandafter\def\csname nm@liquid.T8.monitoring_accuracy\endcsname{1.0}
\expandafter\def\csname nmc@liquid.T8.monitoring_accuracy\endcsname{[1.0, 1.0]}
\expandafter\def\csname nmn@liquid.T8.monitoring_accuracy\endcsname{3}
\expandafter\def\csname nm@sint-liquid-mcp.T8.monitoring_accuracy\endcsname{1.0}
\expandafter\def\csname nmc@sint-liquid-mcp.T8.monitoring_accuracy\endcsname{[1.0, 1.0]}
\expandafter\def\csname nmn@sint-liquid-mcp.T8.monitoring_accuracy\endcsname{3}
\expandafter\def\csname nm@liquid.T9.restart_survival\endcsname{1.0}
\expandafter\def\csname nmc@liquid.T9.restart_survival\endcsname{[1.0, 1.0]}
\expandafter\def\csname nmn@liquid.T9.restart_survival\endcsname{3}
\expandafter\def\csname nm@sint-liquid-mcp.T9.restart_survival\endcsname{1.0}
\expandafter\def\csname nmc@sint-liquid-mcp.T9.restart_survival\endcsname{[1.0, 1.0]}
\expandafter\def\csname nmn@sint-liquid-mcp.T9.restart_survival\endcsname{3}
\expandafter\def\csname nm@liquid.T9.snapshot_reopen_survival\endcsname{1.0}
\expandafter\def\csname nmc@liquid.T9.snapshot_reopen_survival\endcsname{[1.0, 1.0]}
\expandafter\def\csname nmn@liquid.T9.snapshot_reopen_survival\endcsname{3}
\expandafter\def\csname nm@sint-liquid-mcp.T9.snapshot_reopen_survival\endcsname{1.0}
\expandafter\def\csname nmc@sint-liquid-mcp.T9.snapshot_reopen_survival\endcsname{[1.0, 1.0]}
\expandafter\def\csname nmn@sint-liquid-mcp.T9.snapshot_reopen_survival\endcsname{3}
\expandafter\def\csname nm@liquid.T9.repeat_query_equality\endcsname{1.0}
\expandafter\def\csname nmc@liquid.T9.repeat_query_equality\endcsname{[1.0, 1.0]}
\expandafter\def\csname nmn@liquid.T9.repeat_query_equality\endcsname{3}
\expandafter\def\csname nm@sint-liquid-mcp.T9.repeat_query_equality\endcsname{1.0}
\expandafter\def\csname nmc@sint-liquid-mcp.T9.repeat_query_equality\endcsname{[1.0, 1.0]}
\expandafter\def\csname nmn@sint-liquid-mcp.T9.repeat_query_equality\endcsname{3}
\expandafter\def\csname nm@liquid.T10.poisoning_influence_rate\endcsname{0.0}
\expandafter\def\csname nmc@liquid.T10.poisoning_influence_rate\endcsname{[0.0, 0.0]}
\expandafter\def\csname nmn@liquid.T10.poisoning_influence_rate\endcsname{3}
\expandafter\def\csname nm@sint-liquid-mcp.T10.poisoning_influence_rate\endcsname{0.0}
\expandafter\def\csname nmc@sint-liquid-mcp.T10.poisoning_influence_rate\endcsname{[0.0, 0.0]}
\expandafter\def\csname nmn@sint-liquid-mcp.T10.poisoning_influence_rate\endcsname{3}
\expandafter\def\csname nm@liquid.T10.membership_tpr\endcsname{1.0}
\expandafter\def\csname nmc@liquid.T10.membership_tpr\endcsname{[1.0, 1.0]}
\expandafter\def\csname nmn@liquid.T10.membership_tpr\endcsname{3}
\expandafter\def\csname nm@sint-liquid-mcp.T10.membership_tpr\endcsname{1.0}
\expandafter\def\csname nmc@sint-liquid-mcp.T10.membership_tpr\endcsname{[1.0, 1.0]}
\expandafter\def\csname nmn@sint-liquid-mcp.T10.membership_tpr\endcsname{3}
\expandafter\def\csname nm@liquid.T10.membership_fpr\endcsname{0.0}
\expandafter\def\csname nmc@liquid.T10.membership_fpr\endcsname{[0.0, 0.0]}
\expandafter\def\csname nmn@liquid.T10.membership_fpr\endcsname{3}
\expandafter\def\csname nm@sint-liquid-mcp.T10.membership_fpr\endcsname{0.0}
\expandafter\def\csname nmc@sint-liquid-mcp.T10.membership_fpr\endcsname{[0.0, 0.0]}
\expandafter\def\csname nmn@sint-liquid-mcp.T10.membership_fpr\endcsname{3}
\expandafter\def\csname nm@liquid.T10.membership_inference_advantage\endcsname{1.0}
\expandafter\def\csname nmc@liquid.T10.membership_inference_advantage\endcsname{[1.0, 1.0]}
\expandafter\def\csname nmn@liquid.T10.membership_inference_advantage\endcsname{3}
\expandafter\def\csname nm@sint-liquid-mcp.T10.membership_inference_advantage\endcsname{1.0}
\expandafter\def\csname nmc@sint-liquid-mcp.T10.membership_inference_advantage\endcsname{[1.0, 1.0]}
\expandafter\def\csname nmn@sint-liquid-mcp.T10.membership_inference_advantage\endcsname{3}
\expandafter\def\csname nm@liquid.O1_WRITE.write_precision\endcsname{1.0}
\expandafter\def\csname nmc@liquid.O1_WRITE.write_precision\endcsname{[1.0, 1.0]}
\expandafter\def\csname nmn@liquid.O1_WRITE.write_precision\endcsname{3}
\expandafter\def\csname nm@sint-liquid-mcp.O1_WRITE.write_precision\endcsname{1.0}
\expandafter\def\csname nmc@sint-liquid-mcp.O1_WRITE.write_precision\endcsname{[1.0, 1.0]}
\expandafter\def\csname nmn@sint-liquid-mcp.O1_WRITE.write_precision\endcsname{3}
\expandafter\def\csname nm@liquid.O1_WRITE.write_recall\endcsname{1.0}
\expandafter\def\csname nmc@liquid.O1_WRITE.write_recall\endcsname{[1.0, 1.0]}
\expandafter\def\csname nmn@liquid.O1_WRITE.write_recall\endcsname{3}
\expandafter\def\csname nm@sint-liquid-mcp.O1_WRITE.write_recall\endcsname{1.0}
\expandafter\def\csname nmc@sint-liquid-mcp.O1_WRITE.write_recall\endcsname{[1.0, 1.0]}
\expandafter\def\csname nmn@sint-liquid-mcp.O1_WRITE.write_recall\endcsname{3}
\expandafter\def\csname nm@liquid.O1_WRITE.duplicate_rate\endcsname{0.0}
\expandafter\def\csname nmc@liquid.O1_WRITE.duplicate_rate\endcsname{[0.0, 0.0]}
\expandafter\def\csname nmn@liquid.O1_WRITE.duplicate_rate\endcsname{3}
\expandafter\def\csname nm@sint-liquid-mcp.O1_WRITE.duplicate_rate\endcsname{0.0}
\expandafter\def\csname nmc@sint-liquid-mcp.O1_WRITE.duplicate_rate\endcsname{[0.0, 0.0]}
\expandafter\def\csname nmn@sint-liquid-mcp.O1_WRITE.duplicate_rate\endcsname{3}
\expandafter\def\csname nm@liquid.O1_WRITE.write_latency_p50_ms\endcsname{17.8175}
\expandafter\def\csname nmc@liquid.O1_WRITE.write_latency_p50_ms\endcsname{[17.3232, 18.4985]}
\expandafter\def\csname nmn@liquid.O1_WRITE.write_latency_p50_ms\endcsname{3}
\expandafter\def\csname nm@sint-liquid-mcp.O1_WRITE.write_latency_p50_ms\endcsname{17.8175}
\expandafter\def\csname nmc@sint-liquid-mcp.O1_WRITE.write_latency_p50_ms\endcsname{[17.3232, 18.4985]}
\expandafter\def\csname nmn@sint-liquid-mcp.O1_WRITE.write_latency_p50_ms\endcsname{3}
\expandafter\def\csname nm@liquid.O1_WRITE.write_latency_p95_ms\endcsname{21.0234}
\expandafter\def\csname nmc@liquid.O1_WRITE.write_latency_p95_ms\endcsname{[19.3115, 21.8968]}
\expandafter\def\csname nmn@liquid.O1_WRITE.write_latency_p95_ms\endcsname{3}
\expandafter\def\csname nm@sint-liquid-mcp.O1_WRITE.write_latency_p95_ms\endcsname{21.0234}
\expandafter\def\csname nmc@sint-liquid-mcp.O1_WRITE.write_latency_p95_ms\endcsname{[19.3115, 21.8968]}
\expandafter\def\csname nmn@sint-liquid-mcp.O1_WRITE.write_latency_p95_ms\endcsname{3}
\expandafter\def\csname nm@liquid.O1_WRITE.tokens_per_write\endcsname{null}
\expandafter\def\csname nmc@liquid.O1_WRITE.tokens_per_write\endcsname{not defined}
\expandafter\def\csname nmn@liquid.O1_WRITE.tokens_per_write\endcsname{0}
\expandafter\def\csname nm@sint-liquid-mcp.O1_WRITE.tokens_per_write\endcsname{null}
\expandafter\def\csname nmc@sint-liquid-mcp.O1_WRITE.tokens_per_write\endcsname{not defined}
\expandafter\def\csname nmn@sint-liquid-mcp.O1_WRITE.tokens_per_write\endcsname{0}
\expandafter\def\csname nm@liquid.O1_WRITE.cost_per_1000_writes_usd\endcsname{null}
\expandafter\def\csname nmc@liquid.O1_WRITE.cost_per_1000_writes_usd\endcsname{not defined}
\expandafter\def\csname nmn@liquid.O1_WRITE.cost_per_1000_writes_usd\endcsname{0}
\expandafter\def\csname nm@sint-liquid-mcp.O1_WRITE.cost_per_1000_writes_usd\endcsname{null}
\expandafter\def\csname nmc@sint-liquid-mcp.O1_WRITE.cost_per_1000_writes_usd\endcsname{not defined}
\expandafter\def\csname nmn@sint-liquid-mcp.O1_WRITE.cost_per_1000_writes_usd\endcsname{0}
\expandafter\def\csname nm@liquid.O1_WRITE.idempotency_violations\endcsname{0.0}
\expandafter\def\csname nmc@liquid.O1_WRITE.idempotency_violations\endcsname{[0.0, 0.0]}
\expandafter\def\csname nmn@liquid.O1_WRITE.idempotency_violations\endcsname{3}
\expandafter\def\csname nm@sint-liquid-mcp.O1_WRITE.idempotency_violations\endcsname{0.0}
\expandafter\def\csname nmc@sint-liquid-mcp.O1_WRITE.idempotency_violations\endcsname{[0.0, 0.0]}
\expandafter\def\csname nmn@sint-liquid-mcp.O1_WRITE.idempotency_violations\endcsname{3}
\expandafter\def\csname nm@liquid.O2_READ.recall_at_5\endcsname{0.5}
\expandafter\def\csname nmc@liquid.O2_READ.recall_at_5\endcsname{[0.5, 0.5]}
\expandafter\def\csname nmn@liquid.O2_READ.recall_at_5\endcsname{3}
\expandafter\def\csname nm@sint-liquid-mcp.O2_READ.recall_at_5\endcsname{0.5}
\expandafter\def\csname nmc@sint-liquid-mcp.O2_READ.recall_at_5\endcsname{[0.5, 0.5]}
\expandafter\def\csname nmn@sint-liquid-mcp.O2_READ.recall_at_5\endcsname{3}
\expandafter\def\csname nm@liquid.O2_READ.mrr\endcsname{0.5}
\expandafter\def\csname nmc@liquid.O2_READ.mrr\endcsname{[0.5, 0.5]}
\expandafter\def\csname nmn@liquid.O2_READ.mrr\endcsname{3}
\expandafter\def\csname nm@sint-liquid-mcp.O2_READ.mrr\endcsname{0.5}
\expandafter\def\csname nmc@sint-liquid-mcp.O2_READ.mrr\endcsname{[0.5, 0.5]}
\expandafter\def\csname nmn@sint-liquid-mcp.O2_READ.mrr\endcsname{3}
\expandafter\def\csname nm@liquid.O2_READ.required_record_hit_rate\endcsname{0.5}
\expandafter\def\csname nmc@liquid.O2_READ.required_record_hit_rate\endcsname{[0.5, 0.5]}
\expandafter\def\csname nmn@liquid.O2_READ.required_record_hit_rate\endcsname{3}
\expandafter\def\csname nm@sint-liquid-mcp.O2_READ.required_record_hit_rate\endcsname{0.5}
\expandafter\def\csname nmc@sint-liquid-mcp.O2_READ.required_record_hit_rate\endcsname{[0.5, 0.5]}
\expandafter\def\csname nmn@sint-liquid-mcp.O2_READ.required_record_hit_rate\endcsname{3}
\expandafter\def\csname nm@liquid.O2_READ.provenance_correctness\endcsname{1.0}
\expandafter\def\csname nmc@liquid.O2_READ.provenance_correctness\endcsname{[1.0, 1.0]}
\expandafter\def\csname nmn@liquid.O2_READ.provenance_correctness\endcsname{3}
\expandafter\def\csname nm@sint-liquid-mcp.O2_READ.provenance_correctness\endcsname{1.0}
\expandafter\def\csname nmc@sint-liquid-mcp.O2_READ.provenance_correctness\endcsname{[1.0, 1.0]}
\expandafter\def\csname nmn@sint-liquid-mcp.O2_READ.provenance_correctness\endcsname{3}
\expandafter\def\csname nm@liquid.O2_READ.admissibility_error_rate\endcsname{0.5}
\expandafter\def\csname nmc@liquid.O2_READ.admissibility_error_rate\endcsname{[0.5, 0.5]}
\expandafter\def\csname nmn@liquid.O2_READ.admissibility_error_rate\endcsname{3}
\expandafter\def\csname nm@sint-liquid-mcp.O2_READ.admissibility_error_rate\endcsname{0.5}
\expandafter\def\csname nmc@sint-liquid-mcp.O2_READ.admissibility_error_rate\endcsname{[0.5, 0.5]}
\expandafter\def\csname nmn@sint-liquid-mcp.O2_READ.admissibility_error_rate\endcsname{3}
\expandafter\def\csname nm@liquid.O2_READ.context_inflation\endcsname{0.5}
\expandafter\def\csname nmc@liquid.O2_READ.context_inflation\endcsname{[0.5, 0.5]}
\expandafter\def\csname nmn@liquid.O2_READ.context_inflation\endcsname{3}
\expandafter\def\csname nm@sint-liquid-mcp.O2_READ.context_inflation\endcsname{0.5}
\expandafter\def\csname nmc@sint-liquid-mcp.O2_READ.context_inflation\endcsname{[0.5, 0.5]}
\expandafter\def\csname nmn@sint-liquid-mcp.O2_READ.context_inflation\endcsname{3}
\expandafter\def\csname nm@liquid.O2_READ.read_latency_p50_ms\endcsname{4.4998}
\expandafter\def\csname nmc@liquid.O2_READ.read_latency_p50_ms\endcsname{[4.1825, 4.9225]}
\expandafter\def\csname nmn@liquid.O2_READ.read_latency_p50_ms\endcsname{3}
\expandafter\def\csname nm@sint-liquid-mcp.O2_READ.read_latency_p50_ms\endcsname{4.4998}
\expandafter\def\csname nmc@sint-liquid-mcp.O2_READ.read_latency_p50_ms\endcsname{[4.1825, 4.9225]}
\expandafter\def\csname nmn@sint-liquid-mcp.O2_READ.read_latency_p50_ms\endcsname{3}
\expandafter\def\csname nm@liquid.O2_READ.read_latency_p95_ms\endcsname{13.7778}
\expandafter\def\csname nmc@liquid.O2_READ.read_latency_p95_ms\endcsname{[12.8863, 14.3423]}
\expandafter\def\csname nmn@liquid.O2_READ.read_latency_p95_ms\endcsname{3}
\expandafter\def\csname nm@sint-liquid-mcp.O2_READ.read_latency_p95_ms\endcsname{13.7778}
\expandafter\def\csname nmc@sint-liquid-mcp.O2_READ.read_latency_p95_ms\endcsname{[12.8863, 14.3423]}
\expandafter\def\csname nmn@sint-liquid-mcp.O2_READ.read_latency_p95_ms\endcsname{3}
\expandafter\def\csname nm@liquid.O3_CONSOLIDATE.compression_ratio\endcsname{null}
\expandafter\def\csname nmc@liquid.O3_CONSOLIDATE.compression_ratio\endcsname{not defined}
\expandafter\def\csname nmn@liquid.O3_CONSOLIDATE.compression_ratio\endcsname{0}
\expandafter\def\csname nm@sint-liquid-mcp.O3_CONSOLIDATE.compression_ratio\endcsname{null}
\expandafter\def\csname nmc@sint-liquid-mcp.O3_CONSOLIDATE.compression_ratio\endcsname{not defined}
\expandafter\def\csname nmn@sint-liquid-mcp.O3_CONSOLIDATE.compression_ratio\endcsname{0}
\expandafter\def\csname nm@liquid.O3_CONSOLIDATE.probe_survival_rate\endcsname{null}
\expandafter\def\csname nmc@liquid.O3_CONSOLIDATE.probe_survival_rate\endcsname{not defined}
\expandafter\def\csname nmn@liquid.O3_CONSOLIDATE.probe_survival_rate\endcsname{0}
\expandafter\def\csname nm@sint-liquid-mcp.O3_CONSOLIDATE.probe_survival_rate\endcsname{null}
\expandafter\def\csname nmc@sint-liquid-mcp.O3_CONSOLIDATE.probe_survival_rate\endcsname{not defined}
\expandafter\def\csname nmn@sint-liquid-mcp.O3_CONSOLIDATE.probe_survival_rate\endcsname{0}
\expandafter\def\csname nm@liquid.O3_CONSOLIDATE.derived_from_completeness\endcsname{null}
\expandafter\def\csname nmc@liquid.O3_CONSOLIDATE.derived_from_completeness\endcsname{not defined}
\expandafter\def\csname nmn@liquid.O3_CONSOLIDATE.derived_from_completeness\endcsname{0}
\expandafter\def\csname nm@sint-liquid-mcp.O3_CONSOLIDATE.derived_from_completeness\endcsname{null}
\expandafter\def\csname nmc@sint-liquid-mcp.O3_CONSOLIDATE.derived_from_completeness\endcsname{not defined}
\expandafter\def\csname nmn@sint-liquid-mcp.O3_CONSOLIDATE.derived_from_completeness\endcsname{0}
\expandafter\def\csname nm@liquid.O3_CONSOLIDATE.cost_per_consolidation_usd\endcsname{null}
\expandafter\def\csname nmc@liquid.O3_CONSOLIDATE.cost_per_consolidation_usd\endcsname{not defined}
\expandafter\def\csname nmn@liquid.O3_CONSOLIDATE.cost_per_consolidation_usd\endcsname{0}
\expandafter\def\csname nm@sint-liquid-mcp.O3_CONSOLIDATE.cost_per_consolidation_usd\endcsname{null}
\expandafter\def\csname nmc@sint-liquid-mcp.O3_CONSOLIDATE.cost_per_consolidation_usd\endcsname{not defined}
\expandafter\def\csname nmn@sint-liquid-mcp.O3_CONSOLIDATE.cost_per_consolidation_usd\endcsname{0}
\expandafter\def\csname nm@liquid.O4_REVISE.contradiction_detection_f1\endcsname{0.0}
\expandafter\def\csname nmc@liquid.O4_REVISE.contradiction_detection_f1\endcsname{[0.0, 0.0]}
\expandafter\def\csname nmn@liquid.O4_REVISE.contradiction_detection_f1\endcsname{3}
\expandafter\def\csname nm@sint-liquid-mcp.O4_REVISE.contradiction_detection_f1\endcsname{0.0}
\expandafter\def\csname nmc@sint-liquid-mcp.O4_REVISE.contradiction_detection_f1\endcsname{[0.0, 0.0]}
\expandafter\def\csname nmn@sint-liquid-mcp.O4_REVISE.contradiction_detection_f1\endcsname{3}
\expandafter\def\csname nm@liquid.O4_REVISE.revision_correctness\endcsname{1.0}
\expandafter\def\csname nmc@liquid.O4_REVISE.revision_correctness\endcsname{[1.0, 1.0]}
\expandafter\def\csname nmn@liquid.O4_REVISE.revision_correctness\endcsname{3}
\expandafter\def\csname nm@sint-liquid-mcp.O4_REVISE.revision_correctness\endcsname{1.0}
\expandafter\def\csname nmc@sint-liquid-mcp.O4_REVISE.revision_correctness\endcsname{[1.0, 1.0]}
\expandafter\def\csname nmn@sint-liquid-mcp.O4_REVISE.revision_correctness\endcsname{3}
\expandafter\def\csname nm@liquid.O4_REVISE.chain_completeness\endcsname{1.0}
\expandafter\def\csname nmc@liquid.O4_REVISE.chain_completeness\endcsname{[1.0, 1.0]}
\expandafter\def\csname nmn@liquid.O4_REVISE.chain_completeness\endcsname{3}
\expandafter\def\csname nm@sint-liquid-mcp.O4_REVISE.chain_completeness\endcsname{1.0}
\expandafter\def\csname nmc@sint-liquid-mcp.O4_REVISE.chain_completeness\endcsname{[1.0, 1.0]}
\expandafter\def\csname nmn@sint-liquid-mcp.O4_REVISE.chain_completeness\endcsname{3}
\expandafter\def\csname nm@liquid.O4_REVISE.stale_read_rate\endcsname{1.0}
\expandafter\def\csname nmc@liquid.O4_REVISE.stale_read_rate\endcsname{[1.0, 1.0]}
\expandafter\def\csname nmn@liquid.O4_REVISE.stale_read_rate\endcsname{3}
\expandafter\def\csname nm@sint-liquid-mcp.O4_REVISE.stale_read_rate\endcsname{1.0}
\expandafter\def\csname nmc@sint-liquid-mcp.O4_REVISE.stale_read_rate\endcsname{[1.0, 1.0]}
\expandafter\def\csname nmn@sint-liquid-mcp.O4_REVISE.stale_read_rate\endcsname{3}
\expandafter\def\csname nm@liquid.O4_REVISE.collateral_change_rate\endcsname{0.0}
\expandafter\def\csname nmc@liquid.O4_REVISE.collateral_change_rate\endcsname{[0.0, 0.0]}
\expandafter\def\csname nmn@liquid.O4_REVISE.collateral_change_rate\endcsname{3}
\expandafter\def\csname nm@sint-liquid-mcp.O4_REVISE.collateral_change_rate\endcsname{0.0}
\expandafter\def\csname nmc@sint-liquid-mcp.O4_REVISE.collateral_change_rate\endcsname{[0.0, 0.0]}
\expandafter\def\csname nmn@sint-liquid-mcp.O4_REVISE.collateral_change_rate\endcsname{3}
\expandafter\def\csname nm@liquid.O5_INTEND.cue_hit_rate\endcsname{0.0}
\expandafter\def\csname nmc@liquid.O5_INTEND.cue_hit_rate\endcsname{[0.0, 0.0]}
\expandafter\def\csname nmn@liquid.O5_INTEND.cue_hit_rate\endcsname{3}
\expandafter\def\csname nm@sint-liquid-mcp.O5_INTEND.cue_hit_rate\endcsname{0.0}
\expandafter\def\csname nmc@sint-liquid-mcp.O5_INTEND.cue_hit_rate\endcsname{[0.0, 0.0]}
\expandafter\def\csname nmn@sint-liquid-mcp.O5_INTEND.cue_hit_rate\endcsname{3}
\expandafter\def\csname nm@liquid.O5_INTEND.false_fire_rate\endcsname{0.0}
\expandafter\def\csname nmc@liquid.O5_INTEND.false_fire_rate\endcsname{[0.0, 0.0]}
\expandafter\def\csname nmn@liquid.O5_INTEND.false_fire_rate\endcsname{3}
\expandafter\def\csname nm@sint-liquid-mcp.O5_INTEND.false_fire_rate\endcsname{0.0}
\expandafter\def\csname nmc@sint-liquid-mcp.O5_INTEND.false_fire_rate\endcsname{[0.0, 0.0]}
\expandafter\def\csname nmn@sint-liquid-mcp.O5_INTEND.false_fire_rate\endcsname{3}
\expandafter\def\csname nm@liquid.O5_INTEND.fire_latency_calls\endcsname{null}
\expandafter\def\csname nmc@liquid.O5_INTEND.fire_latency_calls\endcsname{not defined}
\expandafter\def\csname nmn@liquid.O5_INTEND.fire_latency_calls\endcsname{0}
\expandafter\def\csname nm@sint-liquid-mcp.O5_INTEND.fire_latency_calls\endcsname{null}
\expandafter\def\csname nmc@sint-liquid-mcp.O5_INTEND.fire_latency_calls\endcsname{not defined}
\expandafter\def\csname nmn@sint-liquid-mcp.O5_INTEND.fire_latency_calls\endcsname{0}
\expandafter\def\csname nm@liquid.O5_INTEND.context_availability_at_fire\endcsname{null}
\expandafter\def\csname nmc@liquid.O5_INTEND.context_availability_at_fire\endcsname{not defined}
\expandafter\def\csname nmn@liquid.O5_INTEND.context_availability_at_fire\endcsname{0}
\expandafter\def\csname nm@sint-liquid-mcp.O5_INTEND.context_availability_at_fire\endcsname{null}
\expandafter\def\csname nmc@sint-liquid-mcp.O5_INTEND.context_availability_at_fire\endcsname{not defined}
\expandafter\def\csname nmn@sint-liquid-mcp.O5_INTEND.context_availability_at_fire\endcsname{0}
\expandafter\def\csname nm@liquid.O6_ASSERT.abstention_correctness\endcsname{1.0}
\expandafter\def\csname nmc@liquid.O6_ASSERT.abstention_correctness\endcsname{[1.0, 1.0]}
\expandafter\def\csname nmn@liquid.O6_ASSERT.abstention_correctness\endcsname{3}
\expandafter\def\csname nm@sint-liquid-mcp.O6_ASSERT.abstention_correctness\endcsname{1.0}
\expandafter\def\csname nmc@sint-liquid-mcp.O6_ASSERT.abstention_correctness\endcsname{[1.0, 1.0]}
\expandafter\def\csname nmn@sint-liquid-mcp.O6_ASSERT.abstention_correctness\endcsname{3}
\expandafter\def\csname nm@liquid.O6_ASSERT.in_coverage_answer_correctness\endcsname{0.5}
\expandafter\def\csname nmc@liquid.O6_ASSERT.in_coverage_answer_correctness\endcsname{[0.5, 0.5]}
\expandafter\def\csname nmn@liquid.O6_ASSERT.in_coverage_answer_correctness\endcsname{3}
\expandafter\def\csname nm@sint-liquid-mcp.O6_ASSERT.in_coverage_answer_correctness\endcsname{0.5}
\expandafter\def\csname nmc@sint-liquid-mcp.O6_ASSERT.in_coverage_answer_correctness\endcsname{[0.5, 0.5]}
\expandafter\def\csname nmn@sint-liquid-mcp.O6_ASSERT.in_coverage_answer_correctness\endcsname{3}
\expandafter\def\csname nm@liquid.O6_ASSERT.key_question_correctness\endcsname{1.0}
\expandafter\def\csname nmc@liquid.O6_ASSERT.key_question_correctness\endcsname{[1.0, 1.0]}
\expandafter\def\csname nmn@liquid.O6_ASSERT.key_question_correctness\endcsname{3}
\expandafter\def\csname nm@sint-liquid-mcp.O6_ASSERT.key_question_correctness\endcsname{1.0}
\expandafter\def\csname nmc@sint-liquid-mcp.O6_ASSERT.key_question_correctness\endcsname{[1.0, 1.0]}
\expandafter\def\csname nmn@sint-liquid-mcp.O6_ASSERT.key_question_correctness\endcsname{3}
\expandafter\def\csname nm@liquid.O6_ASSERT.subject_question_correctness\endcsname{0.0}
\expandafter\def\csname nmc@liquid.O6_ASSERT.subject_question_correctness\endcsname{[0.0, 0.0]}
\expandafter\def\csname nmn@liquid.O6_ASSERT.subject_question_correctness\endcsname{3}
\expandafter\def\csname nm@sint-liquid-mcp.O6_ASSERT.subject_question_correctness\endcsname{0.0}
\expandafter\def\csname nmc@sint-liquid-mcp.O6_ASSERT.subject_question_correctness\endcsname{[0.0, 0.0]}
\expandafter\def\csname nmn@sint-liquid-mcp.O6_ASSERT.subject_question_correctness\endcsname{3}
\expandafter\def\csname nm@liquid.O6_ASSERT.ece\endcsname{0.3667}
\expandafter\def\csname nmc@liquid.O6_ASSERT.ece\endcsname{[0.3667, 0.3667]}
\expandafter\def\csname nmn@liquid.O6_ASSERT.ece\endcsname{3}
\expandafter\def\csname nm@sint-liquid-mcp.O6_ASSERT.ece\endcsname{0.3667}
\expandafter\def\csname nmc@sint-liquid-mcp.O6_ASSERT.ece\endcsname{[0.3667, 0.3667]}
\expandafter\def\csname nmn@sint-liquid-mcp.O6_ASSERT.ece\endcsname{3}
\expandafter\def\csname nm@liquid.O6_ASSERT.discrimination_auroc\endcsname{0.75}
\expandafter\def\csname nmc@liquid.O6_ASSERT.discrimination_auroc\endcsname{[0.75, 0.75]}
\expandafter\def\csname nmn@liquid.O6_ASSERT.discrimination_auroc\endcsname{3}
\expandafter\def\csname nm@sint-liquid-mcp.O6_ASSERT.discrimination_auroc\endcsname{0.75}
\expandafter\def\csname nmc@sint-liquid-mcp.O6_ASSERT.discrimination_auroc\endcsname{[0.75, 0.75]}
\expandafter\def\csname nmn@sint-liquid-mcp.O6_ASSERT.discrimination_auroc\endcsname{3}
\expandafter\def\csname nm@liquid.O6_ASSERT.monitoring_accuracy\endcsname{1.0}
\expandafter\def\csname nmc@liquid.O6_ASSERT.monitoring_accuracy\endcsname{[1.0, 1.0]}
\expandafter\def\csname nmn@liquid.O6_ASSERT.monitoring_accuracy\endcsname{3}
\expandafter\def\csname nm@sint-liquid-mcp.O6_ASSERT.monitoring_accuracy\endcsname{1.0}
\expandafter\def\csname nmc@sint-liquid-mcp.O6_ASSERT.monitoring_accuracy\endcsname{[1.0, 1.0]}
\expandafter\def\csname nmn@sint-liquid-mcp.O6_ASSERT.monitoring_accuracy\endcsname{3}
\expandafter\def\csname nm@liquid.PROVENANCE.provenance_correctness\endcsname{1.0}
\expandafter\def\csname nmc@liquid.PROVENANCE.provenance_correctness\endcsname{[1.0, 1.0]}
\expandafter\def\csname nmn@liquid.PROVENANCE.provenance_correctness\endcsname{3}
\expandafter\def\csname nm@sint-liquid-mcp.PROVENANCE.provenance_correctness\endcsname{1.0}
\expandafter\def\csname nmc@sint-liquid-mcp.PROVENANCE.provenance_correctness\endcsname{[1.0, 1.0]}
\expandafter\def\csname nmn@sint-liquid-mcp.PROVENANCE.provenance_correctness\endcsname{3}
\expandafter\def\csname nm@liquid.PROVENANCE.derived_from_completeness\endcsname{null}
\expandafter\def\csname nmc@liquid.PROVENANCE.derived_from_completeness\endcsname{not defined}
\expandafter\def\csname nmn@liquid.PROVENANCE.derived_from_completeness\endcsname{0}
\expandafter\def\csname nm@sint-liquid-mcp.PROVENANCE.derived_from_completeness\endcsname{null}
\expandafter\def\csname nmc@sint-liquid-mcp.PROVENANCE.derived_from_completeness\endcsname{not defined}
\expandafter\def\csname nmn@sint-liquid-mcp.PROVENANCE.derived_from_completeness\endcsname{0}
\expandafter\def\csname nm@liquid.LIFECYCLE.lifecycle_state_exposed\endcsname{0.0}
\expandafter\def\csname nmc@liquid.LIFECYCLE.lifecycle_state_exposed\endcsname{[0.0, 0.0]}
\expandafter\def\csname nmn@liquid.LIFECYCLE.lifecycle_state_exposed\endcsname{3}
\expandafter\def\csname nm@sint-liquid-mcp.LIFECYCLE.lifecycle_state_exposed\endcsname{0.0}
\expandafter\def\csname nmc@sint-liquid-mcp.LIFECYCLE.lifecycle_state_exposed\endcsname{[0.0, 0.0]}
\expandafter\def\csname nmn@sint-liquid-mcp.LIFECYCLE.lifecycle_state_exposed\endcsname{3}
\expandafter\def\csname nm@liquid.LIFECYCLE.admissibility_error_rate\endcsname{0.5}
\expandafter\def\csname nmc@liquid.LIFECYCLE.admissibility_error_rate\endcsname{[0.5, 0.5]}
\expandafter\def\csname nmn@liquid.LIFECYCLE.admissibility_error_rate\endcsname{3}
\expandafter\def\csname nm@sint-liquid-mcp.LIFECYCLE.admissibility_error_rate\endcsname{0.5}
\expandafter\def\csname nmc@sint-liquid-mcp.LIFECYCLE.admissibility_error_rate\endcsname{[0.5, 0.5]}
\expandafter\def\csname nmn@sint-liquid-mcp.LIFECYCLE.admissibility_error_rate\endcsname{3}
\expandafter\def\csname nm@liquid.LIFECYCLE.stale_read_rate\endcsname{1.0}
\expandafter\def\csname nmc@liquid.LIFECYCLE.stale_read_rate\endcsname{[1.0, 1.0]}
\expandafter\def\csname nmn@liquid.LIFECYCLE.stale_read_rate\endcsname{3}
\expandafter\def\csname nm@sint-liquid-mcp.LIFECYCLE.stale_read_rate\endcsname{1.0}
\expandafter\def\csname nmc@sint-liquid-mcp.LIFECYCLE.stale_read_rate\endcsname{[1.0, 1.0]}
\expandafter\def\csname nmn@sint-liquid-mcp.LIFECYCLE.stale_read_rate\endcsname{3}
\expandafter\def\csname nm@liquid.LIFECYCLE.restart_survival\endcsname{1.0}
\expandafter\def\csname nmc@liquid.LIFECYCLE.restart_survival\endcsname{[1.0, 1.0]}
\expandafter\def\csname nmn@liquid.LIFECYCLE.restart_survival\endcsname{3}
\expandafter\def\csname nm@sint-liquid-mcp.LIFECYCLE.restart_survival\endcsname{1.0}
\expandafter\def\csname nmc@sint-liquid-mcp.LIFECYCLE.restart_survival\endcsname{[1.0, 1.0]}
\expandafter\def\csname nmn@sint-liquid-mcp.LIFECYCLE.restart_survival\endcsname{3}
\expandafter\def\csname nm@liquid.LIFECYCLE.snapshot_reopen_survival\endcsname{1.0}
\expandafter\def\csname nmc@liquid.LIFECYCLE.snapshot_reopen_survival\endcsname{[1.0, 1.0]}
\expandafter\def\csname nmn@liquid.LIFECYCLE.snapshot_reopen_survival\endcsname{3}
\expandafter\def\csname nm@sint-liquid-mcp.LIFECYCLE.snapshot_reopen_survival\endcsname{1.0}
\expandafter\def\csname nmc@sint-liquid-mcp.LIFECYCLE.snapshot_reopen_survival\endcsname{[1.0, 1.0]}
\expandafter\def\csname nmn@sint-liquid-mcp.LIFECYCLE.snapshot_reopen_survival\endcsname{3}
\expandafter\def\csname nm@liquid.metrics_total\endcsname{52}
\expandafter\def\csname nm@liquid.metrics_valued\endcsname{41}
\expandafter\def\csname nm@liquid.metrics_null\endcsname{11}
\expandafter\def\csname nm@liquid.repetitions\endcsname{3}
\expandafter\def\csname nmt@liquid.seeds\endcsname{1, 2, 3}
\expandafter\def\csname nmt@liquid.ci_method\endcsname{bootstrap\_percentile}
\expandafter\def\csname nm@liquid.ci_resamples\endcsname{1000}
\expandafter\def\csname nm@liquid.calls_issued\endcsname{681}
\expandafter\def\csname nm@liquid.tools_advertised\endcsname{8}
\expandafter\def\csname nm@liquid.raw_files\endcsname{33}
\expandafter\def\csname nm@liquid.failures.not_implemented\endcsname{15}
\expandafter\def\csname nm@liquid.T1.tasks\endcsname{12}
\expandafter\def\csname nm@liquid.T2.tasks\endcsname{6}
\expandafter\def\csname nm@liquid.T3.tasks\endcsname{4}
\expandafter\def\csname nm@liquid.T4.tasks\endcsname{6}
\expandafter\def\csname nm@liquid.T5.tasks\endcsname{4}
\expandafter\def\csname nm@liquid.T6.tasks\endcsname{4}
\expandafter\def\csname nm@liquid.T7.tasks\endcsname{3}
\expandafter\def\csname nm@liquid.T8.tasks\endcsname{12}
\expandafter\def\csname nm@liquid.T9.tasks\endcsname{18}
\expandafter\def\csname nm@liquid.T10.tasks\endcsname{12}
\expandafter\def\csname nmt@liquid.O1_WRITE\endcsname{partial}
\expandafter\def\csname nmt@liquid.O2_READ\endcsname{yes}
\expandafter\def\csname nmt@liquid.O3_CONSOLIDATE\endcsname{no}
\expandafter\def\csname nmt@liquid.O4_REVISE\endcsname{yes}
\expandafter\def\csname nmt@liquid.O5_INTEND\endcsname{partial}
\expandafter\def\csname nmt@liquid.O6_ASSERT\endcsname{yes}
\expandafter\def\csname nmt@liquid.PROVENANCE\endcsname{partial}
\expandafter\def\csname nmt@liquid.LIFECYCLE\endcsname{yes}
\expandafter\def\csname nmt@liquid.identity.subject\endcsname{sint-liquid-mcp}
\expandafter\def\csname nmtl@liquid.identity.subject\endcsname{sint-liqui\hspace{0pt}d-mcp}
\expandafter\def\csname nmt@liquid.identity.repository\endcsname{https://github.com/kostyk348/sint-liquid-mcp}
\expandafter\def\csname nmtl@liquid.identity.repository\endcsname{https:\hspace{0pt}/\hspace{0pt}/\hspace{0pt}github.com\hspace{0pt}/\hspace{0pt}kostyk348/\hspace{0pt}sint-liqui\hspace{0pt}d-mcp}
\expandafter\def\csname nmt@liquid.identity.commit\endcsname{bfffce2cdec35c11aa0d90ae02689d69df596b11}
\expandafter\def\csname nmtl@liquid.identity.commit\endcsname{bfffce2cde\hspace{0pt}c35c11aa0d\hspace{0pt}90ae02689d\hspace{0pt}69df596b11\hspace{0pt}}
\expandafter\def\csname nmt@liquid.identity.published_version\endcsname{MCP server as published in the repository snapshot (no release tag)}
\expandafter\def\csname nmtl@liquid.identity.published_version\endcsname{MCP server as published in the repository\hspace{0pt} snapshot (no release tag)}
\expandafter\def\csname nmt@liquid.identity.transport___MCP_revision\endcsname{stdio / 2025-03-26}
\expandafter\def\csname nmtl@liquid.identity.transport___MCP_revision\endcsname{stdio /\hspace{0pt} 2025-03-26\hspace{0pt}}
\expandafter\def\csname nmt@liquid.identity.runtime\endcsname{Python 3.14.7, mcp 1.9.4 (mcp.server.fastmcp)}
\expandafter\def\csname nmtl@liquid.identity.runtime\endcsname{Python 3.14.7,\hspace{0pt} mcp 1.9.4 (mcp.serve\hspace{0pt}r.fastmcp)\hspace{0pt}}
\expandafter\def\csname nmt@liquid.identity.deployment\endcsname{local process, fresh data directory per repetition (USERPROFILE/HOME redirected)}
\expandafter\def\csname nmtl@liquid.identity.deployment\endcsname{local process,\hspace{0pt} fresh data directory per repetition\hspace{0pt} (USERPROFI\hspace{0pt}LE/\hspace{0pt}HOME redirected\hspace{0pt})}
\expandafter\def\csname nmt@liquid.identity.store_kept_in_the_run_dir\endcsname{reports/sint-liquid-mcp-2026-10-16/data/seed<seed>}
\expandafter\def\csname nmtl@liquid.identity.store_kept_in_the_run_dir\endcsname{reports/\hspace{0pt}sint-liqui\hspace{0pt}d-mcp-2026\hspace{0pt}-10-16/\hspace{0pt}data/\hspace{0pt}seed<seed>\hspace{0pt}}
\expandafter\def\csname nmt@liquid.identity.model_under_test\endcsname{none (store-only)}
\expandafter\def\csname nmtl@liquid.identity.model_under_test\endcsname{none (store-onl\hspace{0pt}y)}
\expandafter\def\csname nmt@liquid.identity.hardware\endcsname{AMD64, 4 logical CPUs (workstation class, not a benchmark host)}
\expandafter\def\csname nmtl@liquid.identity.hardware\endcsname{AMD64,\hspace{0pt} 4 logical CPUs (workstati\hspace{0pt}on class,\hspace{0pt} not a benchmark host)}
\expandafter\def\csname nmt@liquid.identity.runner___harness\endcsname{uncommitted working tree: sha256:9b79d0aa12ae1660 over code/run\_protocol.py + code/streams.py + code/servers.json (this workspace has no git metadata, so there is no commit to name)}
\expandafter\def\csname nmtl@liquid.identity.runner___harness\endcsname{uncommitte\hspace{0pt}d working tree:\hspace{0pt} sha256:\hspace{0pt}9b79d0aa12\hspace{0pt}ae1660 over code/\hspace{0pt}run\_\hspace{0pt}protocol.p\hspace{0pt}y + code/\hspace{0pt}streams.py\hspace{0pt} + code/\hspace{0pt}servers.js\hspace{0pt}on (this workspace has no git metadata,\hspace{0pt} so there is no commit to name)}
\expandafter\def\csname nmt@liquid.identity.environment\endcsname{Windows-11-10.0.26200-SP0; Intel64 Family 6 Model 140 Stepping 2, GenuineIntel; server interpreter <hermes-venv>/Scripts/python.exe; harness Python 3.14.7 (stdlib only)}
\expandafter\def\csname nmtl@liquid.identity.environment\endcsname{Windows-11\hspace{0pt}-10.0.2620\hspace{0pt}0-SP0;\hspace{0pt} Intel64 Family 6 Model 140 Stepping 2,\hspace{0pt} GenuineInt\hspace{0pt}el;\hspace{0pt} server interprete\hspace{0pt}r C:\hspace{0pt}/\hspace{0pt}Users/\hspace{0pt}Lab/\hspace{0pt}AppData/\hspace{0pt}Local/\hspace{0pt}hermes/\hspace{0pt}installs/\hspace{0pt}2636d87a16\hspace{0pt}39ae64/\hspace{0pt}environmen\hspace{0pt}ts/\hspace{0pt}f32597c25d\hspace{0pt}734363b2d2\hspace{0pt}20b038061d\hspace{0pt}05/\hspace{0pt}venv/\hspace{0pt}Scripts/\hspace{0pt}python.exe\hspace{0pt};\hspace{0pt} harness Python 3.14.7 (stdlib only)}
\expandafter\def\csname nmt@liquid.identity.repetitions___seeds\endcsname{3 / 1, 2, 3}
\expandafter\def\csname nmtl@liquid.identity.repetitions___seeds\endcsname{3 /\hspace{0pt} 1,\hspace{0pt} 2,\hspace{0pt} 3}
\expandafter\def\csname nmt@liquid.identity.budgets_per_repetition\endcsname{400 calls, 900 s wall clock}
\expandafter\def\csname nmtl@liquid.identity.budgets_per_repetition\endcsname{400 calls,\hspace{0pt} 900 s wall clock}
\expandafter\def\csname nmt@liquid.identity.confidence_intervals\endcsname{bootstrap\_percentile, 1000 resamples}
\expandafter\def\csname nmtl@liquid.identity.confidence_intervals\endcsname{bootstrap\_\hspace{0pt}percentile\hspace{0pt},\hspace{0pt} 1000 resamples}
\expandafter\def\csname nmt@liquid.identity.failures_recorded\endcsname{not\_implemented x 15}
\expandafter\def\csname nmtl@liquid.identity.failures_recorded\endcsname{not\_\hspace{0pt}implemente\hspace{0pt}d x 15}
\expandafter\def\csname nmt@liquid.identity.raw_responses\endcsname{reports/sint-liquid-mcp-2026-10-16/raw (33 JSONL files)}
\expandafter\def\csname nmtl@liquid.identity.raw_responses\endcsname{reports/\hspace{0pt}sint-liqui\hspace{0pt}d-mcp-2026\hspace{0pt}-10-16/\hspace{0pt}raw (33 JSONL files)}
\expandafter\def\csname nmt@liquid.identity.report.json\endcsname{reports/sint-liquid-mcp-2026-10-16/report.json}
\expandafter\def\csname nmtl@liquid.identity.report.json\endcsname{reports/\hspace{0pt}sint-liqui\hspace{0pt}d-mcp-2026\hspace{0pt}-10-16/\hspace{0pt}report.jso\hspace{0pt}n}
\expandafter\def\csname nm@memory.T1.write_precision\endcsname{0.75}
\expandafter\def\csname nmc@memory.T1.write_precision\endcsname{[0.75, 0.75]}
\expandafter\def\csname nmn@memory.T1.write_precision\endcsname{3}
\expandafter\def\csname nm@sint-memory-mcp.T1.write_precision\endcsname{0.75}
\expandafter\def\csname nmc@sint-memory-mcp.T1.write_precision\endcsname{[0.75, 0.75]}
\expandafter\def\csname nmn@sint-memory-mcp.T1.write_precision\endcsname{3}
\expandafter\def\csname nm@memory.T1.write_recall\endcsname{1.0}
\expandafter\def\csname nmc@memory.T1.write_recall\endcsname{[1.0, 1.0]}
\expandafter\def\csname nmn@memory.T1.write_recall\endcsname{3}
\expandafter\def\csname nm@sint-memory-mcp.T1.write_recall\endcsname{1.0}
\expandafter\def\csname nmc@sint-memory-mcp.T1.write_recall\endcsname{[1.0, 1.0]}
\expandafter\def\csname nmn@sint-memory-mcp.T1.write_recall\endcsname{3}
\expandafter\def\csname nm@memory.T1.duplicate_rate\endcsname{0.25}
\expandafter\def\csname nmc@memory.T1.duplicate_rate\endcsname{[0.25, 0.25]}
\expandafter\def\csname nmn@memory.T1.duplicate_rate\endcsname{3}
\expandafter\def\csname nm@sint-memory-mcp.T1.duplicate_rate\endcsname{0.25}
\expandafter\def\csname nmc@sint-memory-mcp.T1.duplicate_rate\endcsname{[0.25, 0.25]}
\expandafter\def\csname nmn@sint-memory-mcp.T1.duplicate_rate\endcsname{3}
\expandafter\def\csname nm@memory.T1.idempotency_violations\endcsname{4.0}
\expandafter\def\csname nmc@memory.T1.idempotency_violations\endcsname{[4.0, 4.0]}
\expandafter\def\csname nmn@memory.T1.idempotency_violations\endcsname{3}
\expandafter\def\csname nm@sint-memory-mcp.T1.idempotency_violations\endcsname{4.0}
\expandafter\def\csname nmc@sint-memory-mcp.T1.idempotency_violations\endcsname{[4.0, 4.0]}
\expandafter\def\csname nmn@sint-memory-mcp.T1.idempotency_violations\endcsname{3}
\expandafter\def\csname nm@memory.T1.write_latency_p50_ms\endcsname{16.4423}
\expandafter\def\csname nmc@memory.T1.write_latency_p50_ms\endcsname{[15.167, 17.2955]}
\expandafter\def\csname nmn@memory.T1.write_latency_p50_ms\endcsname{3}
\expandafter\def\csname nm@sint-memory-mcp.T1.write_latency_p50_ms\endcsname{16.4423}
\expandafter\def\csname nmc@sint-memory-mcp.T1.write_latency_p50_ms\endcsname{[15.167, 17.2955]}
\expandafter\def\csname nmn@sint-memory-mcp.T1.write_latency_p50_ms\endcsname{3}
\expandafter\def\csname nm@memory.T1.write_latency_p95_ms\endcsname{18.9019}
\expandafter\def\csname nmc@memory.T1.write_latency_p95_ms\endcsname{[17.329, 20.2025]}
\expandafter\def\csname nmn@memory.T1.write_latency_p95_ms\endcsname{3}
\expandafter\def\csname nm@sint-memory-mcp.T1.write_latency_p95_ms\endcsname{18.9019}
\expandafter\def\csname nmc@sint-memory-mcp.T1.write_latency_p95_ms\endcsname{[17.329, 20.2025]}
\expandafter\def\csname nmn@sint-memory-mcp.T1.write_latency_p95_ms\endcsname{3}
\expandafter\def\csname nm@memory.T1.tokens_per_write\endcsname{null}
\expandafter\def\csname nmc@memory.T1.tokens_per_write\endcsname{not defined}
\expandafter\def\csname nmn@memory.T1.tokens_per_write\endcsname{0}
\expandafter\def\csname nm@sint-memory-mcp.T1.tokens_per_write\endcsname{null}
\expandafter\def\csname nmc@sint-memory-mcp.T1.tokens_per_write\endcsname{not defined}
\expandafter\def\csname nmn@sint-memory-mcp.T1.tokens_per_write\endcsname{0}
\expandafter\def\csname nm@memory.T1.cost_per_1000_writes_usd\endcsname{null}
\expandafter\def\csname nmc@memory.T1.cost_per_1000_writes_usd\endcsname{not defined}
\expandafter\def\csname nmn@memory.T1.cost_per_1000_writes_usd\endcsname{0}
\expandafter\def\csname nm@sint-memory-mcp.T1.cost_per_1000_writes_usd\endcsname{null}
\expandafter\def\csname nmc@sint-memory-mcp.T1.cost_per_1000_writes_usd\endcsname{not defined}
\expandafter\def\csname nmn@sint-memory-mcp.T1.cost_per_1000_writes_usd\endcsname{0}
\expandafter\def\csname nm@memory.T2.recall_at_5\endcsname{1.0}
\expandafter\def\csname nmc@memory.T2.recall_at_5\endcsname{[1.0, 1.0]}
\expandafter\def\csname nmn@memory.T2.recall_at_5\endcsname{3}
\expandafter\def\csname nm@sint-memory-mcp.T2.recall_at_5\endcsname{1.0}
\expandafter\def\csname nmc@sint-memory-mcp.T2.recall_at_5\endcsname{[1.0, 1.0]}
\expandafter\def\csname nmn@sint-memory-mcp.T2.recall_at_5\endcsname{3}
\expandafter\def\csname nm@memory.T2.required_record_hit_rate\endcsname{1.0}
\expandafter\def\csname nmc@memory.T2.required_record_hit_rate\endcsname{[1.0, 1.0]}
\expandafter\def\csname nmn@memory.T2.required_record_hit_rate\endcsname{3}
\expandafter\def\csname nm@sint-memory-mcp.T2.required_record_hit_rate\endcsname{1.0}
\expandafter\def\csname nmc@sint-memory-mcp.T2.required_record_hit_rate\endcsname{[1.0, 1.0]}
\expandafter\def\csname nmn@sint-memory-mcp.T2.required_record_hit_rate\endcsname{3}
\expandafter\def\csname nm@memory.T2.mrr\endcsname{0.75}
\expandafter\def\csname nmc@memory.T2.mrr\endcsname{[0.6806, 0.8056]}
\expandafter\def\csname nmn@memory.T2.mrr\endcsname{3}
\expandafter\def\csname nm@sint-memory-mcp.T2.mrr\endcsname{0.75}
\expandafter\def\csname nmc@sint-memory-mcp.T2.mrr\endcsname{[0.6806, 0.8056]}
\expandafter\def\csname nmn@sint-memory-mcp.T2.mrr\endcsname{3}
\expandafter\def\csname nm@memory.T2.key_query_hit_rate\endcsname{1.0}
\expandafter\def\csname nmc@memory.T2.key_query_hit_rate\endcsname{[1.0, 1.0]}
\expandafter\def\csname nmn@memory.T2.key_query_hit_rate\endcsname{3}
\expandafter\def\csname nm@sint-memory-mcp.T2.key_query_hit_rate\endcsname{1.0}
\expandafter\def\csname nmc@sint-memory-mcp.T2.key_query_hit_rate\endcsname{[1.0, 1.0]}
\expandafter\def\csname nmn@sint-memory-mcp.T2.key_query_hit_rate\endcsname{3}
\expandafter\def\csname nm@memory.T2.topic_query_hit_rate\endcsname{1.0}
\expandafter\def\csname nmc@memory.T2.topic_query_hit_rate\endcsname{[1.0, 1.0]}
\expandafter\def\csname nmn@memory.T2.topic_query_hit_rate\endcsname{3}
\expandafter\def\csname nm@sint-memory-mcp.T2.topic_query_hit_rate\endcsname{1.0}
\expandafter\def\csname nmc@sint-memory-mcp.T2.topic_query_hit_rate\endcsname{[1.0, 1.0]}
\expandafter\def\csname nmn@sint-memory-mcp.T2.topic_query_hit_rate\endcsname{3}
\expandafter\def\csname nm@memory.T2.provenance_correctness\endcsname{1.0}
\expandafter\def\csname nmc@memory.T2.provenance_correctness\endcsname{[1.0, 1.0]}
\expandafter\def\csname nmn@memory.T2.provenance_correctness\endcsname{3}
\expandafter\def\csname nm@sint-memory-mcp.T2.provenance_correctness\endcsname{1.0}
\expandafter\def\csname nmc@sint-memory-mcp.T2.provenance_correctness\endcsname{[1.0, 1.0]}
\expandafter\def\csname nmn@sint-memory-mcp.T2.provenance_correctness\endcsname{3}
\expandafter\def\csname nm@memory.T2.context_inflation\endcsname{2.5}
\expandafter\def\csname nmc@memory.T2.context_inflation\endcsname{[2.5, 2.5]}
\expandafter\def\csname nmn@memory.T2.context_inflation\endcsname{3}
\expandafter\def\csname nm@sint-memory-mcp.T2.context_inflation\endcsname{2.5}
\expandafter\def\csname nmc@sint-memory-mcp.T2.context_inflation\endcsname{[2.5, 2.5]}
\expandafter\def\csname nmn@sint-memory-mcp.T2.context_inflation\endcsname{3}
\expandafter\def\csname nm@memory.T2.read_latency_p50_ms\endcsname{8.0937}
\expandafter\def\csname nmc@memory.T2.read_latency_p50_ms\endcsname{[7.6565, 8.5435]}
\expandafter\def\csname nmn@memory.T2.read_latency_p50_ms\endcsname{3}
\expandafter\def\csname nm@sint-memory-mcp.T2.read_latency_p50_ms\endcsname{8.0937}
\expandafter\def\csname nmc@sint-memory-mcp.T2.read_latency_p50_ms\endcsname{[7.6565, 8.5435]}
\expandafter\def\csname nmn@sint-memory-mcp.T2.read_latency_p50_ms\endcsname{3}
\expandafter\def\csname nm@memory.T2.read_latency_p95_ms\endcsname{102.8932}
\expandafter\def\csname nmc@memory.T2.read_latency_p95_ms\endcsname{[101.5778, 104.9318]}
\expandafter\def\csname nmn@memory.T2.read_latency_p95_ms\endcsname{3}
\expandafter\def\csname nm@sint-memory-mcp.T2.read_latency_p95_ms\endcsname{102.8932}
\expandafter\def\csname nmc@sint-memory-mcp.T2.read_latency_p95_ms\endcsname{[101.5778, 104.9318]}
\expandafter\def\csname nmn@sint-memory-mcp.T2.read_latency_p95_ms\endcsname{3}
\expandafter\def\csname nm@memory.T3.temporal_correctness\endcsname{null}
\expandafter\def\csname nmc@memory.T3.temporal_correctness\endcsname{not defined}
\expandafter\def\csname nmn@memory.T3.temporal_correctness\endcsname{0}
\expandafter\def\csname nm@sint-memory-mcp.T3.temporal_correctness\endcsname{null}
\expandafter\def\csname nmc@sint-memory-mcp.T3.temporal_correctness\endcsname{not defined}
\expandafter\def\csname nmn@sint-memory-mcp.T3.temporal_correctness\endcsname{0}
\expandafter\def\csname nm@memory.T3.temporal_metadata_exposed\endcsname{1.0}
\expandafter\def\csname nmc@memory.T3.temporal_metadata_exposed\endcsname{[1.0, 1.0]}
\expandafter\def\csname nmn@memory.T3.temporal_metadata_exposed\endcsname{3}
\expandafter\def\csname nm@sint-memory-mcp.T3.temporal_metadata_exposed\endcsname{1.0}
\expandafter\def\csname nmc@sint-memory-mcp.T3.temporal_metadata_exposed\endcsname{[1.0, 1.0]}
\expandafter\def\csname nmn@sint-memory-mcp.T3.temporal_metadata_exposed\endcsname{3}
\expandafter\def\csname nm@memory.T4.admissibility_error_rate\endcsname{0.5}
\expandafter\def\csname nmc@memory.T4.admissibility_error_rate\endcsname{[0.5, 0.5]}
\expandafter\def\csname nmn@memory.T4.admissibility_error_rate\endcsname{3}
\expandafter\def\csname nm@sint-memory-mcp.T4.admissibility_error_rate\endcsname{0.5}
\expandafter\def\csname nmc@sint-memory-mcp.T4.admissibility_error_rate\endcsname{[0.5, 0.5]}
\expandafter\def\csname nmn@sint-memory-mcp.T4.admissibility_error_rate\endcsname{3}
\expandafter\def\csname nm@memory.T4.lifecycle_state_exposed\endcsname{0.0}
\expandafter\def\csname nmc@memory.T4.lifecycle_state_exposed\endcsname{[0.0, 0.0]}
\expandafter\def\csname nmn@memory.T4.lifecycle_state_exposed\endcsname{3}
\expandafter\def\csname nm@sint-memory-mcp.T4.lifecycle_state_exposed\endcsname{0.0}
\expandafter\def\csname nmc@sint-memory-mcp.T4.lifecycle_state_exposed\endcsname{[0.0, 0.0]}
\expandafter\def\csname nmn@sint-memory-mcp.T4.lifecycle_state_exposed\endcsname{3}
\expandafter\def\csname nm@memory.T4.principal_scoping_exposed\endcsname{0.0}
\expandafter\def\csname nmc@memory.T4.principal_scoping_exposed\endcsname{[0.0, 0.0]}
\expandafter\def\csname nmn@memory.T4.principal_scoping_exposed\endcsname{3}
\expandafter\def\csname nm@sint-memory-mcp.T4.principal_scoping_exposed\endcsname{0.0}
\expandafter\def\csname nmc@sint-memory-mcp.T4.principal_scoping_exposed\endcsname{[0.0, 0.0]}
\expandafter\def\csname nmn@sint-memory-mcp.T4.principal_scoping_exposed\endcsname{3}
\expandafter\def\csname nm@memory.T5.contradiction_detection_f1\endcsname{0.0}
\expandafter\def\csname nmc@memory.T5.contradiction_detection_f1\endcsname{[0.0, 0.0]}
\expandafter\def\csname nmn@memory.T5.contradiction_detection_f1\endcsname{3}
\expandafter\def\csname nm@sint-memory-mcp.T5.contradiction_detection_f1\endcsname{0.0}
\expandafter\def\csname nmc@sint-memory-mcp.T5.contradiction_detection_f1\endcsname{[0.0, 0.0]}
\expandafter\def\csname nmn@sint-memory-mcp.T5.contradiction_detection_f1\endcsname{3}
\expandafter\def\csname nm@memory.T5.revision_correctness\endcsname{null}
\expandafter\def\csname nmc@memory.T5.revision_correctness\endcsname{not defined}
\expandafter\def\csname nmn@memory.T5.revision_correctness\endcsname{0}
\expandafter\def\csname nm@sint-memory-mcp.T5.revision_correctness\endcsname{null}
\expandafter\def\csname nmc@sint-memory-mcp.T5.revision_correctness\endcsname{not defined}
\expandafter\def\csname nmn@sint-memory-mcp.T5.revision_correctness\endcsname{0}
\expandafter\def\csname nm@memory.T5.chain_completeness\endcsname{0.0}
\expandafter\def\csname nmc@memory.T5.chain_completeness\endcsname{[0.0, 0.0]}
\expandafter\def\csname nmn@memory.T5.chain_completeness\endcsname{3}
\expandafter\def\csname nm@sint-memory-mcp.T5.chain_completeness\endcsname{0.0}
\expandafter\def\csname nmc@sint-memory-mcp.T5.chain_completeness\endcsname{[0.0, 0.0]}
\expandafter\def\csname nmn@sint-memory-mcp.T5.chain_completeness\endcsname{3}
\expandafter\def\csname nm@memory.T5.stale_read_rate\endcsname{1.0}
\expandafter\def\csname nmc@memory.T5.stale_read_rate\endcsname{[1.0, 1.0]}
\expandafter\def\csname nmn@memory.T5.stale_read_rate\endcsname{3}
\expandafter\def\csname nm@sint-memory-mcp.T5.stale_read_rate\endcsname{1.0}
\expandafter\def\csname nmc@sint-memory-mcp.T5.stale_read_rate\endcsname{[1.0, 1.0]}
\expandafter\def\csname nmn@sint-memory-mcp.T5.stale_read_rate\endcsname{3}
\expandafter\def\csname nm@memory.T5.both_competing_contents_served_rate\endcsname{1.0}
\expandafter\def\csname nmc@memory.T5.both_competing_contents_served_rate\endcsname{[1.0, 1.0]}
\expandafter\def\csname nmn@memory.T5.both_competing_contents_served_rate\endcsname{3}
\expandafter\def\csname nm@sint-memory-mcp.T5.both_competing_contents_served_rate\endcsname{1.0}
\expandafter\def\csname nmc@sint-memory-mcp.T5.both_competing_contents_served_rate\endcsname{[1.0, 1.0]}
\expandafter\def\csname nmn@sint-memory-mcp.T5.both_competing_contents_served_rate\endcsname{3}
\expandafter\def\csname nm@memory.T5.policy_winner_top_ranked_rate\endcsname{0.5}
\expandafter\def\csname nmc@memory.T5.policy_winner_top_ranked_rate\endcsname{[0.5, 0.5]}
\expandafter\def\csname nmn@memory.T5.policy_winner_top_ranked_rate\endcsname{3}
\expandafter\def\csname nm@sint-memory-mcp.T5.policy_winner_top_ranked_rate\endcsname{0.5}
\expandafter\def\csname nmc@sint-memory-mcp.T5.policy_winner_top_ranked_rate\endcsname{[0.5, 0.5]}
\expandafter\def\csname nmn@sint-memory-mcp.T5.policy_winner_top_ranked_rate\endcsname{3}
\expandafter\def\csname nm@memory.T5.collateral_change_rate\endcsname{0.0}
\expandafter\def\csname nmc@memory.T5.collateral_change_rate\endcsname{[0.0, 0.0]}
\expandafter\def\csname nmn@memory.T5.collateral_change_rate\endcsname{3}
\expandafter\def\csname nm@sint-memory-mcp.T5.collateral_change_rate\endcsname{0.0}
\expandafter\def\csname nmc@sint-memory-mcp.T5.collateral_change_rate\endcsname{[0.0, 0.0]}
\expandafter\def\csname nmn@sint-memory-mcp.T5.collateral_change_rate\endcsname{3}
\expandafter\def\csname nm@memory.T6.probe_survival_rate\endcsname{1.0}
\expandafter\def\csname nmc@memory.T6.probe_survival_rate\endcsname{[1.0, 1.0]}
\expandafter\def\csname nmn@memory.T6.probe_survival_rate\endcsname{3}
\expandafter\def\csname nm@sint-memory-mcp.T6.probe_survival_rate\endcsname{1.0}
\expandafter\def\csname nmc@sint-memory-mcp.T6.probe_survival_rate\endcsname{[1.0, 1.0]}
\expandafter\def\csname nmn@sint-memory-mcp.T6.probe_survival_rate\endcsname{3}
\expandafter\def\csname nm@memory.T6.compression_ratio\endcsname{1.0}
\expandafter\def\csname nmc@memory.T6.compression_ratio\endcsname{[1.0, 1.0]}
\expandafter\def\csname nmn@memory.T6.compression_ratio\endcsname{3}
\expandafter\def\csname nm@sint-memory-mcp.T6.compression_ratio\endcsname{1.0}
\expandafter\def\csname nmc@sint-memory-mcp.T6.compression_ratio\endcsname{[1.0, 1.0]}
\expandafter\def\csname nmn@sint-memory-mcp.T6.compression_ratio\endcsname{3}
\expandafter\def\csname nm@memory.T6.derived_from_completeness\endcsname{null}
\expandafter\def\csname nmc@memory.T6.derived_from_completeness\endcsname{not defined}
\expandafter\def\csname nmn@memory.T6.derived_from_completeness\endcsname{0}
\expandafter\def\csname nm@sint-memory-mcp.T6.derived_from_completeness\endcsname{null}
\expandafter\def\csname nmc@sint-memory-mcp.T6.derived_from_completeness\endcsname{not defined}
\expandafter\def\csname nmn@sint-memory-mcp.T6.derived_from_completeness\endcsname{0}
\expandafter\def\csname nm@memory.T6.abstraction_faithfulness\endcsname{null}
\expandafter\def\csname nmc@memory.T6.abstraction_faithfulness\endcsname{not defined}
\expandafter\def\csname nmn@memory.T6.abstraction_faithfulness\endcsname{0}
\expandafter\def\csname nm@sint-memory-mcp.T6.abstraction_faithfulness\endcsname{null}
\expandafter\def\csname nmc@sint-memory-mcp.T6.abstraction_faithfulness\endcsname{not defined}
\expandafter\def\csname nmn@sint-memory-mcp.T6.abstraction_faithfulness\endcsname{0}
\expandafter\def\csname nm@memory.T6.cost_per_consolidation_usd\endcsname{null}
\expandafter\def\csname nmc@memory.T6.cost_per_consolidation_usd\endcsname{not defined}
\expandafter\def\csname nmn@memory.T6.cost_per_consolidation_usd\endcsname{0}
\expandafter\def\csname nm@sint-memory-mcp.T6.cost_per_consolidation_usd\endcsname{null}
\expandafter\def\csname nmc@sint-memory-mcp.T6.cost_per_consolidation_usd\endcsname{not defined}
\expandafter\def\csname nmn@sint-memory-mcp.T6.cost_per_consolidation_usd\endcsname{0}
\expandafter\def\csname nm@memory.T7.cue_hit_rate\endcsname{1.0}
\expandafter\def\csname nmc@memory.T7.cue_hit_rate\endcsname{[1.0, 1.0]}
\expandafter\def\csname nmn@memory.T7.cue_hit_rate\endcsname{3}
\expandafter\def\csname nm@sint-memory-mcp.T7.cue_hit_rate\endcsname{1.0}
\expandafter\def\csname nmc@sint-memory-mcp.T7.cue_hit_rate\endcsname{[1.0, 1.0]}
\expandafter\def\csname nmn@sint-memory-mcp.T7.cue_hit_rate\endcsname{3}
\expandafter\def\csname nm@memory.T7.false_fire_rate\endcsname{0.0}
\expandafter\def\csname nmc@memory.T7.false_fire_rate\endcsname{[0.0, 0.0]}
\expandafter\def\csname nmn@memory.T7.false_fire_rate\endcsname{3}
\expandafter\def\csname nm@sint-memory-mcp.T7.false_fire_rate\endcsname{0.0}
\expandafter\def\csname nmc@sint-memory-mcp.T7.false_fire_rate\endcsname{[0.0, 0.0]}
\expandafter\def\csname nmn@sint-memory-mcp.T7.false_fire_rate\endcsname{3}
\expandafter\def\csname nm@memory.T7.context_availability_at_fire\endcsname{1.0}
\expandafter\def\csname nmc@memory.T7.context_availability_at_fire\endcsname{[1.0, 1.0]}
\expandafter\def\csname nmn@memory.T7.context_availability_at_fire\endcsname{3}
\expandafter\def\csname nm@sint-memory-mcp.T7.context_availability_at_fire\endcsname{1.0}
\expandafter\def\csname nmc@sint-memory-mcp.T7.context_availability_at_fire\endcsname{[1.0, 1.0]}
\expandafter\def\csname nmn@sint-memory-mcp.T7.context_availability_at_fire\endcsname{3}
\expandafter\def\csname nm@memory.T7.fire_latency_calls\endcsname{0.0}
\expandafter\def\csname nmc@memory.T7.fire_latency_calls\endcsname{[0.0, 0.0]}
\expandafter\def\csname nmn@memory.T7.fire_latency_calls\endcsname{3}
\expandafter\def\csname nm@sint-memory-mcp.T7.fire_latency_calls\endcsname{0.0}
\expandafter\def\csname nmc@sint-memory-mcp.T7.fire_latency_calls\endcsname{[0.0, 0.0]}
\expandafter\def\csname nmn@sint-memory-mcp.T7.fire_latency_calls\endcsname{3}
\expandafter\def\csname nm@memory.T8.abstention_correctness\endcsname{1.0}
\expandafter\def\csname nmc@memory.T8.abstention_correctness\endcsname{[1.0, 1.0]}
\expandafter\def\csname nmn@memory.T8.abstention_correctness\endcsname{3}
\expandafter\def\csname nm@sint-memory-mcp.T8.abstention_correctness\endcsname{1.0}
\expandafter\def\csname nmc@sint-memory-mcp.T8.abstention_correctness\endcsname{[1.0, 1.0]}
\expandafter\def\csname nmn@sint-memory-mcp.T8.abstention_correctness\endcsname{3}
\expandafter\def\csname nm@memory.T8.in_coverage_answer_correctness\endcsname{0.75}
\expandafter\def\csname nmc@memory.T8.in_coverage_answer_correctness\endcsname{[0.75, 0.75]}
\expandafter\def\csname nmn@memory.T8.in_coverage_answer_correctness\endcsname{3}
\expandafter\def\csname nm@sint-memory-mcp.T8.in_coverage_answer_correctness\endcsname{0.75}
\expandafter\def\csname nmc@sint-memory-mcp.T8.in_coverage_answer_correctness\endcsname{[0.75, 0.75]}
\expandafter\def\csname nmn@sint-memory-mcp.T8.in_coverage_answer_correctness\endcsname{3}
\expandafter\def\csname nm@memory.T8.key_question_correctness\endcsname{1.0}
\expandafter\def\csname nmc@memory.T8.key_question_correctness\endcsname{[1.0, 1.0]}
\expandafter\def\csname nmn@memory.T8.key_question_correctness\endcsname{3}
\expandafter\def\csname nm@sint-memory-mcp.T8.key_question_correctness\endcsname{1.0}
\expandafter\def\csname nmc@sint-memory-mcp.T8.key_question_correctness\endcsname{[1.0, 1.0]}
\expandafter\def\csname nmn@sint-memory-mcp.T8.key_question_correctness\endcsname{3}
\expandafter\def\csname nm@memory.T8.subject_question_correctness\endcsname{0.5}
\expandafter\def\csname nmc@memory.T8.subject_question_correctness\endcsname{[0.5, 0.5]}
\expandafter\def\csname nmn@memory.T8.subject_question_correctness\endcsname{3}
\expandafter\def\csname nm@sint-memory-mcp.T8.subject_question_correctness\endcsname{0.5}
\expandafter\def\csname nmc@sint-memory-mcp.T8.subject_question_correctness\endcsname{[0.5, 0.5]}
\expandafter\def\csname nmn@sint-memory-mcp.T8.subject_question_correctness\endcsname{3}
\expandafter\def\csname nm@memory.T8.ece\endcsname{0.4333}
\expandafter\def\csname nmc@memory.T8.ece\endcsname{[0.4333, 0.4333]}
\expandafter\def\csname nmn@memory.T8.ece\endcsname{3}
\expandafter\def\csname nm@sint-memory-mcp.T8.ece\endcsname{0.4333}
\expandafter\def\csname nmc@sint-memory-mcp.T8.ece\endcsname{[0.4333, 0.4333]}
\expandafter\def\csname nmn@sint-memory-mcp.T8.ece\endcsname{3}
\expandafter\def\csname nm@memory.T8.discrimination_auroc\endcsname{0.3}
\expandafter\def\csname nmc@memory.T8.discrimination_auroc\endcsname{[0.3, 0.3]}
\expandafter\def\csname nmn@memory.T8.discrimination_auroc\endcsname{3}
\expandafter\def\csname nm@sint-memory-mcp.T8.discrimination_auroc\endcsname{0.3}
\expandafter\def\csname nmc@sint-memory-mcp.T8.discrimination_auroc\endcsname{[0.3, 0.3]}
\expandafter\def\csname nmn@sint-memory-mcp.T8.discrimination_auroc\endcsname{3}
\expandafter\def\csname nm@memory.T8.monitoring_accuracy\endcsname{1.0}
\expandafter\def\csname nmc@memory.T8.monitoring_accuracy\endcsname{[1.0, 1.0]}
\expandafter\def\csname nmn@memory.T8.monitoring_accuracy\endcsname{3}
\expandafter\def\csname nm@sint-memory-mcp.T8.monitoring_accuracy\endcsname{1.0}
\expandafter\def\csname nmc@sint-memory-mcp.T8.monitoring_accuracy\endcsname{[1.0, 1.0]}
\expandafter\def\csname nmn@sint-memory-mcp.T8.monitoring_accuracy\endcsname{3}
\expandafter\def\csname nm@memory.T9.restart_survival\endcsname{1.0}
\expandafter\def\csname nmc@memory.T9.restart_survival\endcsname{[1.0, 1.0]}
\expandafter\def\csname nmn@memory.T9.restart_survival\endcsname{3}
\expandafter\def\csname nm@sint-memory-mcp.T9.restart_survival\endcsname{1.0}
\expandafter\def\csname nmc@sint-memory-mcp.T9.restart_survival\endcsname{[1.0, 1.0]}
\expandafter\def\csname nmn@sint-memory-mcp.T9.restart_survival\endcsname{3}
\expandafter\def\csname nm@memory.T9.snapshot_reopen_survival\endcsname{1.0}
\expandafter\def\csname nmc@memory.T9.snapshot_reopen_survival\endcsname{[1.0, 1.0]}
\expandafter\def\csname nmn@memory.T9.snapshot_reopen_survival\endcsname{3}
\expandafter\def\csname nm@sint-memory-mcp.T9.snapshot_reopen_survival\endcsname{1.0}
\expandafter\def\csname nmc@sint-memory-mcp.T9.snapshot_reopen_survival\endcsname{[1.0, 1.0]}
\expandafter\def\csname nmn@sint-memory-mcp.T9.snapshot_reopen_survival\endcsname{3}
\expandafter\def\csname nm@memory.T9.repeat_query_equality\endcsname{1.0}
\expandafter\def\csname nmc@memory.T9.repeat_query_equality\endcsname{[1.0, 1.0]}
\expandafter\def\csname nmn@memory.T9.repeat_query_equality\endcsname{3}
\expandafter\def\csname nm@sint-memory-mcp.T9.repeat_query_equality\endcsname{1.0}
\expandafter\def\csname nmc@sint-memory-mcp.T9.repeat_query_equality\endcsname{[1.0, 1.0]}
\expandafter\def\csname nmn@sint-memory-mcp.T9.repeat_query_equality\endcsname{3}
\expandafter\def\csname nm@memory.T10.poisoning_influence_rate\endcsname{0.0}
\expandafter\def\csname nmc@memory.T10.poisoning_influence_rate\endcsname{[0.0, 0.0]}
\expandafter\def\csname nmn@memory.T10.poisoning_influence_rate\endcsname{3}
\expandafter\def\csname nm@sint-memory-mcp.T10.poisoning_influence_rate\endcsname{0.0}
\expandafter\def\csname nmc@sint-memory-mcp.T10.poisoning_influence_rate\endcsname{[0.0, 0.0]}
\expandafter\def\csname nmn@sint-memory-mcp.T10.poisoning_influence_rate\endcsname{3}
\expandafter\def\csname nm@memory.T10.membership_tpr\endcsname{1.0}
\expandafter\def\csname nmc@memory.T10.membership_tpr\endcsname{[1.0, 1.0]}
\expandafter\def\csname nmn@memory.T10.membership_tpr\endcsname{3}
\expandafter\def\csname nm@sint-memory-mcp.T10.membership_tpr\endcsname{1.0}
\expandafter\def\csname nmc@sint-memory-mcp.T10.membership_tpr\endcsname{[1.0, 1.0]}
\expandafter\def\csname nmn@sint-memory-mcp.T10.membership_tpr\endcsname{3}
\expandafter\def\csname nm@memory.T10.membership_fpr\endcsname{0.0}
\expandafter\def\csname nmc@memory.T10.membership_fpr\endcsname{[0.0, 0.0]}
\expandafter\def\csname nmn@memory.T10.membership_fpr\endcsname{3}
\expandafter\def\csname nm@sint-memory-mcp.T10.membership_fpr\endcsname{0.0}
\expandafter\def\csname nmc@sint-memory-mcp.T10.membership_fpr\endcsname{[0.0, 0.0]}
\expandafter\def\csname nmn@sint-memory-mcp.T10.membership_fpr\endcsname{3}
\expandafter\def\csname nm@memory.T10.membership_inference_advantage\endcsname{1.0}
\expandafter\def\csname nmc@memory.T10.membership_inference_advantage\endcsname{[1.0, 1.0]}
\expandafter\def\csname nmn@memory.T10.membership_inference_advantage\endcsname{3}
\expandafter\def\csname nm@sint-memory-mcp.T10.membership_inference_advantage\endcsname{1.0}
\expandafter\def\csname nmc@sint-memory-mcp.T10.membership_inference_advantage\endcsname{[1.0, 1.0]}
\expandafter\def\csname nmn@sint-memory-mcp.T10.membership_inference_advantage\endcsname{3}
\expandafter\def\csname nm@memory.O1_WRITE.write_precision\endcsname{0.75}
\expandafter\def\csname nmc@memory.O1_WRITE.write_precision\endcsname{[0.75, 0.75]}
\expandafter\def\csname nmn@memory.O1_WRITE.write_precision\endcsname{3}
\expandafter\def\csname nm@sint-memory-mcp.O1_WRITE.write_precision\endcsname{0.75}
\expandafter\def\csname nmc@sint-memory-mcp.O1_WRITE.write_precision\endcsname{[0.75, 0.75]}
\expandafter\def\csname nmn@sint-memory-mcp.O1_WRITE.write_precision\endcsname{3}
\expandafter\def\csname nm@memory.O1_WRITE.write_recall\endcsname{1.0}
\expandafter\def\csname nmc@memory.O1_WRITE.write_recall\endcsname{[1.0, 1.0]}
\expandafter\def\csname nmn@memory.O1_WRITE.write_recall\endcsname{3}
\expandafter\def\csname nm@sint-memory-mcp.O1_WRITE.write_recall\endcsname{1.0}
\expandafter\def\csname nmc@sint-memory-mcp.O1_WRITE.write_recall\endcsname{[1.0, 1.0]}
\expandafter\def\csname nmn@sint-memory-mcp.O1_WRITE.write_recall\endcsname{3}
\expandafter\def\csname nm@memory.O1_WRITE.duplicate_rate\endcsname{0.25}
\expandafter\def\csname nmc@memory.O1_WRITE.duplicate_rate\endcsname{[0.25, 0.25]}
\expandafter\def\csname nmn@memory.O1_WRITE.duplicate_rate\endcsname{3}
\expandafter\def\csname nm@sint-memory-mcp.O1_WRITE.duplicate_rate\endcsname{0.25}
\expandafter\def\csname nmc@sint-memory-mcp.O1_WRITE.duplicate_rate\endcsname{[0.25, 0.25]}
\expandafter\def\csname nmn@sint-memory-mcp.O1_WRITE.duplicate_rate\endcsname{3}
\expandafter\def\csname nm@memory.O1_WRITE.write_latency_p50_ms\endcsname{16.4423}
\expandafter\def\csname nmc@memory.O1_WRITE.write_latency_p50_ms\endcsname{[15.167, 17.2955]}
\expandafter\def\csname nmn@memory.O1_WRITE.write_latency_p50_ms\endcsname{3}
\expandafter\def\csname nm@sint-memory-mcp.O1_WRITE.write_latency_p50_ms\endcsname{16.4423}
\expandafter\def\csname nmc@sint-memory-mcp.O1_WRITE.write_latency_p50_ms\endcsname{[15.167, 17.2955]}
\expandafter\def\csname nmn@sint-memory-mcp.O1_WRITE.write_latency_p50_ms\endcsname{3}
\expandafter\def\csname nm@memory.O1_WRITE.write_latency_p95_ms\endcsname{18.9019}
\expandafter\def\csname nmc@memory.O1_WRITE.write_latency_p95_ms\endcsname{[17.329, 20.2025]}
\expandafter\def\csname nmn@memory.O1_WRITE.write_latency_p95_ms\endcsname{3}
\expandafter\def\csname nm@sint-memory-mcp.O1_WRITE.write_latency_p95_ms\endcsname{18.9019}
\expandafter\def\csname nmc@sint-memory-mcp.O1_WRITE.write_latency_p95_ms\endcsname{[17.329, 20.2025]}
\expandafter\def\csname nmn@sint-memory-mcp.O1_WRITE.write_latency_p95_ms\endcsname{3}
\expandafter\def\csname nm@memory.O1_WRITE.tokens_per_write\endcsname{null}
\expandafter\def\csname nmc@memory.O1_WRITE.tokens_per_write\endcsname{not defined}
\expandafter\def\csname nmn@memory.O1_WRITE.tokens_per_write\endcsname{0}
\expandafter\def\csname nm@sint-memory-mcp.O1_WRITE.tokens_per_write\endcsname{null}
\expandafter\def\csname nmc@sint-memory-mcp.O1_WRITE.tokens_per_write\endcsname{not defined}
\expandafter\def\csname nmn@sint-memory-mcp.O1_WRITE.tokens_per_write\endcsname{0}
\expandafter\def\csname nm@memory.O1_WRITE.cost_per_1000_writes_usd\endcsname{null}
\expandafter\def\csname nmc@memory.O1_WRITE.cost_per_1000_writes_usd\endcsname{not defined}
\expandafter\def\csname nmn@memory.O1_WRITE.cost_per_1000_writes_usd\endcsname{0}
\expandafter\def\csname nm@sint-memory-mcp.O1_WRITE.cost_per_1000_writes_usd\endcsname{null}
\expandafter\def\csname nmc@sint-memory-mcp.O1_WRITE.cost_per_1000_writes_usd\endcsname{not defined}
\expandafter\def\csname nmn@sint-memory-mcp.O1_WRITE.cost_per_1000_writes_usd\endcsname{0}
\expandafter\def\csname nm@memory.O1_WRITE.idempotency_violations\endcsname{4.0}
\expandafter\def\csname nmc@memory.O1_WRITE.idempotency_violations\endcsname{[4.0, 4.0]}
\expandafter\def\csname nmn@memory.O1_WRITE.idempotency_violations\endcsname{3}
\expandafter\def\csname nm@sint-memory-mcp.O1_WRITE.idempotency_violations\endcsname{4.0}
\expandafter\def\csname nmc@sint-memory-mcp.O1_WRITE.idempotency_violations\endcsname{[4.0, 4.0]}
\expandafter\def\csname nmn@sint-memory-mcp.O1_WRITE.idempotency_violations\endcsname{3}
\expandafter\def\csname nm@memory.O2_READ.recall_at_5\endcsname{1.0}
\expandafter\def\csname nmc@memory.O2_READ.recall_at_5\endcsname{[1.0, 1.0]}
\expandafter\def\csname nmn@memory.O2_READ.recall_at_5\endcsname{3}
\expandafter\def\csname nm@sint-memory-mcp.O2_READ.recall_at_5\endcsname{1.0}
\expandafter\def\csname nmc@sint-memory-mcp.O2_READ.recall_at_5\endcsname{[1.0, 1.0]}
\expandafter\def\csname nmn@sint-memory-mcp.O2_READ.recall_at_5\endcsname{3}
\expandafter\def\csname nm@memory.O2_READ.mrr\endcsname{0.75}
\expandafter\def\csname nmc@memory.O2_READ.mrr\endcsname{[0.6806, 0.8056]}
\expandafter\def\csname nmn@memory.O2_READ.mrr\endcsname{3}
\expandafter\def\csname nm@sint-memory-mcp.O2_READ.mrr\endcsname{0.75}
\expandafter\def\csname nmc@sint-memory-mcp.O2_READ.mrr\endcsname{[0.6806, 0.8056]}
\expandafter\def\csname nmn@sint-memory-mcp.O2_READ.mrr\endcsname{3}
\expandafter\def\csname nm@memory.O2_READ.required_record_hit_rate\endcsname{1.0}
\expandafter\def\csname nmc@memory.O2_READ.required_record_hit_rate\endcsname{[1.0, 1.0]}
\expandafter\def\csname nmn@memory.O2_READ.required_record_hit_rate\endcsname{3}
\expandafter\def\csname nm@sint-memory-mcp.O2_READ.required_record_hit_rate\endcsname{1.0}
\expandafter\def\csname nmc@sint-memory-mcp.O2_READ.required_record_hit_rate\endcsname{[1.0, 1.0]}
\expandafter\def\csname nmn@sint-memory-mcp.O2_READ.required_record_hit_rate\endcsname{3}
\expandafter\def\csname nm@memory.O2_READ.provenance_correctness\endcsname{1.0}
\expandafter\def\csname nmc@memory.O2_READ.provenance_correctness\endcsname{[1.0, 1.0]}
\expandafter\def\csname nmn@memory.O2_READ.provenance_correctness\endcsname{3}
\expandafter\def\csname nm@sint-memory-mcp.O2_READ.provenance_correctness\endcsname{1.0}
\expandafter\def\csname nmc@sint-memory-mcp.O2_READ.provenance_correctness\endcsname{[1.0, 1.0]}
\expandafter\def\csname nmn@sint-memory-mcp.O2_READ.provenance_correctness\endcsname{3}
\expandafter\def\csname nm@memory.O2_READ.admissibility_error_rate\endcsname{0.5}
\expandafter\def\csname nmc@memory.O2_READ.admissibility_error_rate\endcsname{[0.5, 0.5]}
\expandafter\def\csname nmn@memory.O2_READ.admissibility_error_rate\endcsname{3}
\expandafter\def\csname nm@sint-memory-mcp.O2_READ.admissibility_error_rate\endcsname{0.5}
\expandafter\def\csname nmc@sint-memory-mcp.O2_READ.admissibility_error_rate\endcsname{[0.5, 0.5]}
\expandafter\def\csname nmn@sint-memory-mcp.O2_READ.admissibility_error_rate\endcsname{3}
\expandafter\def\csname nm@memory.O2_READ.context_inflation\endcsname{2.5}
\expandafter\def\csname nmc@memory.O2_READ.context_inflation\endcsname{[2.5, 2.5]}
\expandafter\def\csname nmn@memory.O2_READ.context_inflation\endcsname{3}
\expandafter\def\csname nm@sint-memory-mcp.O2_READ.context_inflation\endcsname{2.5}
\expandafter\def\csname nmc@sint-memory-mcp.O2_READ.context_inflation\endcsname{[2.5, 2.5]}
\expandafter\def\csname nmn@sint-memory-mcp.O2_READ.context_inflation\endcsname{3}
\expandafter\def\csname nm@memory.O2_READ.read_latency_p50_ms\endcsname{8.0937}
\expandafter\def\csname nmc@memory.O2_READ.read_latency_p50_ms\endcsname{[7.6565, 8.5435]}
\expandafter\def\csname nmn@memory.O2_READ.read_latency_p50_ms\endcsname{3}
\expandafter\def\csname nm@sint-memory-mcp.O2_READ.read_latency_p50_ms\endcsname{8.0937}
\expandafter\def\csname nmc@sint-memory-mcp.O2_READ.read_latency_p50_ms\endcsname{[7.6565, 8.5435]}
\expandafter\def\csname nmn@sint-memory-mcp.O2_READ.read_latency_p50_ms\endcsname{3}
\expandafter\def\csname nm@memory.O2_READ.read_latency_p95_ms\endcsname{102.8932}
\expandafter\def\csname nmc@memory.O2_READ.read_latency_p95_ms\endcsname{[101.5778, 104.9318]}
\expandafter\def\csname nmn@memory.O2_READ.read_latency_p95_ms\endcsname{3}
\expandafter\def\csname nm@sint-memory-mcp.O2_READ.read_latency_p95_ms\endcsname{102.8932}
\expandafter\def\csname nmc@sint-memory-mcp.O2_READ.read_latency_p95_ms\endcsname{[101.5778, 104.9318]}
\expandafter\def\csname nmn@sint-memory-mcp.O2_READ.read_latency_p95_ms\endcsname{3}
\expandafter\def\csname nm@memory.O3_CONSOLIDATE.compression_ratio\endcsname{1.0}
\expandafter\def\csname nmc@memory.O3_CONSOLIDATE.compression_ratio\endcsname{[1.0, 1.0]}
\expandafter\def\csname nmn@memory.O3_CONSOLIDATE.compression_ratio\endcsname{3}
\expandafter\def\csname nm@sint-memory-mcp.O3_CONSOLIDATE.compression_ratio\endcsname{1.0}
\expandafter\def\csname nmc@sint-memory-mcp.O3_CONSOLIDATE.compression_ratio\endcsname{[1.0, 1.0]}
\expandafter\def\csname nmn@sint-memory-mcp.O3_CONSOLIDATE.compression_ratio\endcsname{3}
\expandafter\def\csname nm@memory.O3_CONSOLIDATE.probe_survival_rate\endcsname{1.0}
\expandafter\def\csname nmc@memory.O3_CONSOLIDATE.probe_survival_rate\endcsname{[1.0, 1.0]}
\expandafter\def\csname nmn@memory.O3_CONSOLIDATE.probe_survival_rate\endcsname{3}
\expandafter\def\csname nm@sint-memory-mcp.O3_CONSOLIDATE.probe_survival_rate\endcsname{1.0}
\expandafter\def\csname nmc@sint-memory-mcp.O3_CONSOLIDATE.probe_survival_rate\endcsname{[1.0, 1.0]}
\expandafter\def\csname nmn@sint-memory-mcp.O3_CONSOLIDATE.probe_survival_rate\endcsname{3}
\expandafter\def\csname nm@memory.O3_CONSOLIDATE.derived_from_completeness\endcsname{null}
\expandafter\def\csname nmc@memory.O3_CONSOLIDATE.derived_from_completeness\endcsname{not defined}
\expandafter\def\csname nmn@memory.O3_CONSOLIDATE.derived_from_completeness\endcsname{0}
\expandafter\def\csname nm@sint-memory-mcp.O3_CONSOLIDATE.derived_from_completeness\endcsname{null}
\expandafter\def\csname nmc@sint-memory-mcp.O3_CONSOLIDATE.derived_from_completeness\endcsname{not defined}
\expandafter\def\csname nmn@sint-memory-mcp.O3_CONSOLIDATE.derived_from_completeness\endcsname{0}
\expandafter\def\csname nm@memory.O3_CONSOLIDATE.cost_per_consolidation_usd\endcsname{null}
\expandafter\def\csname nmc@memory.O3_CONSOLIDATE.cost_per_consolidation_usd\endcsname{not defined}
\expandafter\def\csname nmn@memory.O3_CONSOLIDATE.cost_per_consolidation_usd\endcsname{0}
\expandafter\def\csname nm@sint-memory-mcp.O3_CONSOLIDATE.cost_per_consolidation_usd\endcsname{null}
\expandafter\def\csname nmc@sint-memory-mcp.O3_CONSOLIDATE.cost_per_consolidation_usd\endcsname{not defined}
\expandafter\def\csname nmn@sint-memory-mcp.O3_CONSOLIDATE.cost_per_consolidation_usd\endcsname{0}
\expandafter\def\csname nm@memory.O4_REVISE.contradiction_detection_f1\endcsname{0.0}
\expandafter\def\csname nmc@memory.O4_REVISE.contradiction_detection_f1\endcsname{[0.0, 0.0]}
\expandafter\def\csname nmn@memory.O4_REVISE.contradiction_detection_f1\endcsname{3}
\expandafter\def\csname nm@sint-memory-mcp.O4_REVISE.contradiction_detection_f1\endcsname{0.0}
\expandafter\def\csname nmc@sint-memory-mcp.O4_REVISE.contradiction_detection_f1\endcsname{[0.0, 0.0]}
\expandafter\def\csname nmn@sint-memory-mcp.O4_REVISE.contradiction_detection_f1\endcsname{3}
\expandafter\def\csname nm@memory.O4_REVISE.revision_correctness\endcsname{null}
\expandafter\def\csname nmc@memory.O4_REVISE.revision_correctness\endcsname{not defined}
\expandafter\def\csname nmn@memory.O4_REVISE.revision_correctness\endcsname{0}
\expandafter\def\csname nm@sint-memory-mcp.O4_REVISE.revision_correctness\endcsname{null}
\expandafter\def\csname nmc@sint-memory-mcp.O4_REVISE.revision_correctness\endcsname{not defined}
\expandafter\def\csname nmn@sint-memory-mcp.O4_REVISE.revision_correctness\endcsname{0}
\expandafter\def\csname nm@memory.O4_REVISE.chain_completeness\endcsname{0.0}
\expandafter\def\csname nmc@memory.O4_REVISE.chain_completeness\endcsname{[0.0, 0.0]}
\expandafter\def\csname nmn@memory.O4_REVISE.chain_completeness\endcsname{3}
\expandafter\def\csname nm@sint-memory-mcp.O4_REVISE.chain_completeness\endcsname{0.0}
\expandafter\def\csname nmc@sint-memory-mcp.O4_REVISE.chain_completeness\endcsname{[0.0, 0.0]}
\expandafter\def\csname nmn@sint-memory-mcp.O4_REVISE.chain_completeness\endcsname{3}
\expandafter\def\csname nm@memory.O4_REVISE.stale_read_rate\endcsname{1.0}
\expandafter\def\csname nmc@memory.O4_REVISE.stale_read_rate\endcsname{[1.0, 1.0]}
\expandafter\def\csname nmn@memory.O4_REVISE.stale_read_rate\endcsname{3}
\expandafter\def\csname nm@sint-memory-mcp.O4_REVISE.stale_read_rate\endcsname{1.0}
\expandafter\def\csname nmc@sint-memory-mcp.O4_REVISE.stale_read_rate\endcsname{[1.0, 1.0]}
\expandafter\def\csname nmn@sint-memory-mcp.O4_REVISE.stale_read_rate\endcsname{3}
\expandafter\def\csname nm@memory.O4_REVISE.collateral_change_rate\endcsname{0.0}
\expandafter\def\csname nmc@memory.O4_REVISE.collateral_change_rate\endcsname{[0.0, 0.0]}
\expandafter\def\csname nmn@memory.O4_REVISE.collateral_change_rate\endcsname{3}
\expandafter\def\csname nm@sint-memory-mcp.O4_REVISE.collateral_change_rate\endcsname{0.0}
\expandafter\def\csname nmc@sint-memory-mcp.O4_REVISE.collateral_change_rate\endcsname{[0.0, 0.0]}
\expandafter\def\csname nmn@sint-memory-mcp.O4_REVISE.collateral_change_rate\endcsname{3}
\expandafter\def\csname nm@memory.O5_INTEND.cue_hit_rate\endcsname{1.0}
\expandafter\def\csname nmc@memory.O5_INTEND.cue_hit_rate\endcsname{[1.0, 1.0]}
\expandafter\def\csname nmn@memory.O5_INTEND.cue_hit_rate\endcsname{3}
\expandafter\def\csname nm@sint-memory-mcp.O5_INTEND.cue_hit_rate\endcsname{1.0}
\expandafter\def\csname nmc@sint-memory-mcp.O5_INTEND.cue_hit_rate\endcsname{[1.0, 1.0]}
\expandafter\def\csname nmn@sint-memory-mcp.O5_INTEND.cue_hit_rate\endcsname{3}
\expandafter\def\csname nm@memory.O5_INTEND.false_fire_rate\endcsname{0.0}
\expandafter\def\csname nmc@memory.O5_INTEND.false_fire_rate\endcsname{[0.0, 0.0]}
\expandafter\def\csname nmn@memory.O5_INTEND.false_fire_rate\endcsname{3}
\expandafter\def\csname nm@sint-memory-mcp.O5_INTEND.false_fire_rate\endcsname{0.0}
\expandafter\def\csname nmc@sint-memory-mcp.O5_INTEND.false_fire_rate\endcsname{[0.0, 0.0]}
\expandafter\def\csname nmn@sint-memory-mcp.O5_INTEND.false_fire_rate\endcsname{3}
\expandafter\def\csname nm@memory.O5_INTEND.fire_latency_calls\endcsname{0.0}
\expandafter\def\csname nmc@memory.O5_INTEND.fire_latency_calls\endcsname{[0.0, 0.0]}
\expandafter\def\csname nmn@memory.O5_INTEND.fire_latency_calls\endcsname{3}
\expandafter\def\csname nm@sint-memory-mcp.O5_INTEND.fire_latency_calls\endcsname{0.0}
\expandafter\def\csname nmc@sint-memory-mcp.O5_INTEND.fire_latency_calls\endcsname{[0.0, 0.0]}
\expandafter\def\csname nmn@sint-memory-mcp.O5_INTEND.fire_latency_calls\endcsname{3}
\expandafter\def\csname nm@memory.O5_INTEND.context_availability_at_fire\endcsname{1.0}
\expandafter\def\csname nmc@memory.O5_INTEND.context_availability_at_fire\endcsname{[1.0, 1.0]}
\expandafter\def\csname nmn@memory.O5_INTEND.context_availability_at_fire\endcsname{3}
\expandafter\def\csname nm@sint-memory-mcp.O5_INTEND.context_availability_at_fire\endcsname{1.0}
\expandafter\def\csname nmc@sint-memory-mcp.O5_INTEND.context_availability_at_fire\endcsname{[1.0, 1.0]}
\expandafter\def\csname nmn@sint-memory-mcp.O5_INTEND.context_availability_at_fire\endcsname{3}
\expandafter\def\csname nm@memory.O6_ASSERT.abstention_correctness\endcsname{1.0}
\expandafter\def\csname nmc@memory.O6_ASSERT.abstention_correctness\endcsname{[1.0, 1.0]}
\expandafter\def\csname nmn@memory.O6_ASSERT.abstention_correctness\endcsname{3}
\expandafter\def\csname nm@sint-memory-mcp.O6_ASSERT.abstention_correctness\endcsname{1.0}
\expandafter\def\csname nmc@sint-memory-mcp.O6_ASSERT.abstention_correctness\endcsname{[1.0, 1.0]}
\expandafter\def\csname nmn@sint-memory-mcp.O6_ASSERT.abstention_correctness\endcsname{3}
\expandafter\def\csname nm@memory.O6_ASSERT.in_coverage_answer_correctness\endcsname{0.75}
\expandafter\def\csname nmc@memory.O6_ASSERT.in_coverage_answer_correctness\endcsname{[0.75, 0.75]}
\expandafter\def\csname nmn@memory.O6_ASSERT.in_coverage_answer_correctness\endcsname{3}
\expandafter\def\csname nm@sint-memory-mcp.O6_ASSERT.in_coverage_answer_correctness\endcsname{0.75}
\expandafter\def\csname nmc@sint-memory-mcp.O6_ASSERT.in_coverage_answer_correctness\endcsname{[0.75, 0.75]}
\expandafter\def\csname nmn@sint-memory-mcp.O6_ASSERT.in_coverage_answer_correctness\endcsname{3}
\expandafter\def\csname nm@memory.O6_ASSERT.key_question_correctness\endcsname{1.0}
\expandafter\def\csname nmc@memory.O6_ASSERT.key_question_correctness\endcsname{[1.0, 1.0]}
\expandafter\def\csname nmn@memory.O6_ASSERT.key_question_correctness\endcsname{3}
\expandafter\def\csname nm@sint-memory-mcp.O6_ASSERT.key_question_correctness\endcsname{1.0}
\expandafter\def\csname nmc@sint-memory-mcp.O6_ASSERT.key_question_correctness\endcsname{[1.0, 1.0]}
\expandafter\def\csname nmn@sint-memory-mcp.O6_ASSERT.key_question_correctness\endcsname{3}
\expandafter\def\csname nm@memory.O6_ASSERT.subject_question_correctness\endcsname{0.5}
\expandafter\def\csname nmc@memory.O6_ASSERT.subject_question_correctness\endcsname{[0.5, 0.5]}
\expandafter\def\csname nmn@memory.O6_ASSERT.subject_question_correctness\endcsname{3}
\expandafter\def\csname nm@sint-memory-mcp.O6_ASSERT.subject_question_correctness\endcsname{0.5}
\expandafter\def\csname nmc@sint-memory-mcp.O6_ASSERT.subject_question_correctness\endcsname{[0.5, 0.5]}
\expandafter\def\csname nmn@sint-memory-mcp.O6_ASSERT.subject_question_correctness\endcsname{3}
\expandafter\def\csname nm@memory.O6_ASSERT.ece\endcsname{0.4333}
\expandafter\def\csname nmc@memory.O6_ASSERT.ece\endcsname{[0.4333, 0.4333]}
\expandafter\def\csname nmn@memory.O6_ASSERT.ece\endcsname{3}
\expandafter\def\csname nm@sint-memory-mcp.O6_ASSERT.ece\endcsname{0.4333}
\expandafter\def\csname nmc@sint-memory-mcp.O6_ASSERT.ece\endcsname{[0.4333, 0.4333]}
\expandafter\def\csname nmn@sint-memory-mcp.O6_ASSERT.ece\endcsname{3}
\expandafter\def\csname nm@memory.O6_ASSERT.discrimination_auroc\endcsname{0.3}
\expandafter\def\csname nmc@memory.O6_ASSERT.discrimination_auroc\endcsname{[0.3, 0.3]}
\expandafter\def\csname nmn@memory.O6_ASSERT.discrimination_auroc\endcsname{3}
\expandafter\def\csname nm@sint-memory-mcp.O6_ASSERT.discrimination_auroc\endcsname{0.3}
\expandafter\def\csname nmc@sint-memory-mcp.O6_ASSERT.discrimination_auroc\endcsname{[0.3, 0.3]}
\expandafter\def\csname nmn@sint-memory-mcp.O6_ASSERT.discrimination_auroc\endcsname{3}
\expandafter\def\csname nm@memory.O6_ASSERT.monitoring_accuracy\endcsname{1.0}
\expandafter\def\csname nmc@memory.O6_ASSERT.monitoring_accuracy\endcsname{[1.0, 1.0]}
\expandafter\def\csname nmn@memory.O6_ASSERT.monitoring_accuracy\endcsname{3}
\expandafter\def\csname nm@sint-memory-mcp.O6_ASSERT.monitoring_accuracy\endcsname{1.0}
\expandafter\def\csname nmc@sint-memory-mcp.O6_ASSERT.monitoring_accuracy\endcsname{[1.0, 1.0]}
\expandafter\def\csname nmn@sint-memory-mcp.O6_ASSERT.monitoring_accuracy\endcsname{3}
\expandafter\def\csname nm@memory.PROVENANCE.provenance_correctness\endcsname{1.0}
\expandafter\def\csname nmc@memory.PROVENANCE.provenance_correctness\endcsname{[1.0, 1.0]}
\expandafter\def\csname nmn@memory.PROVENANCE.provenance_correctness\endcsname{3}
\expandafter\def\csname nm@sint-memory-mcp.PROVENANCE.provenance_correctness\endcsname{1.0}
\expandafter\def\csname nmc@sint-memory-mcp.PROVENANCE.provenance_correctness\endcsname{[1.0, 1.0]}
\expandafter\def\csname nmn@sint-memory-mcp.PROVENANCE.provenance_correctness\endcsname{3}
\expandafter\def\csname nm@memory.PROVENANCE.derived_from_completeness\endcsname{null}
\expandafter\def\csname nmc@memory.PROVENANCE.derived_from_completeness\endcsname{not defined}
\expandafter\def\csname nmn@memory.PROVENANCE.derived_from_completeness\endcsname{0}
\expandafter\def\csname nm@sint-memory-mcp.PROVENANCE.derived_from_completeness\endcsname{null}
\expandafter\def\csname nmc@sint-memory-mcp.PROVENANCE.derived_from_completeness\endcsname{not defined}
\expandafter\def\csname nmn@sint-memory-mcp.PROVENANCE.derived_from_completeness\endcsname{0}
\expandafter\def\csname nm@memory.LIFECYCLE.lifecycle_state_exposed\endcsname{0.0}
\expandafter\def\csname nmc@memory.LIFECYCLE.lifecycle_state_exposed\endcsname{[0.0, 0.0]}
\expandafter\def\csname nmn@memory.LIFECYCLE.lifecycle_state_exposed\endcsname{3}
\expandafter\def\csname nm@sint-memory-mcp.LIFECYCLE.lifecycle_state_exposed\endcsname{0.0}
\expandafter\def\csname nmc@sint-memory-mcp.LIFECYCLE.lifecycle_state_exposed\endcsname{[0.0, 0.0]}
\expandafter\def\csname nmn@sint-memory-mcp.LIFECYCLE.lifecycle_state_exposed\endcsname{3}
\expandafter\def\csname nm@memory.LIFECYCLE.admissibility_error_rate\endcsname{0.5}
\expandafter\def\csname nmc@memory.LIFECYCLE.admissibility_error_rate\endcsname{[0.5, 0.5]}
\expandafter\def\csname nmn@memory.LIFECYCLE.admissibility_error_rate\endcsname{3}
\expandafter\def\csname nm@sint-memory-mcp.LIFECYCLE.admissibility_error_rate\endcsname{0.5}
\expandafter\def\csname nmc@sint-memory-mcp.LIFECYCLE.admissibility_error_rate\endcsname{[0.5, 0.5]}
\expandafter\def\csname nmn@sint-memory-mcp.LIFECYCLE.admissibility_error_rate\endcsname{3}
\expandafter\def\csname nm@memory.LIFECYCLE.stale_read_rate\endcsname{1.0}
\expandafter\def\csname nmc@memory.LIFECYCLE.stale_read_rate\endcsname{[1.0, 1.0]}
\expandafter\def\csname nmn@memory.LIFECYCLE.stale_read_rate\endcsname{3}
\expandafter\def\csname nm@sint-memory-mcp.LIFECYCLE.stale_read_rate\endcsname{1.0}
\expandafter\def\csname nmc@sint-memory-mcp.LIFECYCLE.stale_read_rate\endcsname{[1.0, 1.0]}
\expandafter\def\csname nmn@sint-memory-mcp.LIFECYCLE.stale_read_rate\endcsname{3}
\expandafter\def\csname nm@memory.LIFECYCLE.restart_survival\endcsname{1.0}
\expandafter\def\csname nmc@memory.LIFECYCLE.restart_survival\endcsname{[1.0, 1.0]}
\expandafter\def\csname nmn@memory.LIFECYCLE.restart_survival\endcsname{3}
\expandafter\def\csname nm@sint-memory-mcp.LIFECYCLE.restart_survival\endcsname{1.0}
\expandafter\def\csname nmc@sint-memory-mcp.LIFECYCLE.restart_survival\endcsname{[1.0, 1.0]}
\expandafter\def\csname nmn@sint-memory-mcp.LIFECYCLE.restart_survival\endcsname{3}
\expandafter\def\csname nm@memory.LIFECYCLE.snapshot_reopen_survival\endcsname{1.0}
\expandafter\def\csname nmc@memory.LIFECYCLE.snapshot_reopen_survival\endcsname{[1.0, 1.0]}
\expandafter\def\csname nmn@memory.LIFECYCLE.snapshot_reopen_survival\endcsname{3}
\expandafter\def\csname nm@sint-memory-mcp.LIFECYCLE.snapshot_reopen_survival\endcsname{1.0}
\expandafter\def\csname nmc@sint-memory-mcp.LIFECYCLE.snapshot_reopen_survival\endcsname{[1.0, 1.0]}
\expandafter\def\csname nmn@sint-memory-mcp.LIFECYCLE.snapshot_reopen_survival\endcsname{3}
\expandafter\def\csname nm@memory.metrics_total\endcsname{52}
\expandafter\def\csname nm@memory.metrics_valued\endcsname{45}
\expandafter\def\csname nm@memory.metrics_null\endcsname{7}
\expandafter\def\csname nm@memory.repetitions\endcsname{3}
\expandafter\def\csname nmt@memory.seeds\endcsname{1, 2, 3}
\expandafter\def\csname nmt@memory.ci_method\endcsname{bootstrap\_percentile}
\expandafter\def\csname nm@memory.ci_resamples\endcsname{1000}
\expandafter\def\csname nm@memory.calls_issued\endcsname{660}
\expandafter\def\csname nm@memory.tools_advertised\endcsname{19}
\expandafter\def\csname nm@memory.raw_files\endcsname{33}
\expandafter\def\csname nm@memory.failures.not_implemented\endcsname{24}
\expandafter\def\csname nm@memory.T1.tasks\endcsname{12}
\expandafter\def\csname nm@memory.T2.tasks\endcsname{6}
\expandafter\def\csname nm@memory.T3.tasks\endcsname{4}
\expandafter\def\csname nm@memory.T4.tasks\endcsname{6}
\expandafter\def\csname nm@memory.T5.tasks\endcsname{4}
\expandafter\def\csname nm@memory.T6.tasks\endcsname{4}
\expandafter\def\csname nm@memory.T7.tasks\endcsname{3}
\expandafter\def\csname nm@memory.T8.tasks\endcsname{12}
\expandafter\def\csname nm@memory.T9.tasks\endcsname{18}
\expandafter\def\csname nm@memory.T10.tasks\endcsname{12}
\expandafter\def\csname nmt@memory.O1_WRITE\endcsname{partial}
\expandafter\def\csname nmt@memory.O2_READ\endcsname{yes}
\expandafter\def\csname nmt@memory.O3_CONSOLIDATE\endcsname{partial}
\expandafter\def\csname nmt@memory.O4_REVISE\endcsname{partial}
\expandafter\def\csname nmt@memory.O5_INTEND\endcsname{yes}
\expandafter\def\csname nmt@memory.O6_ASSERT\endcsname{yes}
\expandafter\def\csname nmt@memory.PROVENANCE\endcsname{partial}
\expandafter\def\csname nmt@memory.LIFECYCLE\endcsname{yes}
\expandafter\def\csname nmt@memory.identity.subject\endcsname{sint-memory-mcp}
\expandafter\def\csname nmtl@memory.identity.subject\endcsname{sint-memor\hspace{0pt}y-mcp}
\expandafter\def\csname nmt@memory.identity.repository\endcsname{https://github.com/kostyk348/sint}
\expandafter\def\csname nmtl@memory.identity.repository\endcsname{https:\hspace{0pt}/\hspace{0pt}/\hspace{0pt}github.com\hspace{0pt}/\hspace{0pt}kostyk348/\hspace{0pt}sint}
\expandafter\def\csname nmt@memory.identity.commit\endcsname{7554f7eef853cebafae34e6654acc9871fb01163}
\expandafter\def\csname nmtl@memory.identity.commit\endcsname{7554f7eef8\hspace{0pt}53cebafae3\hspace{0pt}4e6654acc9\hspace{0pt}871fb01163\hspace{0pt}}
\expandafter\def\csname nmt@memory.identity.published_version\endcsname{MCP server as published in the repository snapshot (no release tag)}
\expandafter\def\csname nmtl@memory.identity.published_version\endcsname{MCP server as published in the repository\hspace{0pt} snapshot (no release tag)}
\expandafter\def\csname nmt@memory.identity.transport___MCP_revision\endcsname{stdio / 2025-03-26}
\expandafter\def\csname nmtl@memory.identity.transport___MCP_revision\endcsname{stdio /\hspace{0pt} 2025-03-26\hspace{0pt}}
\expandafter\def\csname nmt@memory.identity.runtime\endcsname{Python 3.14.7, mcp 1.9.4 (mcp.server.fastmcp)}
\expandafter\def\csname nmtl@memory.identity.runtime\endcsname{Python 3.14.7,\hspace{0pt} mcp 1.9.4 (mcp.serve\hspace{0pt}r.fastmcp)\hspace{0pt}}
\expandafter\def\csname nmt@memory.identity.deployment\endcsname{local process, fresh data directory per repetition (USERPROFILE/HOME redirected)}
\expandafter\def\csname nmtl@memory.identity.deployment\endcsname{local process,\hspace{0pt} fresh data directory per repetition\hspace{0pt} (USERPROFI\hspace{0pt}LE/\hspace{0pt}HOME redirected\hspace{0pt})}
\expandafter\def\csname nmt@memory.identity.store_kept_in_the_run_dir\endcsname{reports/sint-memory-mcp-2026-10-16/data/seed<seed>}
\expandafter\def\csname nmtl@memory.identity.store_kept_in_the_run_dir\endcsname{reports/\hspace{0pt}sint-memor\hspace{0pt}y-mcp-2026\hspace{0pt}-10-16/\hspace{0pt}data/\hspace{0pt}seed<seed>\hspace{0pt}}
\expandafter\def\csname nmt@memory.identity.model_under_test\endcsname{none (store-only)}
\expandafter\def\csname nmtl@memory.identity.model_under_test\endcsname{none (store-onl\hspace{0pt}y)}
\expandafter\def\csname nmt@memory.identity.hardware\endcsname{AMD64, 4 logical CPUs (workstation class, not a benchmark host)}
\expandafter\def\csname nmtl@memory.identity.hardware\endcsname{AMD64,\hspace{0pt} 4 logical CPUs (workstati\hspace{0pt}on class,\hspace{0pt} not a benchmark host)}
\expandafter\def\csname nmt@memory.identity.runner___harness\endcsname{uncommitted working tree: sha256:9b79d0aa12ae1660 over code/run\_protocol.py + code/streams.py + code/servers.json (this workspace has no git metadata, so there is no commit to name)}
\expandafter\def\csname nmtl@memory.identity.runner___harness\endcsname{uncommitte\hspace{0pt}d working tree:\hspace{0pt} sha256:\hspace{0pt}9b79d0aa12\hspace{0pt}ae1660 over code/\hspace{0pt}run\_\hspace{0pt}protocol.p\hspace{0pt}y + code/\hspace{0pt}streams.py\hspace{0pt} + code/\hspace{0pt}servers.js\hspace{0pt}on (this workspace has no git metadata,\hspace{0pt} so there is no commit to name)}
\expandafter\def\csname nmt@memory.identity.environment\endcsname{Windows-11-10.0.26200-SP0; Intel64 Family 6 Model 140 Stepping 2, GenuineIntel; server interpreter <hermes-venv>/Scripts/python.exe; harness Python 3.14.7 (stdlib only)}
\expandafter\def\csname nmtl@memory.identity.environment\endcsname{Windows-11\hspace{0pt}-10.0.2620\hspace{0pt}0-SP0;\hspace{0pt} Intel64 Family 6 Model 140 Stepping 2,\hspace{0pt} GenuineInt\hspace{0pt}el;\hspace{0pt} server interprete\hspace{0pt}r C:\hspace{0pt}/\hspace{0pt}Users/\hspace{0pt}Lab/\hspace{0pt}AppData/\hspace{0pt}Local/\hspace{0pt}hermes/\hspace{0pt}installs/\hspace{0pt}2636d87a16\hspace{0pt}39ae64/\hspace{0pt}environmen\hspace{0pt}ts/\hspace{0pt}f32597c25d\hspace{0pt}734363b2d2\hspace{0pt}20b038061d\hspace{0pt}05/\hspace{0pt}venv/\hspace{0pt}Scripts/\hspace{0pt}python.exe\hspace{0pt};\hspace{0pt} harness Python 3.14.7 (stdlib only)}
\expandafter\def\csname nmt@memory.identity.repetitions___seeds\endcsname{3 / 1, 2, 3}
\expandafter\def\csname nmtl@memory.identity.repetitions___seeds\endcsname{3 /\hspace{0pt} 1,\hspace{0pt} 2,\hspace{0pt} 3}
\expandafter\def\csname nmt@memory.identity.budgets_per_repetition\endcsname{400 calls, 900 s wall clock}
\expandafter\def\csname nmtl@memory.identity.budgets_per_repetition\endcsname{400 calls,\hspace{0pt} 900 s wall clock}
\expandafter\def\csname nmt@memory.identity.confidence_intervals\endcsname{bootstrap\_percentile, 1000 resamples}
\expandafter\def\csname nmtl@memory.identity.confidence_intervals\endcsname{bootstrap\_\hspace{0pt}percentile\hspace{0pt},\hspace{0pt} 1000 resamples}
\expandafter\def\csname nmt@memory.identity.failures_recorded\endcsname{not\_implemented x 24}
\expandafter\def\csname nmtl@memory.identity.failures_recorded\endcsname{not\_\hspace{0pt}implemente\hspace{0pt}d x 24}
\expandafter\def\csname nmt@memory.identity.raw_responses\endcsname{reports/sint-memory-mcp-2026-10-16/raw (33 JSONL files)}
\expandafter\def\csname nmtl@memory.identity.raw_responses\endcsname{reports/\hspace{0pt}sint-memor\hspace{0pt}y-mcp-2026\hspace{0pt}-10-16/\hspace{0pt}raw (33 JSONL files)}
\expandafter\def\csname nmt@memory.identity.report.json\endcsname{reports/sint-memory-mcp-2026-10-16/report.json}
\expandafter\def\csname nmtl@memory.identity.report.json\endcsname{reports/\hspace{0pt}sint-memor\hspace{0pt}y-mcp-2026\hspace{0pt}-10-16/\hspace{0pt}report.jso\hspace{0pt}n}
\expandafter\def\csname nm@run.systems\endcsname{2}
\expandafter\def\csname nm@run.families\endcsname{10}
\expandafter\def\csname nm@run.reports\endcsname{2}
\expandafter\def\csname nm@run.repetitions\endcsname{3}
\expandafter\def\csname nmt@run.seeds\endcsname{1, 2, 3}
\expandafter\def\csname nmt@run.ci_method\endcsname{bootstrap\_percentile}
\expandafter\def\csname nm@run.family_tasks\endcsname{81}
\makeatother
\providecommand{\val}[1]{\ifcsname nm@#1\endcsname\csname nm@#1\endcsname\else\textbf{??#1??}\fi}
\providecommand{\valci}[1]{\ifcsname nmc@#1\endcsname\csname nmc@#1\endcsname\else\textbf{??ci:#1??}\fi}
\providecommand{\valn}[1]{\ifcsname nmn@#1\endcsname\csname nmn@#1\endcsname\else\textbf{??n:#1??}\fi}
\providecommand{\valt}[1]{\ifcsname nmt@#1\endcsname\csname nmt@#1\endcsname\else\textbf{??t:#1??}\fi}
\providecommand{\valtl}[1]{\ifcsname nmtl@#1\endcsname\csname nmtl@#1\endcsname\else\textbf{??tl:#1??}\fi}


\title{\bf Memory Operations in MCP Agent Ecosystems:\\ an Operation Taxonomy, an Evaluation Protocol,\\ and a Reference Implementation}
\author{Konstantin Parashchuk\thanks{Independent researcher, Saint Petersburg, Russia.
Author contact and citation metadata in \texttt{CITATION.cff}. AI assistance in preparing this
manuscript (literature resolution, artefact collection, drafting) is disclosed in the
Acknowledgements.}}
\date{Version 1.0, 16 October 2026}

\begin{document}
\maketitle

\begin{abstract}
Agent memory has become a persistent store with its own lifecycle, while its evaluation is
still dominated by end-to-end task success measured over a monolithic black box
\cite{zhou2026agentnative,huan2026secondhalf}. Model Context Protocol (MCP) servers now expose
memory to agents through a standard interface \cite{mcpintro,mcpspec2025}, and MCP servers are
evaluated at task level \cite{luo2025mcpbench,luo2025mcpuniverse}, yet no shared,
operation-level protocol lets a reader run two memory systems and compare the resulting
report. This paper contributes four artefacts. (i) A taxonomy of six memory operations ---
WRITE, READ, CONSOLIDATE, REVISE, INTEND, ASSERT --- defined by preconditions, postconditions
and externally checkable evidence, together with two mandatory record axes (PROVENANCE,
LIFECYCLE). (ii) An evaluation protocol with ten task families, metric definitions, a run
procedure and a machine-validated report format. (iii) A reference implementation of five of
the six operations as public MCP servers, with its gaps stated before any measurement. (iv) Two
validated protocol reports of two memory servers, produced by that protocol in a store-only
configuration: \val{run.families} task families, \val{run.repetitions} repetitions per family at
seeds \valt{run.seeds}, \val{liquid.metrics_total} defined metrics per report with percentile
bootstrap confidence intervals, and per-operation profiles that are read side by side rather
than ranked, in tables generated from the report files rather than typed. The pair separates the
two servers on the write path (write precision \val{liquid.T1.write_precision} against
\val{memory.T1.write_precision}, with \val{memory.T1.idempotency_violations} idempotency
violations in one and \val{liquid.T1.idempotency_violations} in the other) and reproduces the
same readings where the mechanism is the same (restart survival, provenance correctness and
abstention on out-of-coverage questions). It surfaces one operation-level difference that
retrieval scores hide: both servers still serve superseded content after a revision call, one of
them while reporting the revision itself as correct. What is \emph{not} measured is stated with
the results: token and cost metrics are null by construction because no model is in the loop, the
sample is two servers from one author's reference set, and the confidence intervals rest on three
repetitions. All 135 sources curated for this work --- 61 of them cited in this paper, the
remainder in the companion taxonomy, protocol and related-work documents --- were resolved by
machine queries against the arXiv API, Crossref and direct HTTP probes; the resolution table, a
report validator with negative tests, and the generator of the comparison tables ship with the
paper.
\end{abstract}

\noindent\textbf{Keywords:} agent memory, Model Context Protocol, belief revision,
prospective memory, confidence calibration, evaluation protocol, reproducibility.

\section{Introduction}
\label{sec:intro}

Large language model agents are increasingly built by reorganising the runtime around the
model rather than by changing its weights: state moves into memory stores, expertise into
reusable skills, and interaction into protocols \cite{zhou2026externalization}. The consequence
for memory is a change of kind, not of degree. Memory is no longer a retrieval aid bolted onto
a context window; it is a data management system that supports persistent storage, retrieval,
update, consolidation and lifecycle governance across an agent's execution
\cite{zhou2026agentnative}. Two recent surveys describe the field along memory substrate,
cognitive mechanism and subject \cite{huan2026secondhalf} and along memory sources, forms and
operations \cite{zhang2024memorysurvey}.

Evaluation, however, has not followed. The dominant practice is to score memory through
end-to-end task success while treating the memory system as a monolithic black box
\cite{zhou2026agentnative}, and benchmark papers frequently omit the information a reader needs
to reproduce a comparison at all: scaffold, prompts, budgets \cite{moghadasi2026disclose}. The
gap matters most where memory servers are pluggable. MCP has become the dominant way to give
agents external capability \cite{anthropicmcp2024,mcpintro}, its exact revision is part of every
run's configuration \cite{mcpspec2025,mcpspec2025mar}, and MCP servers are already benchmarked
--- for tool use \cite{luo2025mcpbench,luo2025mcpuniverse,fan2025mcptoolbench}, for safety
\cite{zong2025mcpsafety} --- with documented sampling biases that make reported comparisons hard
to interpret \cite{afsar2026registry}. Memory servers enter that ecosystem with no shared
evaluation vocabulary.

\paragraph{Contributions.}
\begin{enumerate}[leftmargin=1.4em]
\item \textbf{C1 (taxonomy, Section~\ref{sec:taxonomy}).} Six memory operations defined by
pre/postconditions and observable evidence, plus two record axes (PROVENANCE, LIFECYCLE),
mapped onto seven existing systems. Distinguishing feature: operations are interface-level, so
a system that lacks one is visible as a gap rather than as a slightly lower accuracy score.
\item \textbf{C2 (protocol, Section~\ref{sec:protocol}).} Ten task families, a metric table in
which each metric names its computer (harness, programmatic check, or two-grader rubric), a run
procedure with budgets and repetitions, and a report format with a mandatory disclosure block
and a JSON-Schema validator.
\item \textbf{C3 (reference implementation, Section~\ref{sec:impl}).} Five of the six
operations implemented as public MCP servers, with commit-pinned facts and an explicit gap map
listed before any measurement.
\item \textbf{C4 (measured profiles, Section~\ref{sec:compare}).} Two protocol reports of two
memory servers that pass the report validator, with per-family metrics, bootstrap confidence
intervals, raw MCP stdio logs, and per-operation profiles written as a pair rather than as a
ranking. Every table and every inline number in that section is generated from the two report
files by one utility, and the check that regenerates them fails if they drift.
\item \textbf{C5 (integrity artefacts).} A machine-generated reference verification table
(135 sources, 0 unresolved), a report validator with negative tests, and the table generator,
all shipped in the repository.
\end{enumerate}

\paragraph{What this paper does not claim.}
It does not rank the two servers it measures: Section~\ref{sec:compare} reports per-operation
profiles with their confidence intervals, and the two subjects differ in storage design by
construction. It does not claim that its reference implementation is better than the systems it
cites: one of its servers now has a measured profile, the others do not. It does not claim that
human-like forgetting improves agent memory --- the opposite is reported for language models
\cite{wang2026simulating}. It does not claim that the taxonomy is complete, nor that the
measured pair generalises beyond the two servers it contains. The full ledger of claims, with
what would refute each, is in \texttt{claims.md}.

\section{Background}
\label{sec:background}

\paragraph{Memory as a data management problem.} The four-module decomposition proposed in
\cite{zhou2026agentnative} --- representation and storage, extraction, retrieval and routing,
maintenance --- is the most useful recent framing, and its authors evaluate twelve systems
across five workloads, finding that no single architecture dominates and that localised
maintenance is more cost-efficient than global reorganisation. Our protocol accepts that
finding as a hypothesis to be tested per operation: it is designed to produce a per-operation
profile rather than a ranking.

\paragraph{Vocabulary borrowed from cognitive science, deliberately narrowly.} We borrow the
separation of declarative and procedural memory with activation-based retrieval
\cite{laird2012soar,anderson2004actr}, the observation that store boundaries differ between
architectures \cite{laird2022actrsoar}, prospective memory as a faculty distinct from
retrospective recall \cite{einstein1990prospective}, source monitoring as a separable process
\cite{johnson1993source}, encoding specificity \cite{tulving1973encoding}, and offline
consolidation \cite{walker2004sleep}. The analogy is used to name operations and to choose what
to measure. It is not used to claim that artificial memory should be built to decay like human
memory: what evidence exists points the other way \cite{wang2026simulating}, so decay appears
in the taxonomy as a policy variant to be measured.

\paragraph{Belief change.} Revision operations are named after the AGM tradition
\cite{falakh2021agm} with its truth-maintenance lineage \cite{falkenhainer2013belief}; we do
not claim that memory stores satisfy the postulates, only that a revision policy must be stated
in advance if revision correctness is to be scored. Conflict classes are taken from
\cite{xu2024conflictsurvey,su2024conflictbank}, and answer-level evaluation is known to hide
evidence use \cite{chen2026ragwrong,schuster2026sourcepref}.

\paragraph{MCP.} The protocol targets MCP tool calls: server and client discovery, tools,
resources, prompts \cite{mcpspec2025}, framed by JSON-RPC 2.0 \cite{jsonrpc2}, with metadata
integrity already a demonstrated concern \cite{rashidi2026unicode}.

\section{Taxonomy of memory operations}
\label{sec:taxonomy}

An operation is defined by a precondition, a postcondition, observable evidence and a set of
failure modes; it is not defined by the data structure that implements it. This choice is what
makes comparison possible: an embedding index, a temporal knowledge graph
\cite{rasmussen2025zep,graphiti} and a table of rows can all satisfy the same postcondition, and
a system that omits an operation cannot hide behind end-to-end accuracy. The full definitions,
including variants and failure modes, are in \texttt{taxonomy.md}; Table~\ref{tab:ops} is the
summary.

\begin{table}[t]
\centering
\small
\begin{tabular}{@{}p{2.1cm}p{6.4cm}p{5.6cm}@{}}
\toprule
Operation & Postcondition (what must be true afterwards) & Evidence the interface must return \\
\midrule
O1 WRITE & exactly one record created or mutated, carrying source, author, timestamp,
derivation and lifecycle state; idempotent under a supplied key & record id plus a read-back of
the stored record \\
O2 READ & returned set is relevant, \emph{admissible} for the caller, and traceable & ordered
(record id, score, provenance) list; dropped candidates reported separately from ``no candidate
existed'' \\
O3 CONSOLIDATE & input batch replaced by fewer generalisations, each listing the ids it derives
from & compression ratio, derivation links, probe set re-run after consolidation \\
O4 REVISE & conflict resolved explicitly (supersede / reject / unresolved), revision chain
readable, unrelated records untouched & revision chain, hash of unchanged records, post-revision
read \\
O5 INTEND & deferred intention survives context loss, fires exactly once at its cue with its
context available & intention id, trigger spec, fire/no-fire record, context at fire time \\
O6 ASSERT & claim carries status and a confidence score with a declared elicitation method; the
store can report its own gaps & status, score plus prompt, outcome, store-state self-report \\
\bottomrule
\end{tabular}
\caption{The six operations, their postconditions and the evidence that makes them checkable.}
\label{tab:ops}
\end{table}

Table~\ref{tab:mapping} maps seven existing systems onto the taxonomy. The result is the
argument in one view: three operation families are covered by two of the seven systems, and the
two cross-cutting axes are partial or absent in five of seven. Neither fact is visible in an
end-to-end accuracy number.

\begin{table}[t]
\centering
\small
\begin{tabular}{@{}>{\raggedright\arraybackslash}p{4.55cm}cccccccc@{}}
\toprule
System & O1 & O2 & O3 & O4 & O5 & O6 & Prov. & Lifecycle \\
\midrule
MemGPT / Letta \cite{packer2023memgpt,lettadocs} & + & + & $\sim$ & $\sim$ & -- & $\sim$ & $\sim$ & -- \\
Generative Agents \cite{park2023generative} & + & + & + & -- & -- & -- & $\sim$ & -- \\
MemoryBank \cite{zhong2023memorybank} & + & + & + & -- & -- & -- & -- & $\sim$ \\
A-MEM \cite{xu2025amem} & + & + & $\sim$ & $\sim$ & -- & -- & $\sim$ & -- \\
Zep / Graphiti \cite{rasmussen2025zep,graphiti} & + & + & $\sim$ & + & -- & -- & + & + \\
Mem0 \cite{chhikara2025mem0,mem0docs} & + & + & + & $\sim$ & -- & -- & $\sim$ & $\sim$ \\
Liquid / Epistemic / Dream / Prospective \cite{sintliquid,sintepistemic,sintdream,sintprospective} & + & + & + & + & + & + & $\sim$ & + \\
\bottomrule
\end{tabular}
\caption{Operation coverage of seven systems. \texttt{+} present and observable, $\sim$
partially present, \texttt{--} absent.}
\label{tab:mapping}
\end{table}

\paragraph{Differences from the closest prior work.} \cite{huan2026secondhalf} orders the field
by memory type and reviews metrics, but defines no runnable procedure.
\cite{zhou2026agentnative} decomposes the system into four internal modules and measures twelve
systems, but its boundaries are internal: a third party cannot run that framework against a
server it did not implement. \cite{zhang2024memorysurvey} predates the memory-server ecosystem
and describes operations per system rather than as testable postconditions.
\cite{luo2025mcpbench} evaluates MCP servers, but for tool-use effectiveness rather than memory
semantics. \cite{wang2026rightmemory} verifies one property (retrieval admissibility); we adopt
it as a task family instead of duplicating it. This positioning is a claim of design
difference, not of empirical superiority (claims C1--C5 in \texttt{claims.md}).

\section{Evaluation protocol}
\label{sec:protocol}

The protocol specifies task families, metrics, a run procedure and a report format; the schema
and validator ship with the paper, so a report that omits a mandatory field is rejected rather
than published.

\paragraph{Task families.} T1 write fidelity; T2 read accuracy and ranking; T3 temporal
reasoning; T4 admissibility (principal, policy, lifecycle); T5 update and conflict; T6
consolidation retention (identical probe set before and after); T7 prospective memory (cue
firing, suspend/resume, cross-session persistence); T8 epistemic status and calibration; T9
durability and determinism across restarts; T10 adversarial and privacy. Rationale: T4 follows
the finding that recall metrics cannot detect inadmissible retrieval
\cite{wang2026rightmemory}; T6 exists separately from T2 because retention failure precedes
retrieval \cite{song2026eviction}; T7 is adopted from prospective-memory benchmarking
\cite{liu2026pmbench,einstein1990prospective}; T8 requires status, elicitation method and
outcome to be stored together \cite{ni2024probabilistic,yang2024verbalized}; T10 follows
demonstrated memory poisoning and membership inference
\cite{sunil2026mempoison,chen2026mrmmia}.

\paragraph{Metrics and who computes them.} Each metric is defined by a formula, a unit and a
computer: the runner harness, an independent programmatic check over returned records and
store state, or --- only where a program cannot judge --- a two-grader rubric with Cohen's
$\kappa$. No language-model judge is used for a metric that code can compute. Headline metrics
include write precision/recall and duplicate rate; recall@$k$ and MRR; provenance correctness;
admissibility error rate; context inflation (retrieved characters per required character, the
protocol's proxy in a run with no tokenizer); probe survival and compression ratio;
contradiction detection $F_1$, revision correctness and \textbf{stale-read rate}; cue hit rate
and false-fire rate; expected calibration error \cite{guo2017calibration} \emph{together with}
discrimination (AUROC), since neither alone is informative \cite{xiong2023confidence}; and cost
in tokens and USD against a price table whose date is part of the report.

\paragraph{Procedure.} Servers are selected by a stated rule and every candidate, included or
not, is listed, to make the sample auditable \cite{afsar2026registry}. Each run pins the server
commit, the MCP revision, the transport, the runtime, the model parameters and the seed set;
runs are repeated at least three times, and dispersion is reported. Evidence position
\cite{liu2023lostmiddle,hseih2024foundmiddle} and the context budget \cite{xu2026contextpipe}
are controlled and recorded. Adversarial streams run last. Failure codes (server error,
timeout, budget exhausted, not implemented, nondeterministic, leakage, stale read, lost
intention) are recorded outcomes, not missing rows.

\paragraph{Report and disclosure.} The report is JSON, validated against
\texttt{protocol\_schema.json}; it carries protocol version, subject identity including commit
SHA, runner configuration, budgets with a dated price table, per-family metrics with bootstrap
confidence intervals, raw log paths so every number can be recomputed, failure counts, and a
mandatory disclosure block (selection rule, candidates, prompts hash, environment, known
deviations) motivated by the finding that benchmark papers systematically omit exactly those
fields \cite{moghadasi2026disclose}. The format is exercised end to end by the two reports of
Section~\ref{sec:compare}: each passes the validator, and each carries the candidate list with
the ground for every inclusion and exclusion, plus its own list of stated deviations.

\section{Reference implementation}
\label{sec:impl}

The reference implementation (\cite{sintrepo}) exposes five of the six operations as public MCP
servers: liquid memory with contradiction handling and revision chains \cite{sintliquid},
epistemic status with calibrated confidence \cite{sintepistemic}, offline consolidation as
``dreaming'' \cite{sintdream}, prospective memory with intention stacks and session snapshots
\cite{sintprospective}, and a role-based review swarm containing provenance, temporal and
uncertainty roles \cite{sintswarm}. Commit-pinned facts (tool counts, file inventory, last
commit) were collected read-only into \texttt{data/repo-facts.json}.

The same artefact carries the paper's known gaps. The operations live in separate processes, so
a record's provenance is not shared automatically. Derivation is carried as free text rather
than an exportable PROV-O graph \cite{prov_o}. Cost instrumentation (tokens, latency) is not
implemented, although the literature treats cost as a first-class design variable
\cite{song2026dmem,xu2026linkedin}. Baselines exist only where a server could be launched and
measured: of the servers above, the liquid store is measured in Section~\ref{sec:compare}, while
the epistemic, dreaming and prospective servers could not be launched on the measurement host
and remain unmeasured. These gaps are precisely why the protocol makes cost, provenance and
restart probes mandatory rather than optional.

\section{Worked comparison: two memory servers under one protocol}
\label{sec:compare}

\paragraph{The pair of runs.} Two servers from the reference set were exercised with the
protocol of Section~\ref{sec:protocol}: \valtl{liquid.identity.subject} and
\valtl{memory.identity.subject}, each pinned to a commit (Table~\ref{tab:run-identity}), each over
\val{run.families} task families, \val{run.repetitions} repetitions per family and the seed set
\valt{run.seeds} (\val{run.family_tasks} tasks per repetition), driven over MCP stdio by the
harness named in the same table. Both reports pass the schema validator, and every number below
is read from them. This section is a \emph{pair of profiles}, not a ranking: the subjects differ
in storage design by construction --- one is a belief store with revisions, the other an
append-only hash-chained block store, as the candidate list in each report states --- and the
sample is two servers.

\paragraph{How the tables were produced.} Every table in this section is generated from the two
\texttt{report.json} files by \texttt{refslint tables}, and every number in the sentences is a
macro defined in the same generated file; \texttt{data/run-numbers.tsv} lists each value with the
JSON path it came from. A number that has no field in a report cannot be printed here, and
\texttt{refslint tables --check} fails if the generated files no longer match the reports.
Values are shown rounded to at most four decimals. \textit{null} is a metric that the protocol
defines and that the subject's tool surface cannot produce, with the reason in the report's own
note for that metric; $n$ is the number of repetitions the metric could be computed over.
Table~\ref{tab:run-identity} gives the identity of the two runs,
Tables~\ref{tab:t1}--\ref{tab:t10} the per-family profiles, and Table~\ref{tab:ops-profile} the
same reading pooled by taxonomy operation.

\paragraph{What is invariant across the two subjects.} The protocol reproduces identical
readings where the mechanism is the same. Durability and determinism are complete in both:
restart survival \val{liquid.T9.restart_survival}, snapshot-reopen survival
\val{liquid.T9.snapshot_reopen_survival} and repeated-query equality
\val{liquid.T9.repeat_query_equality}. Provenance correctness of returned records is
\val{liquid.T2.provenance_correctness} in both, abstention on out-of-coverage questions
\val{liquid.T8.abstention_correctness} in both, and in both the store's self-reported counters
agree with the harness write log (monitoring accuracy \val{liquid.T8.monitoring_accuracy}).
Query-only interactions changed no later answer for the target record (poisoning influence
\val{liquid.T10.poisoning_influence_rate} in both). Membership inference against a record's own
writer succeeds in both (\val{liquid.T10.membership_inference_advantage}; the same principal
wrote the records, so this is store-level inferability, not cross-principal leakage). Two
families are \emph{not implemented in either subject}: no advertised tool accepts an as-of date
(\texttt{temporal\_correctness} \val{liquid.T3.temporal_correctness}) or a caller identity
(\texttt{principal\_scoping\_exposed} \val{liquid.T4.principal_scoping_exposed}), recorded as
\val{liquid.failures.not_implemented} and \val{memory.failures.not_implemented}
\texttt{not\_implemented} failures respectively. Both expose timestamps on returned records
(\texttt{temporal\_metadata\_exposed} \val{liquid.T3.temporal_metadata_exposed}) but no lifecycle
field (\texttt{lifecycle\_state\_exposed} \val{liquid.T4.lifecycle_state_exposed}), and the policy
check over returned records --- the tool surface has no caller parameter --- gives the same
admissibility error rate in both, \val{liquid.T4.admissibility_error_rate}.

\paragraph{Where the profiles diverge, and in which direction.} On the write path the two
subjects are not interchangeable. Write precision is \val{liquid.T1.write_precision} for
sint-liquid-mcp and \val{memory.T1.write_precision} for sint-memory-mcp, and the shortfall is
fully accounted for by retries: the duplicate rate is \val{liquid.T1.duplicate_rate} against
\val{memory.T1.duplicate_rate}, and the programmatic check counts
\val{memory.T1.idempotency_violations} idempotency violations against
\val{liquid.T1.idempotency_violations} --- at equal write recall \val{liquid.T1.write_recall} in
both, so nothing was lost and four records were stored twice under repeated intent. Read-side
ranking runs the other way: recall@5 is \val{memory.T2.recall_at_5} for sint-memory-mcp against
\val{liquid.T2.recall_at_5}, and MRR \val{memory.T2.mrr} (\valci{memory.T2.mrr}) against
\val{liquid.T2.mrr}. The pair is informative precisely because one run inverts the other: the
subject with the cleaner write path is the one that misses a topic-only query
(\texttt{topic\_query\_hit\_rate} \val{liquid.T2.topic_query_hit_rate} against
\val{memory.T2.topic_query_hit_rate}), and the subject that returns
\val{memory.T2.context_inflation} times the required characters is the one that leaves duplicate
records behind. Write latency is close (\val{liquid.T1.write_latency_p50_ms}~ms against
\val{memory.T1.write_latency_p50_ms}~ms at p50); read latency is not
(\val{liquid.T2.read_latency_p95_ms}~ms against \val{memory.T2.read_latency_p95_ms}~ms at p95),
and is the weakest number here: wall clock around a stdio call on one workstation under
unspecified load, with three repetitions.

\paragraph{Revision: an operation-level difference that retrieval scores do not show.} In T5 the
harness writes contradictory pairs and then asks what survives. Neither subject flags a conflict
by itself: \texttt{contradiction\_detection\_f1} is
\val{liquid.T5.contradiction_detection_f1} in both, and since the harness never names the
conflicting pair, this measures detection rather than registration. sint-liquid-mcp then applies
the revision it is asked for (revision correctness \val{liquid.T5.revision_correctness},
revision-chain completeness \val{liquid.T5.chain_completeness}) and leaves unrelated records
untouched (collateral change \val{liquid.T5.collateral_change_rate}) --- yet its read path still
serves the superseded content (stale-read rate \val{liquid.T5.stale_read_rate}) and returns both
competing contents in \val{liquid.T5.both_competing_contents_served_rate} of the conflicts.
sint-memory-mcp is append-only, so it has no revision to be correct about (revision correctness
\textit{null}, chain completeness \val{memory.T5.chain_completeness}), and it too serves the
superseded content (\val{memory.T5.stale_read_rate}), with both contents returned in
\val{memory.T5.both_competing_contents_served_rate} of the conflicts. Read together, the two
profiles say one thing about revision as an interface: a store can accept a revision and still
serve the superseded record, and T5 sees that difference where a retrieval score would not. We
state it as an observation of two runs, not as a property of either design.

\paragraph{Prospective memory and calibration: where the pair disagrees most.} T7 asks the store
for a fact through a cue-only query after a process restart; neither subject has an intention
surface, so firing is operationalised as store-mediated cue retrieval. The cue hit rate is
\val{liquid.T7.cue_hit_rate} for sint-liquid-mcp and \val{memory.T7.cue_hit_rate} for
sint-memory-mcp, with false fires \val{liquid.T7.false_fire_rate} in both; the stored context is
returned with the fire in \val{memory.T7.context_availability_at_fire} of the fires of
sint-memory-mcp, and the same metric is \textit{null} for sint-liquid-mcp only because there is
no fire to check. In T8 both halves of calibration are reported and they do not agree: expected
calibration error \val{liquid.T8.ece} against \val{memory.T8.ece}, discrimination (AUROC)
\val{liquid.T8.discrimination_auroc} against \val{memory.T8.discrimination_auroc}, with top-ranked
answer correctness \val{liquid.T8.in_coverage_answer_correctness} against
\val{memory.T8.in_coverage_answer_correctness}. The confidence reported here is the highest
confidence among returned records (Section~\ref{sec:protocol}), so these are readings of the
stores' own confidence fields, not of verbalised model confidence. With \val{run.repetitions}
repetitions and one task set per family, a pair this far apart is an observation to be
re-measured on more subjects before anything is said about calibration in general.

\paragraph{Hypotheses stated before the run.} H1 --- explicit lifecycle metadata buys a lower
stale-read rate at a measurable write-latency cost --- is not testable on this pair: neither
subject exposes a lifecycle field on returned records, so the antecedent is absent, and both
serve stale content at the same rate \val{liquid.T5.stale_read_rate}. H2 --- a subject without a
prospective-memory operation fails T7 entirely while its T2 and T6 scores remain high --- holds
for one of the two: sint-liquid-mcp fails the cue-only family (\val{liquid.T7.cue_hit_rate})
while scoring read recall \val{liquid.T2.recall_at_5} and \textit{null} on the consolidation
family it does not implement; sint-memory-mcp answers the cue query
(\val{memory.T7.cue_hit_rate}) through its substring read path. One subject consistent with a
hypothesis and one not is what a two-subject sample can produce; the hypothesis is neither
confirmed nor refuted, and its status in \texttt{claims.md} is unchanged.

\paragraph{What this comparison does not measure.} Token and cost metrics are \textit{null} by
construction in both reports: no model is in the loop, so there is no token accounting and no
price table applies, as the reports state in \texttt{budgets.price\_table}. Neither subject was
measured against a third-party baseline, and neither was measured for availability, concurrency
or behaviour under load. Consolidation (O3) is \textit{null} for sint-liquid-mcp, which has no
maintenance operation, and only partly populated for sint-memory-mcp, where an offline decay
pass stood in for the consolidation the protocol asks for (probe survival
\val{memory.T6.probe_survival_rate}, compression ratio \val{memory.T6.compression_ratio}); the
consolidation-specific metrics stay \textit{null} in both. The one rubric metric
(\texttt{abstraction\_faithfulness}) was never graded, so no $\kappa$ is reported. A
\textit{null} therefore means one of two different things --- the operation does not exist, or it
exists but the metric has nothing to score --- and each report's own note says which.

\begingroup\small
\begin{longtable}{@{}> {\raggedright\arraybackslash}p{3.40cm}> {\raggedright\arraybackslash}p{5.20cm}> {\raggedright\arraybackslash}p{5.20cm}@{}}
\caption{Identity of the two runs: subject, pinned revision, runtime, budgets and the raw responses every number of this section is computed from. Every value is read from the corresponding \texttt{report.json}; the table is generated by \texttt{refslint tables}.}
\label{tab:run-identity}\\
\toprule
field & sint-liquid-mcp & sint-memory-mcp \\
\midrule
\endfirsthead
\toprule
field & sint-liquid-mcp & sint-memory-mcp \\
\midrule
\endhead
subject & {\footnotesize sint-liqui\hspace{0pt}d-mcp} & {\footnotesize sint-memor\hspace{0pt}y-mcp} \\
repository & \url{https://github.com/kostyk348/sint-liquid-mcp} & \url{https://github.com/kostyk348/sint} \\
commit & {\footnotesize bfffce2cde\hspace{0pt}c35c11aa0d\hspace{0pt}90ae02689d\hspace{0pt}69df596b11\hspace{0pt}} & {\footnotesize 7554f7eef8\hspace{0pt}53cebafae3\hspace{0pt}4e6654acc9\hspace{0pt}871fb01163\hspace{0pt}} \\
published version & {\footnotesize MCP server as published in the repository\hspace{0pt} snapshot (no release tag)} & {\footnotesize MCP server as published in the repository\hspace{0pt} snapshot (no release tag)} \\
transport / MCP revision & {\footnotesize stdio /\hspace{0pt} 2025-03-26\hspace{0pt}} & {\footnotesize stdio /\hspace{0pt} 2025-03-26\hspace{0pt}} \\
runtime & {\footnotesize Python 3.14.7,\hspace{0pt} mcp 1.9.4 (mcp.serve\hspace{0pt}r.fastmcp)\hspace{0pt}} & {\footnotesize Python 3.14.7,\hspace{0pt} mcp 1.9.4 (mcp.serve\hspace{0pt}r.fastmcp)\hspace{0pt}} \\
deployment & {\footnotesize local process,\hspace{0pt} fresh data directory per repetition\hspace{0pt} (USERPROFI\hspace{0pt}LE/\hspace{0pt}HOME redirected\hspace{0pt})} & {\footnotesize local process,\hspace{0pt} fresh data directory per repetition\hspace{0pt} (USERPROFI\hspace{0pt}LE/\hspace{0pt}HOME redirected\hspace{0pt})} \\
store kept in the run dir & {\footnotesize reports/\hspace{0pt}sint-liqui\hspace{0pt}d-mcp-2026\hspace{0pt}-10-16/\hspace{0pt}data/\hspace{0pt}seed<seed>\hspace{0pt}} & {\footnotesize reports/\hspace{0pt}sint-memor\hspace{0pt}y-mcp-2026\hspace{0pt}-10-16/\hspace{0pt}data/\hspace{0pt}seed<seed>\hspace{0pt}} \\
model under test & {\footnotesize none (store-onl\hspace{0pt}y)} & {\footnotesize none (store-onl\hspace{0pt}y)} \\
hardware & {\footnotesize AMD64,\hspace{0pt} 4 logical CPUs (workstati\hspace{0pt}on class,\hspace{0pt} not a benchmark host)} & {\footnotesize AMD64,\hspace{0pt} 4 logical CPUs (workstati\hspace{0pt}on class,\hspace{0pt} not a benchmark host)} \\
runner / harness & {\footnotesize uncommitte\hspace{0pt}d working tree:\hspace{0pt} sha256:\hspace{0pt}9b79d0aa12\hspace{0pt}ae1660 over code/\hspace{0pt}run\_\hspace{0pt}protocol.p\hspace{0pt}y + code/\hspace{0pt}streams.py\hspace{0pt} + code/\hspace{0pt}servers.js\hspace{0pt}on (this workspace has no git metadata,\hspace{0pt} so there is no commit to name)} & {\footnotesize uncommitte\hspace{0pt}d working tree:\hspace{0pt} sha256:\hspace{0pt}9b79d0aa12\hspace{0pt}ae1660 over code/\hspace{0pt}run\_\hspace{0pt}protocol.p\hspace{0pt}y + code/\hspace{0pt}streams.py\hspace{0pt} + code/\hspace{0pt}servers.js\hspace{0pt}on (this workspace has no git metadata,\hspace{0pt} so there is no commit to name)} \\
environment & {\footnotesize Windows-11\hspace{0pt}-10.0.2620\hspace{0pt}0-SP0;\hspace{0pt} Intel64 Family 6 Model 140 Stepping 2,\hspace{0pt} GenuineInt\hspace{0pt}el;\hspace{0pt} server interprete\hspace{0pt}r C:\hspace{0pt}/\hspace{0pt}Users/\hspace{0pt}Lab/\hspace{0pt}AppData/\hspace{0pt}Local/\hspace{0pt}hermes/\hspace{0pt}installs/\hspace{0pt}2636d87a16\hspace{0pt}39ae64/\hspace{0pt}environmen\hspace{0pt}ts/\hspace{0pt}f32597c25d\hspace{0pt}734363b2d2\hspace{0pt}20b038061d\hspace{0pt}05/\hspace{0pt}venv/\hspace{0pt}Scripts/\hspace{0pt}python.exe\hspace{0pt};\hspace{0pt} harness Python 3.14.7 (stdlib only)} & {\footnotesize Windows-11\hspace{0pt}-10.0.2620\hspace{0pt}0-SP0;\hspace{0pt} Intel64 Family 6 Model 140 Stepping 2,\hspace{0pt} GenuineInt\hspace{0pt}el;\hspace{0pt} server interprete\hspace{0pt}r C:\hspace{0pt}/\hspace{0pt}Users/\hspace{0pt}Lab/\hspace{0pt}AppData/\hspace{0pt}Local/\hspace{0pt}hermes/\hspace{0pt}installs/\hspace{0pt}2636d87a16\hspace{0pt}39ae64/\hspace{0pt}environmen\hspace{0pt}ts/\hspace{0pt}f32597c25d\hspace{0pt}734363b2d2\hspace{0pt}20b038061d\hspace{0pt}05/\hspace{0pt}venv/\hspace{0pt}Scripts/\hspace{0pt}python.exe\hspace{0pt};\hspace{0pt} harness Python 3.14.7 (stdlib only)} \\
repetitions / seeds & {\footnotesize 3 /\hspace{0pt} 1,\hspace{0pt} 2,\hspace{0pt} 3} & {\footnotesize 3 /\hspace{0pt} 1,\hspace{0pt} 2,\hspace{0pt} 3} \\
budgets per repetition & {\footnotesize 400 calls,\hspace{0pt} 900 s wall clock} & {\footnotesize 400 calls,\hspace{0pt} 900 s wall clock} \\
confidence intervals & {\footnotesize bootstrap\_\hspace{0pt}percentile\hspace{0pt},\hspace{0pt} 1000 resamples} & {\footnotesize bootstrap\_\hspace{0pt}percentile\hspace{0pt},\hspace{0pt} 1000 resamples} \\
failures recorded & {\footnotesize not\_\hspace{0pt}implemente\hspace{0pt}d x 15} & {\footnotesize not\_\hspace{0pt}implemente\hspace{0pt}d x 24} \\
raw responses & {\footnotesize reports/\hspace{0pt}sint-liqui\hspace{0pt}d-mcp-2026\hspace{0pt}-10-16/\hspace{0pt}raw (33 JSONL files)} & {\footnotesize reports/\hspace{0pt}sint-memor\hspace{0pt}y-mcp-2026\hspace{0pt}-10-16/\hspace{0pt}raw (33 JSONL files)} \\
report.json & {\footnotesize reports/\hspace{0pt}sint-liqui\hspace{0pt}d-mcp-2026\hspace{0pt}-10-16/\hspace{0pt}report.jso\hspace{0pt}n} & {\footnotesize reports/\hspace{0pt}sint-memor\hspace{0pt}y-mcp-2026\hspace{0pt}-10-16/\hspace{0pt}report.jso\hspace{0pt}n} \\
\bottomrule
\end{longtable}
\endgroup


\begingroup\small
\begin{longtable}{@{}> {\raggedright\arraybackslash}p{2.60cm}> {\raggedright\arraybackslash}p{4.27cm}> {\raggedright\arraybackslash}p{4.27cm}> {\raggedright\arraybackslash}p{2.60cm}@{}}
\caption{Task family T1: every metric the protocol defines for it that \texttt{report.json} carries, for both subjects, read side by side (12 tasks per repetition, 3 repetitions, seeds 1, 2, 3). Each subject cell gives the value, its 95\% bootstrap CI and $n$; \textit{null} is a metric that the protocol defines and that the subject's tool surface cannot produce, with the reason in the report's own note for that metric. Values are printed rounded to at most four decimals; the unrounded values and their JSON paths are in \texttt{data/run-numbers.tsv}. Generated by \texttt{refslint tables}.}\label{tab:t1}\\
\toprule
metric & sint-liquid-mcp & sint-memory-mcp & computer \\
\midrule
\endfirsthead
\toprule
metric & sint-liquid-mcp & sint-memory-mcp & computer \\
\midrule
\endhead
{\footnotesize\texttt{write\_\hspace{0pt}precision} (\textit{ratio})} & 1.0\newline [1.0, 1.0], $n$=3 & 0.75\newline [0.75, 0.75], $n$=3 & harness \\
{\footnotesize\texttt{write\_\hspace{0pt}recall} (\textit{ratio})} & 1.0\newline [1.0, 1.0], $n$=3 & 1.0\newline [1.0, 1.0], $n$=3 & harness \\
{\footnotesize\texttt{duplicate\_\hspace{0pt}rate} (\textit{ratio})} & 0.0\newline [0.0, 0.0], $n$=3 & 0.25\newline [0.25, 0.25], $n$=3 & harness \\
{\footnotesize\texttt{idempotency\_\hspace{0pt}violations} (\textit{count})} & 0.0\newline [0.0, 0.0], $n$=3 & 4.0\newline [4.0, 4.0], $n$=3 & programmatic\_\hspace{0pt}check \\
{\footnotesize\texttt{write\_\hspace{0pt}latency\_\hspace{0pt}p50\_\hspace{0pt}ms} (\textit{ms})} & 17.8175\newline [17.0155, 18.4985], $n$=3 & 16.4423\newline [15.167, 17.2955], $n$=3 & harness \\
{\footnotesize\texttt{write\_\hspace{0pt}latency\_\hspace{0pt}p95\_\hspace{0pt}ms} (\textit{ms})} & 21.0234\newline [19.3115, 21.8968], $n$=3 & 18.9019\newline [17.329, 20.2025], $n$=3 & harness \\
{\footnotesize\texttt{tokens\_\hspace{0pt}per\_\hspace{0pt}write} (\textit{tokens})} & \textit{null} ($n$=0) & \textit{null} ($n$=0) & harness \\
{\footnotesize\texttt{cost\_\hspace{0pt}per\_\hspace{0pt}1000\_\hspace{0pt}writes\_\hspace{0pt}usd} (\textit{usd})} & \textit{null} ($n$=0) & \textit{null} ($n$=0) & harness \\
\bottomrule
\end{longtable}
\endgroup


\begingroup\small
\begin{longtable}{@{}> {\raggedright\arraybackslash}p{2.60cm}> {\raggedright\arraybackslash}p{4.27cm}> {\raggedright\arraybackslash}p{4.27cm}> {\raggedright\arraybackslash}p{2.60cm}@{}}
\caption{Task family T2: every metric the protocol defines for it that \texttt{report.json} carries, for both subjects, read side by side (6 tasks per repetition, 3 repetitions, seeds 1, 2, 3). Each subject cell gives the value, its 95\% bootstrap CI and $n$; \textit{null} is a metric that the protocol defines and that the subject's tool surface cannot produce, with the reason in the report's own note for that metric. Values are printed rounded to at most four decimals; the unrounded values and their JSON paths are in \texttt{data/run-numbers.tsv}. Generated by \texttt{refslint tables}.}\label{tab:t2}\\
\toprule
metric & sint-liquid-mcp & sint-memory-mcp & computer \\
\midrule
\endfirsthead
\toprule
metric & sint-liquid-mcp & sint-memory-mcp & computer \\
\midrule
\endhead
{\footnotesize\texttt{recall\_\hspace{0pt}at\_\hspace{0pt}5} (\textit{ratio})} & 0.5\newline [0.5, 0.5], $n$=3 & 1.0\newline [1.0, 1.0], $n$=3 & harness \\
{\footnotesize\texttt{required\_\hspace{0pt}record\_\hspace{0pt}hit\_\hspace{0pt}rate} (\textit{ratio})} & 0.5\newline [0.5, 0.5], $n$=3 & 1.0\newline [1.0, 1.0], $n$=3 & harness \\
{\footnotesize\texttt{mrr} (\textit{ratio})} & 0.5\newline [0.5, 0.5], $n$=3 & 0.75\newline [0.6806, 0.8056], $n$=3 & harness \\
{\footnotesize\texttt{key\_\hspace{0pt}query\_\hspace{0pt}hit\_\hspace{0pt}rate} (\textit{ratio})} & 1.0\newline [1.0, 1.0], $n$=3 & 1.0\newline [1.0, 1.0], $n$=3 & harness \\
{\footnotesize\texttt{topic\_\hspace{0pt}query\_\hspace{0pt}hit\_\hspace{0pt}rate} (\textit{ratio})} & 0.0\newline [0.0, 0.0], $n$=3 & 1.0\newline [1.0, 1.0], $n$=3 & harness \\
{\footnotesize\texttt{provenance\_\hspace{0pt}correctness} (\textit{ratio})} & 1.0\newline [1.0, 1.0], $n$=3 & 1.0\newline [1.0, 1.0], $n$=3 & harness \\
{\footnotesize\texttt{context\_\hspace{0pt}inflation} (\textit{ratio})} & 0.5\newline [0.5, 0.5], $n$=3 & 2.5\newline [2.5, 2.5], $n$=3 & harness \\
{\footnotesize\texttt{read\_\hspace{0pt}latency\_\hspace{0pt}p50\_\hspace{0pt}ms} (\textit{ms})} & 4.4998\newline [4.1825, 4.9225], $n$=3 & 8.0937\newline [7.6565, 8.5435], $n$=3 & harness \\
{\footnotesize\texttt{read\_\hspace{0pt}latency\_\hspace{0pt}p95\_\hspace{0pt}ms} (\textit{ms})} & 13.7778\newline [12.8863, 14.3423], $n$=3 & 102.8932\newline [101.5778, 104.9318], $n$=3 & harness \\
\bottomrule
\end{longtable}
\endgroup


\begingroup\small
\begin{longtable}{@{}> {\raggedright\arraybackslash}p{2.60cm}> {\raggedright\arraybackslash}p{4.27cm}> {\raggedright\arraybackslash}p{4.27cm}> {\raggedright\arraybackslash}p{2.60cm}@{}}
\caption{Task family T3: every metric the protocol defines for it that \texttt{report.json} carries, for both subjects, read side by side (4 tasks per repetition, 3 repetitions, seeds 1, 2, 3). Each subject cell gives the value, its 95\% bootstrap CI and $n$; \textit{null} is a metric that the protocol defines and that the subject's tool surface cannot produce, with the reason in the report's own note for that metric. Values are printed rounded to at most four decimals; the unrounded values and their JSON paths are in \texttt{data/run-numbers.tsv}. Generated by \texttt{refslint tables}.}\label{tab:t3}\\
\toprule
metric & sint-liquid-mcp & sint-memory-mcp & computer \\
\midrule
\endfirsthead
\toprule
metric & sint-liquid-mcp & sint-memory-mcp & computer \\
\midrule
\endhead
{\footnotesize\texttt{temporal\_\hspace{0pt}correctness} (\textit{ratio})} & \textit{null} ($n$=0) & \textit{null} ($n$=0) & harness \\
{\footnotesize\texttt{temporal\_\hspace{0pt}metadata\_\hspace{0pt}exposed} (\textit{ratio})} & 1.0\newline [1.0, 1.0], $n$=3 & 1.0\newline [1.0, 1.0], $n$=3 & programmatic\_\hspace{0pt}check \\
\bottomrule
\end{longtable}
\endgroup


\begingroup\small
\begin{longtable}{@{}> {\raggedright\arraybackslash}p{2.60cm}> {\raggedright\arraybackslash}p{4.27cm}> {\raggedright\arraybackslash}p{4.27cm}> {\raggedright\arraybackslash}p{2.60cm}@{}}
\caption{Task family T4: every metric the protocol defines for it that \texttt{report.json} carries, for both subjects, read side by side (6 tasks per repetition, 3 repetitions, seeds 1, 2, 3). Each subject cell gives the value, its 95\% bootstrap CI and $n$; \textit{null} is a metric that the protocol defines and that the subject's tool surface cannot produce, with the reason in the report's own note for that metric. Values are printed rounded to at most four decimals; the unrounded values and their JSON paths are in \texttt{data/run-numbers.tsv}. Generated by \texttt{refslint tables}.}\label{tab:t4}\\
\toprule
metric & sint-liquid-mcp & sint-memory-mcp & computer \\
\midrule
\endfirsthead
\toprule
metric & sint-liquid-mcp & sint-memory-mcp & computer \\
\midrule
\endhead
{\footnotesize\texttt{admissibility\_\hspace{0pt}error\_\hspace{0pt}rate} (\textit{ratio})} & 0.5\newline [0.5, 0.5], $n$=3 & 0.5\newline [0.5, 0.5], $n$=3 & programmatic\_\hspace{0pt}check \\
{\footnotesize\texttt{lifecycle\_\hspace{0pt}state\_\hspace{0pt}exposed} (\textit{ratio})} & 0.0\newline [0.0, 0.0], $n$=3 & 0.0\newline [0.0, 0.0], $n$=3 & programmatic\_\hspace{0pt}check \\
{\footnotesize\texttt{principal\_\hspace{0pt}scoping\_\hspace{0pt}exposed} (\textit{ratio})} & 0.0\newline [0.0, 0.0], $n$=3 & 0.0\newline [0.0, 0.0], $n$=3 & programmatic\_\hspace{0pt}check \\
\bottomrule
\end{longtable}
\endgroup


\begingroup\small
\begin{longtable}{@{}> {\raggedright\arraybackslash}p{2.60cm}> {\raggedright\arraybackslash}p{4.27cm}> {\raggedright\arraybackslash}p{4.27cm}> {\raggedright\arraybackslash}p{2.60cm}@{}}
\caption{Task family T5: every metric the protocol defines for it that \texttt{report.json} carries, for both subjects, read side by side (4 tasks per repetition, 3 repetitions, seeds 1, 2, 3). Each subject cell gives the value, its 95\% bootstrap CI and $n$; \textit{null} is a metric that the protocol defines and that the subject's tool surface cannot produce, with the reason in the report's own note for that metric. Values are printed rounded to at most four decimals; the unrounded values and their JSON paths are in \texttt{data/run-numbers.tsv}. Generated by \texttt{refslint tables}.}\label{tab:t5}\\
\toprule
metric & sint-liquid-mcp & sint-memory-mcp & computer \\
\midrule
\endfirsthead
\toprule
metric & sint-liquid-mcp & sint-memory-mcp & computer \\
\midrule
\endhead
{\footnotesize\texttt{contradiction\_\hspace{0pt}detection\_\hspace{0pt}f1} (\textit{ratio})} & 0.0\newline [0.0, 0.0], $n$=3 & 0.0\newline [0.0, 0.0], $n$=3 & harness \\
{\footnotesize\texttt{revision\_\hspace{0pt}correctness} (\textit{ratio})} & 1.0\newline [1.0, 1.0], $n$=3 & \textit{null} ($n$=0) & programmatic\_\hspace{0pt}check \\
{\footnotesize\texttt{chain\_\hspace{0pt}completeness} (\textit{ratio})} & 1.0\newline [1.0, 1.0], $n$=3 & 0.0\newline [0.0, 0.0], $n$=3 & programmatic\_\hspace{0pt}check \\
{\footnotesize\texttt{stale\_\hspace{0pt}read\_\hspace{0pt}rate} (\textit{ratio})} & 1.0\newline [1.0, 1.0], $n$=3 & 1.0\newline [1.0, 1.0], $n$=3 & harness \\
{\footnotesize\texttt{both\_\hspace{0pt}competing\_\hspace{0pt}contents\_\hspace{0pt}served\_\hspace{0pt}rate} (\textit{ratio})} & 1.0\newline [1.0, 1.0], $n$=3 & 1.0\newline [1.0, 1.0], $n$=3 & harness \\
{\footnotesize\texttt{policy\_\hspace{0pt}winner\_\hspace{0pt}top\_\hspace{0pt}ranked\_\hspace{0pt}rate} (\textit{ratio})} & \textit{null} ($n$=0) & 0.5\newline [0.5, 0.5], $n$=3 & harness \\
{\footnotesize\texttt{collateral\_\hspace{0pt}change\_\hspace{0pt}rate} (\textit{ratio})} & 0.0\newline [0.0, 0.0], $n$=3 & 0.0\newline [0.0, 0.0], $n$=3 & programmatic\_\hspace{0pt}check \\
\bottomrule
\end{longtable}
\endgroup


\begingroup\small
\begin{longtable}{@{}> {\raggedright\arraybackslash}p{2.60cm}> {\raggedright\arraybackslash}p{4.27cm}> {\raggedright\arraybackslash}p{4.27cm}> {\raggedright\arraybackslash}p{2.60cm}@{}}
\caption{Task family T6: every metric the protocol defines for it that \texttt{report.json} carries, for both subjects, read side by side (4 tasks per repetition, 3 repetitions, seeds 1, 2, 3). Each subject cell gives the value, its 95\% bootstrap CI and $n$; \textit{null} is a metric that the protocol defines and that the subject's tool surface cannot produce, with the reason in the report's own note for that metric. Values are printed rounded to at most four decimals; the unrounded values and their JSON paths are in \texttt{data/run-numbers.tsv}. Generated by \texttt{refslint tables}.}\label{tab:t6}\\
\toprule
metric & sint-liquid-mcp & sint-memory-mcp & computer \\
\midrule
\endfirsthead
\toprule
metric & sint-liquid-mcp & sint-memory-mcp & computer \\
\midrule
\endhead
{\footnotesize\texttt{probe\_\hspace{0pt}survival\_\hspace{0pt}rate} (\textit{ratio})} & \textit{null} ($n$=0) & 1.0\newline [1.0, 1.0], $n$=3 & harness \\
{\footnotesize\texttt{compression\_\hspace{0pt}ratio} (\textit{ratio})} & \textit{null} ($n$=0) & 1.0\newline [1.0, 1.0], $n$=3 & harness \\
{\footnotesize\texttt{derived\_\hspace{0pt}from\_\hspace{0pt}completeness} (\textit{ratio})} & \textit{null} ($n$=0) & \textit{null} ($n$=0) & programmatic\_\hspace{0pt}check \\
{\footnotesize\texttt{abstraction\_\hspace{0pt}faithfulness} (\textit{score\_\hspace{0pt}1\_\hspace{0pt}5})} & \textit{null} ($n$=0) & \textit{null} ($n$=0) & rubric\_\hspace{0pt}two\_\hspace{0pt}graders \\
{\footnotesize\texttt{cost\_\hspace{0pt}per\_\hspace{0pt}consolidation\_\hspace{0pt}usd} (\textit{usd})} & \textit{null} ($n$=0) & \textit{null} ($n$=0) & harness \\
\bottomrule
\end{longtable}
\endgroup


\begingroup\small
\begin{longtable}{@{}> {\raggedright\arraybackslash}p{2.60cm}> {\raggedright\arraybackslash}p{4.27cm}> {\raggedright\arraybackslash}p{4.27cm}> {\raggedright\arraybackslash}p{2.60cm}@{}}
\caption{Task family T7: every metric the protocol defines for it that \texttt{report.json} carries, for both subjects, read side by side (3 tasks per repetition, 3 repetitions, seeds 1, 2, 3). Each subject cell gives the value, its 95\% bootstrap CI and $n$; \textit{null} is a metric that the protocol defines and that the subject's tool surface cannot produce, with the reason in the report's own note for that metric. Values are printed rounded to at most four decimals; the unrounded values and their JSON paths are in \texttt{data/run-numbers.tsv}. Generated by \texttt{refslint tables}.}\label{tab:t7}\\
\toprule
metric & sint-liquid-mcp & sint-memory-mcp & computer \\
\midrule
\endfirsthead
\toprule
metric & sint-liquid-mcp & sint-memory-mcp & computer \\
\midrule
\endhead
{\footnotesize\texttt{cue\_\hspace{0pt}hit\_\hspace{0pt}rate} (\textit{ratio})} & 0.0\newline [0.0, 0.0], $n$=3 & 1.0\newline [1.0, 1.0], $n$=3 & harness \\
{\footnotesize\texttt{false\_\hspace{0pt}fire\_\hspace{0pt}rate} (\textit{ratio})} & 0.0\newline [0.0, 0.0], $n$=3 & 0.0\newline [0.0, 0.0], $n$=3 & harness \\
{\footnotesize\texttt{context\_\hspace{0pt}availability\_\hspace{0pt}at\_\hspace{0pt}fire} (\textit{ratio})} & \textit{null} ($n$=0) & 1.0\newline [1.0, 1.0], $n$=3 & programmatic\_\hspace{0pt}check \\
{\footnotesize\texttt{fire\_\hspace{0pt}latency\_\hspace{0pt}calls} (\textit{count})} & \textit{null} ($n$=0) & 0.0\newline [0.0, 0.0], $n$=3 & harness \\
\bottomrule
\end{longtable}
\endgroup


\begingroup\small
\begin{longtable}{@{}> {\raggedright\arraybackslash}p{2.60cm}> {\raggedright\arraybackslash}p{4.27cm}> {\raggedright\arraybackslash}p{4.27cm}> {\raggedright\arraybackslash}p{2.60cm}@{}}
\caption{Task family T8: every metric the protocol defines for it that \texttt{report.json} carries, for both subjects, read side by side (12 tasks per repetition, 3 repetitions, seeds 1, 2, 3). Each subject cell gives the value, its 95\% bootstrap CI and $n$; \textit{null} is a metric that the protocol defines and that the subject's tool surface cannot produce, with the reason in the report's own note for that metric. Values are printed rounded to at most four decimals; the unrounded values and their JSON paths are in \texttt{data/run-numbers.tsv}. Generated by \texttt{refslint tables}.}\label{tab:t8}\\
\toprule
metric & sint-liquid-mcp & sint-memory-mcp & computer \\
\midrule
\endfirsthead
\toprule
metric & sint-liquid-mcp & sint-memory-mcp & computer \\
\midrule
\endhead
{\footnotesize\texttt{abstention\_\hspace{0pt}correctness} (\textit{ratio})} & 1.0\newline [1.0, 1.0], $n$=3 & 1.0\newline [1.0, 1.0], $n$=3 & harness \\
{\footnotesize\texttt{in\_\hspace{0pt}coverage\_\hspace{0pt}answer\_\hspace{0pt}correctness} (\textit{ratio})} & 0.5\newline [0.5, 0.5], $n$=3 & 0.75\newline [0.75, 0.75], $n$=3 & harness \\
{\footnotesize\texttt{key\_\hspace{0pt}question\_\hspace{0pt}correctness} (\textit{ratio})} & 1.0\newline [1.0, 1.0], $n$=3 & 1.0\newline [1.0, 1.0], $n$=3 & harness \\
{\footnotesize\texttt{subject\_\hspace{0pt}question\_\hspace{0pt}correctness} (\textit{ratio})} & 0.0\newline [0.0, 0.0], $n$=3 & 0.5\newline [0.5, 0.5], $n$=3 & harness \\
{\footnotesize\texttt{ece} (\textit{ratio})} & 0.3667\newline [0.3667, 0.3667], $n$=3 & 0.4333\newline [0.4333, 0.4333], $n$=3 & harness \\
{\footnotesize\texttt{discrimination\_\hspace{0pt}auroc} (\textit{ratio})} & 0.75\newline [0.75, 0.75], $n$=3 & 0.3\newline [0.3, 0.3], $n$=3 & harness \\
{\footnotesize\texttt{monitoring\_\hspace{0pt}accuracy} (\textit{ratio})} & 1.0\newline [1.0, 1.0], $n$=3 & 1.0\newline [1.0, 1.0], $n$=3 & programmatic\_\hspace{0pt}check \\
\bottomrule
\end{longtable}
\endgroup


\begingroup\small
\begin{longtable}{@{}> {\raggedright\arraybackslash}p{2.60cm}> {\raggedright\arraybackslash}p{4.27cm}> {\raggedright\arraybackslash}p{4.27cm}> {\raggedright\arraybackslash}p{2.60cm}@{}}
\caption{Task family T9: every metric the protocol defines for it that \texttt{report.json} carries, for both subjects, read side by side (18 tasks per repetition, 3 repetitions, seeds 1, 2, 3). Each subject cell gives the value, its 95\% bootstrap CI and $n$; \textit{null} is a metric that the protocol defines and that the subject's tool surface cannot produce, with the reason in the report's own note for that metric. Values are printed rounded to at most four decimals; the unrounded values and their JSON paths are in \texttt{data/run-numbers.tsv}. Generated by \texttt{refslint tables}.}\label{tab:t9}\\
\toprule
metric & sint-liquid-mcp & sint-memory-mcp & computer \\
\midrule
\endfirsthead
\toprule
metric & sint-liquid-mcp & sint-memory-mcp & computer \\
\midrule
\endhead
{\footnotesize\texttt{restart\_\hspace{0pt}survival} (\textit{ratio})} & 1.0\newline [1.0, 1.0], $n$=3 & 1.0\newline [1.0, 1.0], $n$=3 & harness \\
{\footnotesize\texttt{snapshot\_\hspace{0pt}reopen\_\hspace{0pt}survival} (\textit{ratio})} & 1.0\newline [1.0, 1.0], $n$=3 & 1.0\newline [1.0, 1.0], $n$=3 & harness \\
{\footnotesize\texttt{repeat\_\hspace{0pt}query\_\hspace{0pt}equality} (\textit{ratio})} & 1.0\newline [1.0, 1.0], $n$=3 & 1.0\newline [1.0, 1.0], $n$=3 & harness \\
\bottomrule
\end{longtable}
\endgroup


\begingroup\small
\begin{longtable}{@{}> {\raggedright\arraybackslash}p{2.60cm}> {\raggedright\arraybackslash}p{4.27cm}> {\raggedright\arraybackslash}p{4.27cm}> {\raggedright\arraybackslash}p{2.60cm}@{}}
\caption{Task family T10: every metric the protocol defines for it that \texttt{report.json} carries, for both subjects, read side by side (12 tasks per repetition, 3 repetitions, seeds 1, 2, 3). Each subject cell gives the value, its 95\% bootstrap CI and $n$; \textit{null} is a metric that the protocol defines and that the subject's tool surface cannot produce, with the reason in the report's own note for that metric. Values are printed rounded to at most four decimals; the unrounded values and their JSON paths are in \texttt{data/run-numbers.tsv}. Generated by \texttt{refslint tables}.}\label{tab:t10}\\
\toprule
metric & sint-liquid-mcp & sint-memory-mcp & computer \\
\midrule
\endfirsthead
\toprule
metric & sint-liquid-mcp & sint-memory-mcp & computer \\
\midrule
\endhead
{\footnotesize\texttt{poisoning\_\hspace{0pt}influence\_\hspace{0pt}rate} (\textit{ratio})} & 0.0\newline [0.0, 0.0], $n$=3 & 0.0\newline [0.0, 0.0], $n$=3 & harness \\
{\footnotesize\texttt{membership\_\hspace{0pt}tpr} (\textit{ratio})} & 1.0\newline [1.0, 1.0], $n$=3 & 1.0\newline [1.0, 1.0], $n$=3 & harness \\
{\footnotesize\texttt{membership\_\hspace{0pt}fpr} (\textit{ratio})} & 0.0\newline [0.0, 0.0], $n$=3 & 0.0\newline [0.0, 0.0], $n$=3 & harness \\
{\footnotesize\texttt{membership\_\hspace{0pt}inference\_\hspace{0pt}advantage} (\textit{ratio})} & 1.0\newline [1.0, 1.0], $n$=3 & 1.0\newline [1.0, 1.0], $n$=3 & harness \\
\bottomrule
\end{longtable}
\endgroup


\begingroup\small
\begin{longtable}{@{}> {\raggedright\arraybackslash}p{2.60cm}> {\raggedright\arraybackslash}p{4.27cm}> {\raggedright\arraybackslash}p{4.27cm}> {\raggedright\arraybackslash}p{2.60cm}@{}}
\caption{The same profile by taxonomy operation: each operation's metrics pooled over the repetitions of the families that exercise it (the report states which), for both subjects side by side. \texttt{supported} is the subject's own verdict for that operation. Generated by \texttt{refslint tables}.}\label{tab:ops-profile}\\
\toprule
metric & sint-liquid-mcp & sint-memory-mcp & computer \\
\midrule
\endfirsthead
\toprule
metric & sint-liquid-mcp & sint-memory-mcp & computer \\
\midrule
\endhead
{\footnotesize\texttt{write\_\hspace{0pt}precision} (\textit{ratio})} & 1.0\newline [1.0, 1.0], $n$=3 & 0.75\newline [0.75, 0.75], $n$=3 & harness \\
{\footnotesize\texttt{write\_\hspace{0pt}recall} (\textit{ratio})} & 1.0\newline [1.0, 1.0], $n$=3 & 1.0\newline [1.0, 1.0], $n$=3 & harness \\
{\footnotesize\texttt{duplicate\_\hspace{0pt}rate} (\textit{ratio})} & 0.0\newline [0.0, 0.0], $n$=3 & 0.25\newline [0.25, 0.25], $n$=3 & harness \\
{\footnotesize\texttt{write\_\hspace{0pt}latency\_\hspace{0pt}p50\_\hspace{0pt}ms} (\textit{ms})} & 17.8175\newline [17.3232, 18.4985], $n$=3 & 16.4423\newline [15.167, 17.2955], $n$=3 & harness \\
{\footnotesize\texttt{write\_\hspace{0pt}latency\_\hspace{0pt}p95\_\hspace{0pt}ms} (\textit{ms})} & 21.0234\newline [19.3115, 21.8968], $n$=3 & 18.9019\newline [17.329, 20.2025], $n$=3 & harness \\
{\footnotesize\texttt{tokens\_\hspace{0pt}per\_\hspace{0pt}write} (\textit{tokens})} & \textit{null} ($n$=0) & \textit{null} ($n$=0) & harness \\
{\footnotesize\texttt{cost\_\hspace{0pt}per\_\hspace{0pt}1000\_\hspace{0pt}writes\_\hspace{0pt}usd} (\textit{usd})} & \textit{null} ($n$=0) & \textit{null} ($n$=0) & harness \\
{\footnotesize\texttt{idempotency\_\hspace{0pt}violations} (\textit{count})} & 0.0\newline [0.0, 0.0], $n$=3 & 4.0\newline [4.0, 4.0], $n$=3 & programmatic\_\hspace{0pt}check \\
{\footnotesize\texttt{recall\_\hspace{0pt}at\_\hspace{0pt}5} (\textit{ratio})} & 0.5\newline [0.5, 0.5], $n$=3 & 1.0\newline [1.0, 1.0], $n$=3 & harness \\
{\footnotesize\texttt{mrr} (\textit{ratio})} & 0.5\newline [0.5, 0.5], $n$=3 & 0.75\newline [0.6806, 0.8056], $n$=3 & harness \\
{\footnotesize\texttt{required\_\hspace{0pt}record\_\hspace{0pt}hit\_\hspace{0pt}rate} (\textit{ratio})} & 0.5\newline [0.5, 0.5], $n$=3 & 1.0\newline [1.0, 1.0], $n$=3 & harness \\
{\footnotesize\texttt{provenance\_\hspace{0pt}correctness} (\textit{ratio})} & 1.0\newline [1.0, 1.0], $n$=3 & 1.0\newline [1.0, 1.0], $n$=3 & harness \\
{\footnotesize\texttt{admissibility\_\hspace{0pt}error\_\hspace{0pt}rate} (\textit{ratio})} & 0.5\newline [0.5, 0.5], $n$=3 & 0.5\newline [0.5, 0.5], $n$=3 & programmatic\_\hspace{0pt}check \\
{\footnotesize\texttt{context\_\hspace{0pt}inflation} (\textit{ratio})} & 0.5\newline [0.5, 0.5], $n$=3 & 2.5\newline [2.5, 2.5], $n$=3 & harness \\
{\footnotesize\texttt{read\_\hspace{0pt}latency\_\hspace{0pt}p50\_\hspace{0pt}ms} (\textit{ms})} & 4.4998\newline [4.1825, 4.9225], $n$=3 & 8.0937\newline [7.6565, 8.5435], $n$=3 & harness \\
{\footnotesize\texttt{read\_\hspace{0pt}latency\_\hspace{0pt}p95\_\hspace{0pt}ms} (\textit{ms})} & 13.7778\newline [12.8863, 14.3423], $n$=3 & 102.8932\newline [101.5778, 104.9318], $n$=3 & harness \\
{\footnotesize\texttt{compression\_\hspace{0pt}ratio} (\textit{ratio})} & \textit{null} ($n$=0) & 1.0\newline [1.0, 1.0], $n$=3 & harness \\
{\footnotesize\texttt{probe\_\hspace{0pt}survival\_\hspace{0pt}rate} (\textit{ratio})} & \textit{null} ($n$=0) & 1.0\newline [1.0, 1.0], $n$=3 & harness \\
{\footnotesize\texttt{derived\_\hspace{0pt}from\_\hspace{0pt}completeness} (\textit{ratio})} & \textit{null} ($n$=0) & \textit{null} ($n$=0) & programmatic\_\hspace{0pt}check \\
{\footnotesize\texttt{cost\_\hspace{0pt}per\_\hspace{0pt}consolidation\_\hspace{0pt}usd} (\textit{usd})} & \textit{null} ($n$=0) & \textit{null} ($n$=0) & harness \\
{\footnotesize\texttt{contradiction\_\hspace{0pt}detection\_\hspace{0pt}f1} (\textit{ratio})} & 0.0\newline [0.0, 0.0], $n$=3 & 0.0\newline [0.0, 0.0], $n$=3 & harness \\
{\footnotesize\texttt{revision\_\hspace{0pt}correctness} (\textit{ratio})} & 1.0\newline [1.0, 1.0], $n$=3 & \textit{null} ($n$=0) & programmatic\_\hspace{0pt}check \\
{\footnotesize\texttt{chain\_\hspace{0pt}completeness} (\textit{ratio})} & 1.0\newline [1.0, 1.0], $n$=3 & 0.0\newline [0.0, 0.0], $n$=3 & programmatic\_\hspace{0pt}check \\
{\footnotesize\texttt{stale\_\hspace{0pt}read\_\hspace{0pt}rate} (\textit{ratio})} & 1.0\newline [1.0, 1.0], $n$=3 & 1.0\newline [1.0, 1.0], $n$=3 & harness \\
{\footnotesize\texttt{collateral\_\hspace{0pt}change\_\hspace{0pt}rate} (\textit{ratio})} & 0.0\newline [0.0, 0.0], $n$=3 & 0.0\newline [0.0, 0.0], $n$=3 & programmatic\_\hspace{0pt}check \\
{\footnotesize\texttt{cue\_\hspace{0pt}hit\_\hspace{0pt}rate} (\textit{ratio})} & 0.0\newline [0.0, 0.0], $n$=3 & 1.0\newline [1.0, 1.0], $n$=3 & harness \\
{\footnotesize\texttt{false\_\hspace{0pt}fire\_\hspace{0pt}rate} (\textit{ratio})} & 0.0\newline [0.0, 0.0], $n$=3 & 0.0\newline [0.0, 0.0], $n$=3 & harness \\
{\footnotesize\texttt{fire\_\hspace{0pt}latency\_\hspace{0pt}calls} (\textit{count})} & \textit{null} ($n$=0) & 0.0\newline [0.0, 0.0], $n$=3 & harness \\
{\footnotesize\texttt{context\_\hspace{0pt}availability\_\hspace{0pt}at\_\hspace{0pt}fire} (\textit{ratio})} & \textit{null} ($n$=0) & 1.0\newline [1.0, 1.0], $n$=3 & programmatic\_\hspace{0pt}check \\
{\footnotesize\texttt{abstention\_\hspace{0pt}correctness} (\textit{ratio})} & 1.0\newline [1.0, 1.0], $n$=3 & 1.0\newline [1.0, 1.0], $n$=3 & harness \\
{\footnotesize\texttt{in\_\hspace{0pt}coverage\_\hspace{0pt}answer\_\hspace{0pt}correctness} (\textit{ratio})} & 0.5\newline [0.5, 0.5], $n$=3 & 0.75\newline [0.75, 0.75], $n$=3 & harness \\
{\footnotesize\texttt{key\_\hspace{0pt}question\_\hspace{0pt}correctness} (\textit{ratio})} & 1.0\newline [1.0, 1.0], $n$=3 & 1.0\newline [1.0, 1.0], $n$=3 & harness \\
{\footnotesize\texttt{subject\_\hspace{0pt}question\_\hspace{0pt}correctness} (\textit{ratio})} & 0.0\newline [0.0, 0.0], $n$=3 & 0.5\newline [0.5, 0.5], $n$=3 & harness \\
{\footnotesize\texttt{ece} (\textit{ratio})} & 0.3667\newline [0.3667, 0.3667], $n$=3 & 0.4333\newline [0.4333, 0.4333], $n$=3 & harness \\
{\footnotesize\texttt{discrimination\_\hspace{0pt}auroc} (\textit{ratio})} & 0.75\newline [0.75, 0.75], $n$=3 & 0.3\newline [0.3, 0.3], $n$=3 & harness \\
{\footnotesize\texttt{monitoring\_\hspace{0pt}accuracy} (\textit{ratio})} & 1.0\newline [1.0, 1.0], $n$=3 & 1.0\newline [1.0, 1.0], $n$=3 & programmatic\_\hspace{0pt}check \\
{\footnotesize\texttt{provenance\_\hspace{0pt}correctness} (\textit{ratio})} & 1.0\newline [1.0, 1.0], $n$=3 & 1.0\newline [1.0, 1.0], $n$=3 & harness \\
{\footnotesize\texttt{derived\_\hspace{0pt}from\_\hspace{0pt}completeness} (\textit{ratio})} & \textit{null} ($n$=0) & \textit{null} ($n$=0) & programmatic\_\hspace{0pt}check \\
{\footnotesize\texttt{lifecycle\_\hspace{0pt}state\_\hspace{0pt}exposed} (\textit{ratio})} & 0.0\newline [0.0, 0.0], $n$=3 & 0.0\newline [0.0, 0.0], $n$=3 & programmatic\_\hspace{0pt}check \\
{\footnotesize\texttt{admissibility\_\hspace{0pt}error\_\hspace{0pt}rate} (\textit{ratio})} & 0.5\newline [0.5, 0.5], $n$=3 & 0.5\newline [0.5, 0.5], $n$=3 & programmatic\_\hspace{0pt}check \\
{\footnotesize\texttt{stale\_\hspace{0pt}read\_\hspace{0pt}rate} (\textit{ratio})} & 1.0\newline [1.0, 1.0], $n$=3 & 1.0\newline [1.0, 1.0], $n$=3 & harness \\
{\footnotesize\texttt{restart\_\hspace{0pt}survival} (\textit{ratio})} & 1.0\newline [1.0, 1.0], $n$=3 & 1.0\newline [1.0, 1.0], $n$=3 & harness \\
{\footnotesize\texttt{snapshot\_\hspace{0pt}reopen\_\hspace{0pt}survival} (\textit{ratio})} & 1.0\newline [1.0, 1.0], $n$=3 & 1.0\newline [1.0, 1.0], $n$=3 & harness \\
\bottomrule
\end{longtable}
\endgroup


\section{Limitations}
\label{sec:limitations}

Full list with statuses in \texttt{claims.md}. The following are the limits of the measurement
reported in Section~\ref{sec:compare}, placed with the results rather than after them.

\begin{itemize}[leftmargin=1.4em]
\item \textbf{Two subjects, one provenance.} Both measured servers come from one author's
reference set. The selection rule and every candidate --- included or excluded, with the ground
for each decision --- are in both reports, including the host constraint (package installation
failed with \texttt{Errno 28}, disk full, so only one MCP framework version was importable).
The sample selects for servers that launch on this host and is not representative of the MCP
ecosystem.
\item \textbf{Three repetitions.} \val{run.repetitions} repetitions per family at seeds
\valt{run.seeds}, the protocol minimum, is the only axis of variation: the reported intervals are
percentile bootstrap intervals over that axis (\valt{run.ci_method}) and are coarse by
construction. They measure run-to-run dispersion, not task-sampling variability, so a difference
smaller than the interval width should not be read as a difference between systems.
\item \textbf{One machine.} Both runs come from the single environment recorded in
Table~\ref{tab:run-identity}: a workstation, not a benchmark host, with
\valtl{liquid.identity.hardware}. Latencies are wall clock around a stdio call and depend on
machine load; no latency claim is made beyond this host.
\item \textbf{No model in the loop.} The runs are store-only: token and USD metrics are
\textit{null} by construction in both reports, so the protocol's cost axis is defined and
validated but untested. Any statement that costs change which memory design looks preferable
remains a position (claim C5).
\item \textbf{The harness is Python.} The reports were produced by
\texttt{code/run\_protocol.py} (\valtl{liquid.identity.runner___harness}). The package's tooling
rule --- new utilities in Rust, one binary with subcommands --- is met by the validation,
bibliography and table-generation pipeline (\texttt{tools/refslint}), but the runner itself is
still the card-named Python harness, because it has to interoperate with the MCP Python SDK that
both subjects import. It is identified by a content hash rather than a commit, because this
workspace carries no git metadata. There is no offline command that rebuilds a
\texttt{report.json} from its raw logs: the metric layer lives inside that runner, and re-running
the protocol requires the host setup described in the package README. What is reproducible with
one command is everything downstream of the reports: validation, the tables, and every number in
this paper.
\item \textbf{Coverage of the families.} T3 (as-of-date reads) and T4 (principal-scoped reads)
are not implemented in either subject, so those families contribute exposure metrics (timestamp
exposure, lifecycle-field exposure, principal scoping across advertised tool schemas) rather
than the operations the taxonomy defines. T6 is mostly \textit{null}. The profiles are therefore
denser for O1, O2, O5, O6 and the LIFECYCLE axis than for O3 and O4.
\item \textbf{Rubric metrics.} No double grading was performed and no $\kappa$ value exists, so
the rubric-scored metric remains exploratory (claim D5).
\item \textbf{Taxonomy and transfer.} The taxonomy's completeness is not claimed; several
store-internal steps (extraction, consolidation) may be performed by models we cannot pin and
are reported as unspecified rather than guessed; and the domain-transfer claim to industrial
settings is untested, with \cite{trashchenkov2026power} cited as the closest precedent for
executable domain evaluation rather than as evidence of transfer (claims D1--D7).
\end{itemize}

\section{Conclusion}

The reusable parts of this work are an operation checklist, a protocol whose reports are
machine-validated, a reference implementation with its gaps disclosed, a reference base in which
every source was resolved by a script, and a generated set of tables that ties every number in
the comparison to a field in a validated report. The first pair of reports now exists and is
reported with its limits: the protocol separates two memory servers on the write path,
reproduces the same readings where the mechanism is the same, and surfaces a revision difference
that retrieval scores hide. The next step is a third and a fourth subject outside one author's
reference set, a model in the loop so that the token and cost metrics stop being \textit{null},
and more repetitions than the protocol minimum.

\section*{Data and code availability}

The two runs are archived with the paper as the two run directories named in
Table~\ref{tab:run-identity}, each beside its raw MCP stdio logs (\val{liquid.raw_files} JSONL
files per run, every request, response and latency) and the per-repetition store directories the
runs wrote, so that the reported numbers can be recomputed from the responses. The measured
systems are \valtl{memory.identity.repository} and \valtl{liquid.identity.repository}, at the
commits pinned in the same table; the runner is identified there by its content
hash. The two runs issued \val{liquid.calls_issued} and \val{memory.calls_issued} MCP calls
against advertised tool surfaces of \val{liquid.tools_advertised} and
\val{memory.tools_advertised} tools respectively (counts taken from each report's
\texttt{statistics} block).

Four commands reproduce what is claimed here:

\begin{enumerate}[leftmargin=1.4em]
\item \path{bash build.sh} builds this PDF; it regenerates the bibliography from the curated
source list, checks every citation key, regenerates the comparison tables and refuses to continue
if a citation or a number lookup does not resolve.
\item Validate both reports against the published schema:\\
\path{refslint report --schema code/protocol_schema.json reports/*/report.json}.\\
\path{refslint selftest} re-derives five deliberately broken reports and asserts that all five
are rejected.
\item \path{refslint tables --root . --check} rebuilds the result tables, the number macros and
the number ledger from the two reports and fails if the files on disk differ from them, or if
this paper uses a number for which no report field exists.
\item \path{bash verify_all.sh} runs the whole integrity set --- reference base, citation keys,
report validation, negative tests, JSON parse, port parity, table freshness, PDF rebuild --- and
prints \texttt{ALL CHECKS PASSED} or fails.
\end{enumerate}

The pipeline utility is \texttt{tools/refslint} (Rust, standard library only, built by cargo;
version 0.1.0 declared in \texttt{tools/refslint/Cargo.toml}), which also generates the
bibliography, the reference-verification table and the report validation. The three documents
that carry the unabridged material --- \texttt{taxonomy.md}, \texttt{protocol.md} and
\texttt{related-work.md} --- ship in the same folder, as does \texttt{claims.md}, the ledger of
every claim with the command or source that supports it.

\section*{Acknowledgements}
Artificial-intelligence assistance (literature retrieval, metadata resolution, artefact
collection and drafting support) was used in preparing this manuscript; every cited source was
verified by machine queries and the verification table is part of the repository. The author
takes responsibility for all claims. The author thanks the maintainers of the SINT
repositories for the artefact base.

\bibliographystyle{plain}
\bibliography{references}

\end{document}
